% =====================================================================
%  Harald Høffding, OM GRUNDLAGET FOR DEN HUMANE ETHIK.  Kjøbenhavn: Andr. Fred.
%  Høst & Søns Forlag, 1876 (Sally B. Salomons Tryk).  [IV], 146 pp.
%
%  TRANSCRIPTION -- GENERATED; DO NOT EDIT BY HAND.  Made in two steps, both rerunnable
%  (ebook-method/humane-ethik/, see its scripts' docstrings):
%    python3 ebook-method/humane-ethik/draft.py --emit --force
%      this template + the OCR text layer of Google Books' scan of the University of
%      Wisconsin copy (~/bibliotek/Høffding, Harald/1876-Om_grundlaget_for_den_humane_ethik.pdf),
%      then patches.py (corrections read at the image that a word decision cannot express:
%      notes, page turns, words run together, misprint comments) and spaced.tsv (letterspacing);
%    python3 ebook-method/collate/check.py humane-ethik --apply
%      the decisions in decisions.tsv, taken against an INDEPENDENT reading of ANOTHER copy:
%      tesseract (dan) on the Royal Danish Library's scan (~/bibliotek/Høffding, Harald/
%      1876-Om-grundlaget-KB.pdf).  SCANS.tsv has both scans' sha256.
%  No e-book in hand (SAGA publishes one, ISBN 9788726363098; not used).
%
%  WORD-ACCURACY (2026-10-05).  The draft (Google's OCR) was collated word by word, with
%  punctuation, quote kinds and word division, against tesseract's reading of the KB copy:
%  714 disagreements.  389 were settled by rule: 95 by check.py (glyph faults, noise), 294
%  by rules.py with KB's own OCR layer as a THIRD witness (the draft's reading confirmed
%  205 times; dashes and quotes Google's layer drops, 22; tesseract's misread note marks
%  and quotes, 65).  The rest were read at the image: 327 crops read BLIND (the reader saw
%  only the printed lines, not the two readings) and compared by script (reads.py), the
%  remainder decided by hand from the readers' transcriptions.  OPEN 0; no decision left
%  unapplied (check.py --unapplied).  Then, at the image: every footnote read in full (45
%  notes, notesheets.py/notescmp.py); the words attested in no other transcription of the
%  corpus swept and the suspects read (7 misprints marked, 4 words the OCRs ran together
%  mended); title page, verso, Indhold and the chapter openings read.  Spot-read in full,
%  pp. 47 and 97 against the image, with two planted control errors each (4 of 4 caught):
%  1 fault found (a doubled closing guillemet, a class then mended throughout).  Not
%  established: misprints that both OCRs read alike as real words (no full image pass).
%  Letterspacing: candidates measured from letter gaps in the KB images (386), read yes/no
%  at the image; 113 set as \emph (short words may be missed).  No italics seen.
%  The print uses both »...« and «...» (and „...“ chiefly in notes and quotations); the
%  direction of every guillemet was collated.  Printer's errors are kept and marked
%  "% PRINTED AS IS" (16).
%  SECOND READING BY TRANSLATION (2026-10-05).  The English translation read the whole
%  text; its 59 odd spots were read at the image (crop sheets, KB copy; pp. 2, 10, 40 the
%  Google copy) and 27 corrections made, by patches.POST (draft.py --post, the last step):
%  3 false paragraph breaks at page turns (pp. 44, 66, 103) removed; 2 line-end hyphens
%  read as a point or lost (om-|fatte p. 15, Vil-|kaar p. 77); 6 closing marks (5 spaces
%  before », a stray » p. 25) and Bonald's lost closing mark (p. 64); ɔ: read as 0:
%  (p. 134); a speck read as an accent (Dét, p. 40); 5 names letterspaced in the print
%  (Alex. Bain, Wallace, Kant, Stuart Mill, Pitt); the two verse couplets (pp. 82, 122)
%  set as one block; 5 printer's errors marked (pp. 24, 28, 49, 61, 143: kan wanting, a
%  point for a comma, og doubled twice, a point wanting).  The other 32 spots read as
%  transcribed (among them the print's own slips: Ferguson's "the civil society",
%  Grotius's "et pacis", Longfellow's "Life ... life", "kalon").
%
%  CONVENTIONS.  "% --- p. N ---" + \opage{N} at the first word of printed page N
%  (a word broken over the turn is split with %).  Footnotes as printed, marked
%  *), **), ***) and numbered per page (\footnote[k]).  Italics \textit; letterspacing
%  \emph.  A note running over a page is set whole at its mark.  Folios and signature
%  marks are not transcribed.  The print governs; misprints are kept and marked
%  PRINTED AS IS.  Number ranges are set with --- (the print's dash there is short).  A
%  space the print sets before : ; ! ? is not transcribed.  The fractions on p. 93 are
%  small solidus fractions in the print.  Pencil annotations and library stamps in both
%  copies are not transcribed.  Page 6 is blank.
% =====================================================================
\documentclass[12pt,a4paper,openany]{book}
\usepackage[utf8]{inputenc}
\usepackage[T1]{fontenc}
\usepackage{libertinus}
\usepackage{libertinust1math}
\usepackage[danish]{babel}
\usepackage[protrusion=true,expansion=true]{microtype}
\usepackage{setspace}
\usepackage[margin=1.2in]{geometry}
\usepackage{fancyhdr}
\setlength{\headheight}{28pt}
\usepackage{hyperref}
\hypersetup{pdftitle={Om Grundlaget for den humane Ethik (1876)}, pdfauthor={Harald Høffding}, hidelinks}
\pagestyle{fancy}
\fancyhf{}
\fancyhead[LE,RO]{\thepage}
\fancyhead[RE]{\textit{Harald Høffding: Om Grundlaget for den humane Ethik}}
\fancyhead[LO]{\textit{1876}}
\fancypagestyle{plain}{\fancyhf{}\fancyhead[LE,RO]{\thepage}}
\setstretch{1.3}
\usepackage{marginnote}
\usepackage{graphicx}
\DeclareUnicodeCharacter{0254}{\reflectbox{c}}   % ɔ (id est)
\renewcommand{\thefootnote}{%
  \ifcase\value{footnote}\or *)\or **)\or ***)\else \arabic{footnote})\fi}
\newcommand{\tocline}[2]{\noindent #1\dotfill\ #2\par}
\newcommand{\opage}[1]{\marginnote{\footnotesize\textit{[#1]}}}
% a chapter of the print: new page, centred title, a bookmark and a contents line
\newcommand{\chap}[1]{\clearpage\addcontentsline{toc}{chapter}{#1}\vspace*{2em}\begin{center}{\large #1}\end{center}\medskip}
\pdfstringdefDisableCommands{\renewcommand{\footnote}[2][]{}\renewcommand{\opage}[1]{}}

\begin{document}

% --- title page (Google PDF 7, KB PDF 8) ---
\begin{titlepage}
\begin{center}
\vspace*{3em}
{\large OM GRUNDLAGET}\\[0.8em]
FOR\\[0.8em]
{\LARGE DEN HUMANE ETHIK}\\[2em]
AF\\[1em]
DR. HARALD HØFFDING\\[2em]
% (the publisher's device: an ornamental monogram)
\vfill
KJØBENHAVN\\[0.4em]
ANDR. FRED. HØST \& SØNS FORLAG\\[0.4em]
1876
\end{center}
\end{titlepage}
% --- verso (Google PDF 8, KB PDF 9): the motto, and the printer at the foot ---
\thispagestyle{empty}
\vspace*{2em}
\begin{flushright}
Was ist das Heiligste? --- Das, was heut und ewig die Geister,\\
Tiefer und tiefer gefühlt, immer nur einiger macht.\\
\hfill Goethe.
\end{flushright}
\vfill
\begin{center}{\small Sally B. Salomons Tryk.}\end{center}
\clearpage
% --- Indhold (Google PDF 9, KB PDF 10).  The library's shelfmarks on the Wisconsin copy are not transcribed. ---
\thispagestyle{empty}
\begin{center}{\large \emph{Indhold.}}\end{center}
\hfill Side\par
\tocline{Indledning}{1.}
\tocline{Psykologiske Bemærkninger}{7.}
\tocline{Individuelle Udgangspunkter}{19.}
\tocline{Avtoriteten}{55.}
\tocline{Den ethiske Lov og Fremskridtet}{87.}
\tocline{Villiens Frihed}{108.}
\begin{center}\rule{4em}{0.4pt}\end{center}
\clearpage

% --- p. 1 ---
\setcounter{footnote}{0}%
\chap{Indledning.}
\opage{1}


Skjønt den nyere Tids store Kampe imellem de forskjellige
religiøse, sociale og politiske Retninger ogsaa angaa
ethiske Spørgsmaal, tør man dog paastaa, at der hvad
Ethikens Indhold, selve Pligtloven angaar, hersker langt
større Ensartethed og Overenstemmelse, end Modstanderne % PRINTED AS IS: „Overenstemmelse“ (one s)
i Stridens Hede ere sig bevidste. Paa samme Kulturtrin
ville ifølge Tingenes Gang de almindelige Grundsætninger
og Regler for den menneskelige Handlen være de samme
hos Alle. En fælles Grund af Humanitet udgjør den Arena,
hvor Kampen føres. Fra dette Synspunkt forsvinder endog
Modsætningen mellem den kristelige og den humane Ethik.
„Det Kristelige,“ siger \emph{Martensen}, „er i vore Dage ofte
skjult under det fælles Humane og kun kjendeligt for en
nærmere Betragtning.“ Men hvorom føres da Striden?
Dels om den videre Udvikling og Gjennemførelse af hvad
der allerede er opnaaet, og dels om den dybere Forstaaelse
og Begrundelse deraf. Og her er Forholdet da nærmere
dette, at den Aand og Retning, i hvilken Udviklingen føres
videre, vil afhænge af den Maade, hvorpaa man opfatter og
begrunder det Ethiske i det Hele. Her gjør den store Modsætning
mellem de forskjellige Livsanskuelser sig gjældende.
% --- p. 2 ---
\setcounter{footnote}{0}%
\opage{2}Saa enige to Mennesker kunne være om et ethisk Forhold
eller et sædeligt Bud, ligesaa uenige kunne de være i deres
Begrundelse deraf. Og dog kan paa det ethiske Omraade
Begrundelsen ikke skilles fra Indholdet. Den Maade, hvorpaa
jeg begrunder min Handling, vil tillige være et Udtryk
for det Sindelag, hvormed den udføres, og de Motiver, der
fremkalde den, og det kommer paa det ethiske Omraade
ikke blot an paa, \emph{hvad} jeg gjør, men ogsaa, ja især paa,
\emph{hvorledes} jeg gjør det. Og skjøndt nu den bevidste
Tænkning her lige saa lidt som noget andet Sted kan give
et fuldstændigt Billede af hvad der i den levende Følelse
og den energiske Drift rører sig som bevægende Motiver,
er det dog en Opgave, som den forskende Tanke ikke kan
skyde til Side, at gjøre sig Rede for det Grundlag, hvorpaa
den ethiske Handlen hviler. At give et Bidrag til Belysningen
af dette Spørgsmaal er Øjemedet med dette lille Skrift.

I det praktiske Liv er der hverken Tid til slige Overvejelser
eller direkte Brug for dem. Livsforholdene kræve
Handling, ofte øjeblikkelig Handling. Der maa da træffes
et Valg i Medfør af den bedste Indsigt, som haves paa rede
Haand. Det er Elementer af meget forskjellig Oprindelse,
som i Forening udgjøre den praktiske Overbevisning, der
bestemmer Menneskenes Handlen, og saa vigtigt det end er
at foretage en nøjagtig Undersøgelse og Vurdering af alle
disse Elementer, saa kan en saadan Selvprøvelse i Praxis
kun foretages i meget begrændset Omfang. Dels forudsætter
den en klar Selvbevidsthed og en Færdighed i psykologisk
Analyse, som er de Færrestes Sag, og dels foretages
den ikke med kontemplativ Ro og Uinteresserthed,
men under Paavirkning af mægtige Følelser og Lidenskaber.
Ethvert Menneske er for saa vidt en praktisk Ethiker, som
hans Stræben maa gaa ud paa at bringe hans Handlen i
Harmoni med det højeste Formaal og den højeste Lov, han
kan erkjende; men hans Bestræbelse for Lutring og dybere
Begrundelse af, hvad der er Formaal og Lov for hans
% --- p. 3 ---
\setcounter{footnote}{0}%
\opage{3}Handlen, vil snart hæmmes af Skranker, som ikke blot ligge
i hans egen Individualitet, men ogsaa i de givne sociale
og historiske Forhold. En ung Fidjiboer f. Ex. kan ikke
tænke sig Andet end, at han giver sin sønlige Kjærlighed
det mest passende Udtryk ved at tage Livet af sin Moder,
medens hun endnu er i sin kraftige Alder, for at hun ikke
skal gaa over i den anden Verden som en Affældig.\footnote[1]{Se Samtalen mellem den unge Fidjiboer og Missionæren hos \emph{Lubbock}: Origin of civilisation p. 368.} En
videre gaaende ethisk Kritik vilde her forudsætte et Brud
med Stammens overleverede Udødelighedstro. --- En videre
Horizont finde vi hos dem, man kunde kalde de profetiske
Ethikere, saadanne Personligheder, der formedelst dyb Forstaaelse
af Folkets og Tidernes Vilkaar og af Menneskelivets
Grundforhold i det Hele og ved den Magt over Menneskene,
deres aandelige Overlegenhed giver dem, formaa
at opdage og udbrede nye ethiske Love, skabe nye Former
og Retninger for det ethiske Liv. Store Fremskridt i ethisk
Retning blive ofte kun mulige ved saadanne profetiske
Personligheder, til hvis Tal de store Religionsstiftere, Lovgivere
og Reformatorer høre. Men trods den større Vidde
og Dybde hos disse Aander spores dog ogsaa hos dem den
Begrændsning, som er uadskillelig fra enhver mod enkelte
praktiske Formaal rettet Stræben. Deres Synskreds betinges
ved de historiske Forudsætninger, ved Traditionen og ved
det Folks Natur, af hvilke de udspringe og paa hvilket de
skulle virke. Den psykologiske og historiske Kritik træder
ogsaa hos dem afgjort i Skygge overfor Handlingens bydende
Krav.

Hvis altsaa den Opgave, vi have fremdraget --- at finde
Elementerne i det Ethiske og Loven for deres Forbindelse,
--- i det Hele skal kunne behandles, maa et andet Standpunkt
indtages. Den theoretiske Ethiker forholder sig rent
undersøgende til, hvad den praktiske og den profetiske
% --- p. 4 ---
\setcounter{footnote}{0}%
\opage{4}Ethiker stedse kun opfatter i umiddelbart Forhold til praktisk
Handling. Naar den praktiske Ethiker er Matrosen,
den profetiske Ethiker Styrmanden, saa er den theoretiske
Ethiker Astronomen, som paa sit, i det Mindste forholdsvis,
faste Stade formaar at gjøre anderledes sikre og nøjagtige
Iagttagelser og Undersøgelser end hine, der kæmpe med
Vind og Bølge. Der ligger ikke nogen Anmasselse i den
theoretiske Ethikers Stilling overfor de to andre; ser han
klarere, er det kun, fordi han bevæger sig i de ideelle
Forestillingers Verden i Stedet for i den fyldige Virkelighed,
og vil han stille sin theoretiske Indsigt over den energiske
Overbevisning, som hine have opnaaet ad praktisk
Vej, maa han mindes det bekjendte Vers: „Et er et Søkort
at forstaa, et Andet Skib at føre!“ --- Den theoretiske
Ethiker medbringer ikke andre Forudsætninger end enhver
anden Forsker til det Emne, han vil undersøge. Det gjælder
først og fremmest at opfatte den Gjenstand, man har
valgt sig, i dens karakteristiske Ejendommelighed, at finde
det mest træffende Udtryk for den. Hvad det Ethiske angaar,
vil altsaa det Første være en Skildring af den psykologiske
Tilstand, det Forhold og Sindelag, som med Rette
bærer Navn derefter. Det Næste er at udskille de Elementer,
hvoraf denne Grundform bestaar, og at paavise de
Kilder, hvorfra de stamme; herved vil da tillige det ethiske
Liv ses i sin Sammenhæng med Menneskelivet og Verdenslivet
i det Hele. Den theoretiske Ethik er saaledes en
Erfaringsvidenskab, bygget paa Psykologi og Historie.
Den gaar ud paa Iagttagelse af den virkelig givne Handlen
og slutter derfra til almindelige Love. Men den kan dog
ikke blive staaende ved det Empiriske. Enhver videnskabelig
Erkjendelse er en ideal Konstruktion, som opfatter
de i Erfaringen givne Elementer med Forbigaaelse af
alle forstyrrende og uvæsenlige Omstændigheder. Saaledes
ere f. Ex Keplers Love Idealer, hvortil Planeternes virkelige % PRINTED AS IS: „Ex“
Bevægelser kun tilnærmelsesvis svare. Ved at abstra%
% --- p. 5 ---
\setcounter{footnote}{0}%
\opage{5}here --- d. v. s. ved at se bort fra alle Biomstændigheder ---
idealiserer man tillige, d. v. s. man danner sig et Billede
af Sagen, som den vilde være under visse simple Forudsætninger
og uden forstyrrende Indgriben af andre Forhold.
Paa denne Maade former den theoretiske Ethik ogsaa en
ideal Type, et Udtryk for den menneskelige Handlen, som
den vilde være i sin Fuldendelse. Den Konsekvens og
Renhed, hvormed en saadan Type udvikles, virker rensende
og klarende tilbage paa det virkelige Livs indviklede Forhold,
og selv om det vilde være doktrinært Fantasteri umiddelbart
at anvende en saadan Maalestok under visse givne
Forhold, har det dog sin store Betydning at blive sig bevidst,
i hvilken Retning Bevægelsen maa gaa. Opfattelsen
af det Virkelige og Opfattelsen af det Ideale tjene saaledes
hinanden til gjensidig Berigelse: det Virkelige forklares og
belyses og hæver sig derved til et højere Trin; det Ideale
modtager nærmere, konkretere Bestemmelser og bringes
derved Virkeliggjørelsen nærmere.

\begin{center}\rule{6em}{0.4pt}\end{center}% an ornamental rule ends the Indledning
% --- p. 6 --- (blank in the print)


% --- p. 7 ---
\setcounter{footnote}{0}%
\chap{Psykologiske Bemærkninger.}
\opage{7}
Vore Fornemmelser og Forestillinger opstaa hos
os som Symboler for Tingenes Væsen og Forhold.
Naar vi opfatte eller forestille os Noget, se vi bort
fra os selv og gaa op i Anskuelsen af Gjenstanden,
og jo mere vi formaa dette, desto tydeligere og nøjagtigere
en Opfattelse vinde vi. Men enhver Fornemmelse
og Forestilling virker tillige tilbage paa
vor Bevidsthed, griber paa en eller anden Maade
ind i vor hele sjælelige Tilstand, frembringer en vis
Stemning eller Følelse: --- dersom vort Bevidsthedsliv
i det Hele fremmes og forhøjes, en Følelse af
Lyst, --- dersom det hæmmes og svækkes, en Følelse
af Smerte. Disse Stemninger kunne ofte være aldeles
underordnede i Forhold til den udad rettede
Side ved Fornemmelsen eller Forestillingen; jo mere
de vinde Overhaand, desto mindre energisk og klar
bliver Opfattelsen af Tingene. Den, der er bevæget
af Lyst eller Smerte, er ikke længere en uhildet
Iagttager, og omvendt, for at anstille Iagttagelser,
% --- p. 8 ---
\setcounter{footnote}{0}%
\opage{8}maa Lyst og Smerte trænges tilbage. Der er altsaa
et omvendt Forhold imellem Erkjendelse og Følelse;
begge kunne ikke til samme Tid staa paa Højdepunktet,
men den enes højeste Grad forudsætter den
andens laveste Grad.

Alt hvad der fremkalder en Forhøjelse af den
sjælelige Tilstand, og altsaa en Følelse af Lyst, tilskrive
vi en vis Værdi, hvilket oprindelig viser sig
ved, at vi uvilkaarlig søge at fastholde Forestillingen
om det og bevæge os henimod det. Fornemmelsen
og Forestillingen i sig selv indeholder ingen Grund
til Bevægelse; det er kun ved sit Forhold til vort
sjælelige Liv i det Hele, altsaa ved den Følelse,
den vækker, at en Erkjendelse fremkalder Bevægelse
og Handling, og da der finder et omvendt Forhold
Sted mellem Erkjendelse og Følelse, maa der ogsaa
være et omvendt Forhold mellem Erkjendelse og
Handling.

Men dette er kun den ene Side af Sagen. Vi
have ved de foregaaende Bemærkninger kun taget
Hensyn til \emph{Styrkegraden}. Spørge vi derimod om
det, hvorved Følelsen vækkes, og hvorimod Bevægelsen
rettes, altsaa om \emph{Indholdet}, vil det vise
sig, at Erkjendelse og Følelse, Erkjendelse og Handling
ikke længere forholde sig omvendt, men ligefrem
til hinanden. En Følelse af Lyst fører til at
fastholde den Forestilling eller vedblive med den
Bevægelse, som fremkalder den. Men naar nu visse
Forestillinger og Bevægelser vise sig at staa i nøje
Forbindelse med hverandre, idet den ene fører til
den anden som Aarsag til Virkning, eller der er
% --- p. 9 ---
\setcounter{footnote}{0}%
\opage{9}stor Lighed og Slægtskab imellem dem, saa vil efterhaanden
Følelsen udstrækkes til alle saadanne
sammenhørende Forestillinger. Der danner sig da
en omfattende Verden af Gjenstande, som Individet
tillægger Værdi, fordi de direkte eller indirekte
vække og forhøje det sjælelige Liv. Paa denne
Verdens Indhold beroer det, i hvilken Retning Individets
Liv og Handlen vil gaa. Ethvert højere
Trin, som naaes i Erkjendelsens Udvikling, vil derfor
nødvendigvis virke tilbage paa Følelsens og
Handlingens Liv. Stille, men uafbrudt rykker Følelsen
efter i Tankens Fodspor, nuanceres, præciseres
og beriges. Jo højere og rigere Tankens Liv er,
desto dybere og fyldigere kan derfor ogsaa Følelseslivet
være. Mangfoldige Følelser, som den civiliserede
Europæer kjender, ere fremmede for den
Vilde paa Grund af hans lidet udviklede Forestillingskreds,
og de Følelser, som de have tilfælles, have
en forskjellig Karakter hos hver af dem.

Naar Erfaringen har naaet et vist Omfang og
Erkjendelsen en vis Udvikling, tages der ved Handlingen
ikke blot Hensyn til de umiddelbare og
nærværende Følelser af Lyst og Smerte, men ogsaa
til de mulige og fjernere liggende. Men dette forudsætter
allerede et Udviklingstrin, hvortil flere vilde
Folkeslag knap have hævet sig. De fylde sig med
Spise og Drikke, naar de have Overflod, uden at
tænke paa de skadelige Følger af Fylderi eller paa
den Mangel af Føde, som snart vil afløse Overfloden.
De tænke kun paa den øjeblikkelige Tilfredsstillelse.
Kun Nøden bringer dem til at arbejde igjen, og saa%
% --- p. 10 ---
\setcounter{footnote}{0}%
\opage{10}ledes blive de ved i det samme Kredsløb. Først
naar Bevidsthedslivet naar en vis Udvikling baade
i Styrke og Vidde, kunne Erindringer og Forventninger
bestemme Handlemaaden, selv om den modarbejdes
af den øjeblikkelige Lyst eller Smerte, og
vi have her den første Spore til en udover det
umiddelbart Givne og Nærværende liggende Regel
for Handlingen. Og tillige have vi her den første
klare Fremtræden af en virkelig Villie i den Energi,
hvormed de øjeblikkelige Tilskyndelser trænges tilbage,
for at de Motiver, der have vist sig at være
af væsenligere Betydning for Individets Ve og Vel,
kunne gjøre sig gjældende. I al Opmærksomhed er
der Noget af en Villen, men tydeligst fremtræder
dette paa Handlingens Omraade, saa snart Erfaringen
har fremkaldt en Bevidsthed om Handlingernes forskjellige
Følger for Mennesket.

Man har ofte forsøgt at reducere Opmærksomheden
og alle Villiesytringerne til blotte Modifikationer
af Fornemmelser og Forestillinger, en Slags
Fortætning af disse. Tvertimod har den nyere
Fysiologi og Psykologi ført til det Resultat, at
Bevægelse gaar forud for Sansning og er et mere
primitivt Udtryk for vort Væsen end denne. Paa
det bevidste Livs Trin synes det i det Mindste, som
om Bevægelser og Handlinger ledes af bestemte
Motiver og Øjemed. Men gaa vi tilbage til det
Stadium, hvor der endnu ikke er Tale om bevidste
Fornemmelser og Forestillinger, til Barnet i Moders
Liv, saa kunne dettes vilkaarlige Bevægelser ikke
udledes af noget Øjemed, ikke bestemmes af nogen
% --- p. 11 ---
\setcounter{footnote}{0}%
\opage{11}Forestilling, men have en rent spontan Karakter.
«Det bevæger,» som Johannes Müller siger, «ikke
sine Lemmer for at opnaa et ydre Formaal, men
kun fordi det kan bevæge dem.»\footnote[1]{\emph{Joh. Müller}: „Handbuch der Physiologie“ (I. p. 94). \emph{Alex. Bain} har i sit Skrift „The senses and the intellect“ (Book I) nærmere udviklet og begrundet Læren om den spontane Bevægelse, som man iøvrigt allerede vil finde hos \emph{Descartes} i hans „Traité des passions“ (I. 10---11., II. 107).} Først efterhaanden
føre saa de Erfaringer, der paa denne Maade gjøres,
til at lede de følgende Bevægelser i visse Retninger,
og endelig kan Forestillingen om den mulige Bevægelse
gaa forud for den virkelige Bevægelse og
fremkalde denne. Men selve den Maade, hvorpaa
den oprindelige spontane Bevægelse gaar for sig,
vil være betinget ved Individets medfødte Organisation,
og heraf ses igjen, i hvor høj en Grad Retningen
og Karakteren af Individets Handlen maa
bero paa den Natur, det modtager som Arv fra Forfædrene.
Den individuelle Handlen gaar her tilbage
i hele Slægtens Handlen og bestemmes ved
denne. Det er nemlig ikke blot ved det individuelle
Livs Begyndelse, at den spontane Bevægelse spiller
en Rolle; alt hvad vi kalde instinktmæssig Handlen,
maa henføres dertil. Det er Handlinger, vi foretage
uden at vide hvorfor, og hvis egentlige Karakter vi
først senere blive os bevidste, lige som den unge
Fugl sætter sine Stemmeredskaber i Bevægelse uden
at vide, hvad den gjør, og først derefter lærer sine
egne Toner at kjende. Og selv om bevidste Motiver
gjøre sig gjældende, kjender Enhver dog af egen
% --- p. 12 ---
\setcounter{footnote}{0}%
\opage{12}Erfaring Tilfælde, hvori hverken noget enkelt af
disse Motiver eller alle tilsammen indeholde den
tilstrækkelige Forklaring af Handlingen; man maa
da her gaa tilbage til det dybere Instinkt, hvad enten
dette nu har dannet sig som en anden Natur i
Løbet af Individets Liv eller udspringer af den oprindelige
Organisation.

Den Sætning, at al Handlen udspringer af Lyst
eller Smerte, er altsaa kun delvis sand; Mennesket
handler ikke altid efter Motiver, fordi han ikke altid
handler med Bevidsthed. Det er af Vigtighed at
fastholde dette, da man ellers vil komme til en falsk
Opfattelse af den menneskelige Villen og Handlen
paa dens primitive Trin. Man har nemlig ofte brugt
Sætningen om Lyst og Smerte som Handlingens
Grundmotiver til at bevise, at Egoismen er den
dybeste Drift i den menneskelige Natur. Men virkelig
Egoisme forudsætter en Reflexion, en Beregning
og altsaa en Forstandsudvikling, hvorom der
paa de primitive Trin endnu ikke er Tale. Den
instinktmæssige Selvopholdelse kan man ikke med
Rette kalde Egoisme. Først paa et højere Udviklingstrin
kan Individet løsrive sig fra den Sammenhæng,
hvori det oprindelig staar, gjøre sig selv til Midtpunktet
og nedsætte alt Andet til Middel. Oprindelig
gjør Individet ikke denne Forskjel, skjælner
ikke mellem Lyst og Smerte og det, der fremkalder
dem, men mener i selve Lysten eller Smerten at
lære Tingenes Væsen at kjende. Herpaa beroer det
netop, at Mennesket kan opdrages; han kan ved de
skiftende Stemningers Magt føres ud over sin op%
% --- p. 13 ---
\setcounter{footnote}{0}%
\opage{13}rindelige Begrænsning, hvad der ikke var muligt,
hvis han virkelig fra Begyndelsen af havde isoleret
sig i sig selv.

Lige saa lidt som vort fysiske Verdenssystem
vilde have kunnet udvikle sig af den oprindelige
Taagemasse, dersom den frastødende Kraft havde
været den ene raadende, lige saa lidt vilde Humanitetens
Verden være bleven til, naar kun én enkelt
Drift, den isolerende Selvopholdelsesdrift, havde hersket
i den menneskelige Natur. Ogsaa i denne virker
en Flerhed af Kræfter. Der ligger i den lige saa
vel et Anlæg til Sympathi som til Egoisme; der er
i vor Natur «Noget af Duen ved Siden af Slangen
og Ulven.» Vi finde sympathetiske Drifter allerede
indenfor Dyreverdenen, i Forældrenes Kjærlighed
til og Opofrelse for deres Unger og i Sammenholdet
mellem Individer af samme Flok. Hos Mennesket
udvikles og beriges disse Drifter lige som alle andre,
eftersom den aandelige Udvikling skrider frem. Den
uvilkaarlige Medliden og Medfølelse med Andre kan
ikke udledes af nogen Beregning om, hvilke Følger
deres Lykke eller Ulykke kan have for os. Sympathien
forudsætter kun, at vi selv af Erfaring vide,
hvad Lyst og Smerte er, og at vi virkelig betragte
Andre som Væsener af samme Art som vi. Vi
sætte os uvilkaarlig i deres Sted, gjøre deres Følelser
til vore. I hvilket Omfang Sympathien gjør sig
gjældende, afhænger af, hvorvidt Bevidstheden om
Fællesskabet strækker sig. Sin oprindelige Arne
har den sympathetiske Følelse i Familieforholdet,
og hos meget raa Stammer kan det selv her let se ud,
% --- p. 14 ---
\setcounter{footnote}{0}%
\opage{14}som om den overskyggedes af den tyranniske Selvopholdelsesdrift,
idet Individet i fortvivlet Kamp for
at bestaa blindt raser mod Alt, hvad der staar det
i Vejen, for saa vidt Overmagt ikke holder det
Stangen. Men Sympathien er i sig selv lige saa
naturlig for Mennesket, som Familieforholdet er det,
og lige som Menneslægten vilde gaa til Grunde, naar % PRINTED AS IS: „Menneslægten“ (for Menneskeslægten)
dette Forhold ophørte, vilde det Samme ske, naar
al Sympathi manglede. Fra Familiens snevrere
Kreds udbreder den sig tilsidst til Folket, Stammen
og tilsidst til hele Menneskeheden.

Baade ved at følge Selvopholdelsesdriftens og ved
at følge Sympathiens Tilskyndelser ledes Mennesket
til at betragte sig som Led i en større Helhed.
Hvad enten hans Drifter tilfredsstilles eller møde
Modstand, faar han et Indtryk af bestemte Vilkaar
og fast Sammenhæng, af en Tingenes Orden, under
hvilken han maa bøje sig. Det er, som vi have bemærket,
Erkjendelsens Udvikling, som gjør en videre
Udvikling af Følelserne mulige. Men Erkjendelsen
forholder sig ikke blot som tjenende Redskab. Den
giver selv væsenlige Bidrag til Handlingens Karakter.
Den opstiller almindelige Love og Regler,
gjennemfører konsekvent ét Synspunkt. Følelse og
Lidenskab have Tendens til at gjøre Undtagelser,
til at virke uregelmæssig og kun i enkelte Retninger.
Den udviklede Erkjendelse gjør Handlingen konsekvent
og retter den mod ét væsenligt Formaal. Individet
fornemmer i den erkjendte Sandhed en Myndighed,
der staar uberørt af hans Ønsker og Drifter;
han har i den et fast Punkt, hvortil han kan tage
% --- p. 15 ---
\setcounter{footnote}{0}%
\opage{15}sin Tilflugt for at vinde en uhildet Opfattelse af
sine Vilkaar og sin Færd. Det personlige Liv trænger
til at lutres ved den upersonlige, den over al
Vilkaarlighed og Tilfældighed hævede Sandhed.

Af det nu Udviklede drage vi nogle almindelige
Slutninger, som ville føre os nærmere til vort egenlige
Emne. For at Mennesket skal komme til at
handle, maa der for det Første være tilstrækkelig
stærk Anledning dertil. Anledningens Styrke betyder
her den Grad, i hvilken Menneskets Natur i
et givet Øjeblik disponeres til Handling, og en saadan
Disposition kan, som vi have set, enten være
et oprindeligt Instinkt, en ubevidst Drift, eller en
ved visse Fornemmelser og Forestillinger vakt Følelse
af Lyst eller Smerte. Men for det Andet have vi
set, at Erfaring og Erkjendelse have en stor Indflydelse
paa, hvad der vækker Lyst og Smerte, og
at Mennesket i sin Handling bestemmes ved den
større Helhed, hvoraf han er Led, og som han kan omfatte
med uinteresseret Sympathi. For saa vidt vi
blot se hen til \emph{Styrken} af, hvad der driver Mennesket
til Handling, synes vi kun at have med ham
selv at gjøre; men Handlingens \emph{Indhold og Formaal}
viser os ud over de enkelte Individer. Spørgsmaalet
er nu, om vi i nogen af de her fremhævede
psykologiske Love have Udgangspunkter for Opfattelsen
og Begrundelsen af det Ethiske.

At det Ethiske ikke beroer paa Styrken af, hvad
der leder til Handling, vil let indses. Stærke Følelser
og Instinkter ere umiddelbare Ytringer af Individets
Temperament og medfødte Karakter. Det
% --- p. 16 ---
\setcounter{footnote}{0}%
\opage{16}kan være et Fortrin ogsaa i ethisk Henseende at
være udrustet med stærke Naturdrifter, men Ingen
vil umiddelbart finde det Ethiskes Væsen i dem.
Alle ville være enige i, at det Ethiske altid forudsætter
en Lov, en Regel, et Formaal, som ikke
umiddelbart falder sammen med Drift eller Instinkt.
Det gaaer med det Gode som med det Sande. Vi
kalde ikke en Forestilling eller en Mening for sand,
førend vi have overbevist os om dens Overensstemmelse
med det Virkelige efter den Maalestok, vi i
denne Henseende besidde. Og vi kalde heller ikke
en Handling for god, uden for saa vidt den stemmer
med den almindelige Regel, som efter vor Mening
bør gjælde i det menneskelige Liv. Her forudsættes
altsaa begge Steder en højere Lov eller Maalestok.
Vi vises da over til den anden Side ved den menneskelige
Handlen, --- dens Indhold og Formaal. Heller
ikke her vil man umiddelbart finde det Ethiske.
Der ligger ganske vist, som vi tidligere have set,
allerede i Menneskets umiddelbare Væsen en Evne
til at give sig hen og til at anerkjende Noget, som
ligger ud over hans eget individuelle Jeg. Men kun
naar denne Anerkjendelse bliver bevidst, vil den faa
Navn af ethisk. Og dette vil da netop være den
mest almindelige Eestemmelse, vi kunne opstille af % PRINTED AS IS: „Eestemmelse“ (for Bestemmelse: the first letter reads E at the image and to both OCRs)
den ethiske Handlen: en Handlen, hvis Motiver ere
Forestillinger og Følelser, der pege udover det individuelle
Jeg og godtgjøre, at Individet ser sig som
Led i en mere omfattende Tingenes Orden, hvis
Formaal han gjør til sine, og hvis Lov han betragter
som Regel for sit Liv. Det græske Ord, hvoraf
% --- p. 17 ---
\setcounter{footnote}{0}%
\opage{17}Udtrykket «ethisk» er dannet, lige som det latinske
Ord, hvoraf det ensbetydende Udtryk «moralsk»
kommer, betyder Skik og Sædvane; i det danske
Ord «Sædelighed» have vi den samme Grundbetydning.
Ved disse Udtryk vises vi hen til en almindelig
Lov og Regel, hvorefter Individets Handlen
retter sig, men der ligger egentlig ikke i dem, at
denne Lov er Gjenstand for hans levende Følelse
og Tilslutning. Den blotte Tilvanthed knytter
ganske vist den Enkelte til et større Hele, men kan
i sig selv meget godt have en rent mekanisk Oprindelse.
Sædvane er endnu ikke Sædelighed.

Hvor vidt denne almindelige Definition af den
ethiske Handlen, som motiveret ved Anerkjendelse
af en højere Lov og Ærefrygt for den, er rigtig,
vil først efterhaanden vise sig under vore følgende
Undersøgelser. Først i Løbet af disse vil der ligeledes
blive Spørgsmaal om den højere Lovs Natur,
og hvor vidt vi i det Hele kunne tale om en almengyldig
ethisk Lov, samt om Forholdet mellem den
og de forskjellige Love, som Menneskene til forskjellige
Tider og paa forskjellige Steder have følt
sig bundne til.

Den foreløbige Bestemmelse af det Ethiskes
Grundlag, som vi have opstillet, vil i det Følgende
blive støttet først ved en mere indirekte, dernæst
ved en mere direkte Bevisførelse. For det Første
ville vi nemlig undersøge Berettigelsen af de Forsøg,
som man har gjort paa at udlede det Ethiske
af et eller flere af de i det Foregaaende fremhævede
individuelle Udgangspunkter: den rene Selvophol%
% --- p. 18 ---
\setcounter{footnote}{0}%
\opage{18}delsesdrift, Sympathien og Fornuften. Hvis disse
Elementer indeholde det fulde Grundlag for det
Ethiske, vil vor Definition være urigtig; dersom derimod
hine Forsøg vise sig at være ufuldstændige,
vil den vinde en indirekte Bekræftelse; thi naar
Grundlaget for det Ethiske ikke kan søges i det
rent Individuelle, maa det søges i det Forhold, hvori
Individet staar til en højere Magt. Det Andet, vi
have at gjøre, vil saa være, at vise, hvorledes et
saadant Forhold til en højere Magt bliver til, og
hvorledes den Anerkjendelse og Ærefrygt udvikler
sig, i hvilken vi se det egenlig ethiske Motiv.
--- Den første Undersøgelsesrække vil altsaa angaa Individualiteten,
den anden Avtoriteten.

% --- p. 19 ---
\setcounter{footnote}{0}%
\chap{Individuelle Udgangspunkter.}
\opage{19}
1. Allerede i det Foregaaende have vi set, at
man kun med Urette kunde betegne Menneskets af
Lyst og Smerte bestemte Handlen som egoistisk. Naar
man forsøger at udlede det Ethiske af Lyksalighedsdriften,
af Menneskets oprindelige Stræben efter
Velvære, behøver man ikke derfor at nægte, at der
gives andre Følelser og Drifter i den menneskelige
Natur end de egoistiske. Det er nok, naar man
gaar ud fra, hvad Ingen vel kan nægte, at disse fra
først af ere de stærkeste. For nu at finde en ret
fast og paalidelig Grund at bygge det Ethiske paa,
henvender man sig til den oprindelige Grunddrift,
som under forskjellige Former maa finde sin Tilfredsstillelse
ved enhver menneskelig Handlen. Her er
man sikker paa at raade over tilstrækkelig vægtige
Motiver; man appellerer til den individuelle Interesse
og søger at vise, at den ethiske Handlen kun er en
% --- p. 20 ---
\setcounter{footnote}{0}%
\opage{20}særegen Maade at tilfredsstille denne Interesse paa.
Det Ethiske er den vel forstaaede Interesse.

I Oldtiden repræsenteredes dette Standpunkt af
\emph{Aristip} og \emph{Epikur}. Individets Følelse af Lyst
og Smerte er for dem Grundlag og Princip for al
Handlen. \emph{Aristip} betegnede nærmere Lysten som
øjeblikkelig Nydelse. Opgaven blev da for ham at
besidde en saadan Virtuositet, at man i hvert Øjeblik
kan forholde sig nydende. Men denne Isolation
af Øjeblikket fra Fortid og Fremtid viste sig snart
umulig at gjennemføre, lige som hin Virtuositet kun
er faa Dødeliges Sag, der med ydre heldige Vilkaar
forbinde aandelig Bevægelighed og et sangvinsk Temperament.
Epikur gik derfor ikke ud fra den øjeblikkelige
Lyst, men saa Grundlaget for den rette
Handlen i Menneskets Tragten efter Lykke i Betydning
af en varig og med Klogskab sikret Tilfredshed,
der forudsætter Resignation overfor de uundgaaelige
Smerter. Den mindre opofres over for den
større, den kortere for den varigere. Medens Aristip
glad hengav sig til det Nærværende, trak Epikur
sig tilbage i sig selv, i stille Nydelse af Freden i
sit Indre, idet Lidenskabens Uro er dæmpet af den
fornuftige Overvejelse, og Frygten for det Uventede
og det Overnaturlige er fjernet ved Indsigt i Naturens
nødvendige Love. Til andre Mennesker slutter
Epikuræeren sig kun for at forhøje sin Nydelse
ved at dele den med sine Venner. Samfundslivet
opstaar ved den gjensidige Nytte, Menneskene kunne
yde hverandre; det skaffer Sikkerhed tilveje og er
saaledes en nødvendig Betingelse for den Vises Lyk%
% --- p. 21 ---
\setcounter{footnote}{0}%
\opage{21}salighed, da han ellers vilde leve i en bestandig Uro
og vilde drages bort fra sig selv i anstrengt Virksomhed
udadtil.

De beslægtede Retninger, som ere fremkomne i
den nyere Tid, have trods det fælles Udgangspunkt
to karakteristiske Ejendommeligheder. Hovedsagen
er ikke længere at skaffe Individet et idyllisk Tilflugtssted,
hvor han kan nyde Livet i Ro, men Tilfredsstillelsen
af en rastløs Trang til Fremskridt.
Karakterer som Hobbes og Bentham have lidet tilfælles
med Aristip og Epikur; de ere Polemikere og
Reformatorer. Hermed staar ogsaa i Forbindelse,
at de langt mindre tænke paa den Enkeltes end paa
alle Enkeltes Lykke. Det var naturligt for Aristip
og Epikur at trække sig tilbage til det individuelle
Jeg, da de levede i en Periode, hvor den offenlige
Interesse havde tabt sig og en hel Kulturs Opløsningsperiode
var begyndt. Men det var lige saa
naturligt for Mænd af beslægtet Retning i den nyere
Tid, hvor Kultur- og Samfundslivet er i frodig Fremgang
og nye Kræfter og Ideer røre sig, at anlægge
et mere omfattende Synspunkt.

Det første, af Naturen selv anviste Gode er ifølge
\emph{Hobbes} Selvopholdelsen. Alle stræbe efter at bevare
og nyde Livet og undgaa Døden, især en smertefuld
Død. Af dette første Gode kunne andre Goder
udledes. Saaledes er Magt, Rigdom, Venskab og
Visdom Goder, fordi de tjene til Livets Bevarelse
og Nydelse. Et højeste Gode, et absolut Formaal
kan ikke tænkes; det vilde være en Tilstand, hvor
al Modsætning og Bevægelse var ophørt; men saa
% --- p. 22 ---
\setcounter{footnote}{0}%
\opage{22}vilde ingen Nydelse, ikke engang nogen Opfattelse
være mulig, da al Sansning forudsætter Afvexling
og Flerhed. Den største Lykke bestaar i, under
saa faa Hindringer som mulig, at gaa fremad mod
højere og fjernere Formaal.

Da Selvopholdelsesdriften er lige stærk hos Alle,
men den Ene ofte staar den Anden i Vejen, saa leve
Menneskene, siger Hobbes, i en stadig Frygt for
hverandre indbyrdes. De Svage frygte de Stærkes
Magt; de Stærke frygte de Svages List. Det eneste
Middel til Sikkerhed er at foregribe den Anden ved
at berøve ham Evnen til at skade. Derved opstaar
en Kamp af Alle imod Alle. Menneskene ere som
rasende Ulve mod hverandre. I denne Tilstand
kan ingen Kultur, intet sandt Menneskeliv udvikle
sig. Menneskets Evne til at kaste Blikket ud i
Fremtiden gjør det kun ulykkeligere end Dyret.
Frygten og Bekymringen vaagne stedse paa ny, lige
som Prometheus' Lever voxede ud igjen om Natten
efter at være fortæret om Dagen. Naar nu Fornuften
udvikles, lærer den Mennesket, at Krigstilstanden
er en Kilde til al hans Elendighed og
derfor bør undgaas for enhver Pris. Freden er
altsaa et Gode, om hvis Værdi Alle kunne blive
enige; den stemmer med Alles Tarv og fører dem
ud over Naturtilstandens Ulykker. Fornuften indskærper
de naturlige Love, som Menneskene maa
følge, dersom de ville arbejde sig ud af den oprindelige
dyriske Tilstand. Disse Love sætte Skranker
for den individuelle Frihed, for Lidenskabernes
umiddelbare Indflydelse, og indskærpe visse Dyder
% --- p. 23 ---
\setcounter{footnote}{0}%
\opage{23}som nødvendige Midler til et fredeligt Liv. Saadanne
Dyder ere Retfærdighed, Taknemmelighed,
Forsonlighed, Medlidenhed o. s. v. --- Men for at de
naturlige Love skulle være mere end blotte Forskrifter
og Læresætninger og kunne optræde med
tilstrækkelig Myndighed overfor de ubændige Drifter,
maa Alle underkaste sig en fælles Suveræn, en absolut
Avtoritet, overfor hvilken de hver for sig aldeles
opgive deres egen Villie. Naturloven, som
foreskriver Midlerne for Selvopholdelsesdriftens Tilfredsstillelse,
foreskriver derfor ogsaa fuldstændig
Underkastelse under Statsmagtens Villie som eneste
Middel til at sikre Freden.

\emph{Bentham} er ikke pessimistisk som Hobbes;
han anser det ikke for nødvendigt at svinebinde
Menneskene, for at de ikke skulle sønderrive hverandre
som vilde Dyr; men han fastholder dog Lyksalighedsdriften
som eneste Grundlag for det Ethiske.
Overfor denne Drift er enhver Tale om Pligt og
Selvopofrelse kun tomme Ord; hvorledes skulde saadanne
Fraser kunne overvinde den mægtigste, i
vort Væsen dybest grundede Trang? Og paa den
anden Side, hvad vilde det være for en Dyd, der i
ingensomhelst Henseende bidrog til at fremme menneskelig
Lykke? Og selv om der gaves en saadan
Dyd, ved hvilke Motiver vilde man da kunne faa
Menneskene til at stræbe efter den? Man maa reducere
alle Spørgsmaal om Godt og Ondt til Spørgsmaal
om Lyst og Smerte; først da kan man komme
til en sikker og bestemt Vurdering. Ud fra den
simple Modsætning mellem Lyst og Smerte kunne
% --- p. 24 ---
\setcounter{footnote}{0}%
\opage{24}alle ethiske Dyder, alle ethiske Forhold begrundes.
Giv mig Glæde og Sorg, Lyst og Smerte, saa vil
jeg skabe en ethisk Verden!

Da Grunddriften i Mennesket oprindelig er rettet
mod Lykken, drejer det sig blot om at lede den ret;
thi af Kortsynethed og fejl Beregning gaar Mennesket
ofte glip af, hvad han søger. Ethikens Opgave
er det da at lære Mennesket at regne rigtig. Han
maa lære at ofre det Mindre for det Større. Kun
saadanne Ofre kunne forsvares, som forøge Lysten
eller formindske Smerten, nu eller i Fremtiden.
Klogskab bliver derfor den første eller egenlig den
eneste Dyd. Maadehold og Selvbeherskelse staa i
nøje Sammenhæng med den. Det er imidlertid ikke
enhver Stræben efter Lykke eller Nytte, som fortjener
at kaldes dydig; for at en Stræben skal føre
dette Navn, maa den være forbunden med Anstrengelse
og have en Modstand at overvinde. Anstrengelsen
bestaar i at fastholde den vundne Indsigt i,
hvad der bringer Lyst og Smerte, overfor Øjeblikkets
Tilskyndelser. En uethisk Handling beror paa en
fejl Beregning af den personlige Interesse; men Individet
har ikke blot sine egne Interesser at tage
Hensyn til. For ret at kunne naa sit eget Maal
maa han tjene Andres Interesser. En stor Del af
hans Glæder og Nydelser beror paa Andres Villie,
og han kun sikre sig deres Medvirken ved selv at % PRINTED AS IS: „han kun sikre“ (a word, kan, wanting; both copies)
virke for deres Ønsker. Det er selve den personlige
Interesse, det stærkeste af alle Baand, som knytter
den Enkelte til hele Slægten. Bild Dig ikke ind, at
Menneskene ville bevæge Enden af deres Finger for
% --- p. 25 ---
\setcounter{footnote}{0}%
\opage{25}at tjene Dig, naar de ikke kunne have nogen Fordel
deraf det er aldrig sket og vil aldrig ske, saa
længe Menneskenaturen er, som den nu er. Men
Menneskene ville ønske at tjene Dig, naar de finde
deres Fordel derved, og der er utallige Lejligheder,
hvor de kunne være Dig til Nytte, samtidig med at
de ere sig selv til Nytte. Saadanne Lejligheder opdager
Forstanden, men Mængden bliver den ikke
vaer. I disse gjensidige Tjenester bestaar Dyden.
--- Ved Siden af Klogskaben som Hoveddyd stilles
altsaa den beregnende Velvillie. Men Bentham indskærper
udtrykkelig, at Kjærligheden til sig selv
ligger til Grund for Kjærligheden til Andre, eller
som han udtrykker sig, at det sociale Princip er
underordnet det personlige Princip. Ethiken, saaledes
som Bentham opfatter den, kunde man definere
som Læren om de Omveje, Menneskene nødes til at
gjøre, for at tilfredsstille deres Stræben efter Lykke.
Den vigtigste af disse Omveje betinges ved Hensynet
til Andre. Bentham lægger saa stor Vægt
paa dette Punkt, at han som almindeligt Princip opstiller:
den størst mulige Lykke for det størst mulige
Antal af Mennesker! Dette Princip ophæver
ikke den personlige Klogskabs Bud; thi hvert enkelt
Individs Lykke er jo lige saa vel en Del af Lykkens
fælles Sum, som alle Andres. Principet fyldestgjøres
altsaa ikke bedre, end naar Enhver for sit
Vedkommende bereder sig saa megen Lykke som
mulig.

Bentham polemiserer heftig mod Ord som «Pligt»
og «bør», fordi de efter hans Mening begunstige
% --- p. 26 ---
\setcounter{footnote}{0}%
\opage{26}ubegrundede Avtoritetsbud. Naar man har udtalt,
at Noget er Pligt eller bør ske, anser man Sagen
for endt uden nærmere at begrunde den. Men der
er et lille Ord, som har bødet paa den Skade, hine
Ord have voldet, det er Ordet \emph{hvorfor}. Naar man
kun vedbliver at bruge dette, vil man, haaber Bentham,
gjøre Ende paa Avtoritetstroen og de mystiske
Indskydelsers Herredømme i Ethiken. Og her har han
i Virkeligheden ogsaa en stor Fortjeneste. Han har
saa energisk som maaske ingen anden Ethiker forlangt
Grunde for hvert enkelt ethisk Bud. Han har
fremdeles paavist den nøje Sammenhæng mellem
Avtoritetstroen og Askesen; begge ere lige tyranniske
over for Menneskenaturen, lemlæste denne hver
paa sin Vis. Han har endelig i sit Princip om den
størst mulige Lykke for Alle opstillet en Maalestok
for praktiske Reformer, som allerede har godtgjort
sin store Betydning. Men det centralt ethiske Forhold
har han ikke Øje for i sin Theori. Er der
Noget en Ethiker skal beskjæftige sig med at underdersøge, % PRINTED AS IS: „under-|dersøge“ (for undersøge)
maa det være selve Forpligtelsens Begreb.
At Bentham vil udstøde dette af Ethiken, har sin
naturlige Grund deri, at dette Begreb ikke kan faa
nogen Plads i hans egen Opfattelse; thi man behøver
jo ikke at forpligtes til, hvad man umiddelbart og
uvilkaarligt stræber efter. Paa den anden Side undlader
Bentham ikke at bruge baade Ordet «Pligt»
og Ordet «bør», naar han vil indskærpe Klogskabens
og Velvilliens Dyder. Og den eneste Sanktion, han
da kan henvise til som Grundlag for Forpligtelsen,
bliver Frygten for Straf, enten den kriminelle Straf
% --- p. 27 ---
\setcounter{footnote}{0}%
\opage{27}eller den offenlige Menings Fordømmelse. Her appellerer
han lige som Hobbes til Frygten som sidste Motiv.

De almindelige psykologiske Undersøgelser, vi
tidligere have anstillet, afgive de fornødne Elementer
til en Vurdering af de nu fremstillede Forsøg.
Hvad Bentham (vi holde os til ham, da han har
givet Theorien den mest fuldendte Form) forudsætter,
er, at Menneskets Handlen kun foregaar efter bevidste
Motiver, der mekanisk kunne afvejes mod
hverandre. Han overser den dybere Sammenhæng
mellem det enkelte Individ og hele Slægten, af
hvilken det udspringer, og i hvis Midte det lever.
Det er ifølge Benthams Lære kun den personlige
Interesse, som knytter Menneskene sammen. Skjønt
han ikke deler Hobbes' Lære om den oprindelige
Kontrakt, hvorved Samfundslivet stiftedes, for at
hver Enkelt kunde vinde Sikkerhed og Fred, staar
han dog i Principet endnu paa samme Standpunkt.
Samfundet og Slægten er for ham kun et Resultat,
et mekanisk Aggregat af Individer, ikke tillige det
Moderskjød, hvoraf disse udspringe. Men allerede
Hobbes selv var bleven opmærksom paa, at der dog
var et Samfund, der ikke kunde forklares paa denne
Maade, ikke kunde være opstaaet ved Parternes individuelle
Vilkaarlighed, nemlig Familiesamfundet,
særlig Samfundet mellem Moder og Barn. Her gjør
sig ikke blot et umiddelbart sympathetisk Instinkt
gjældende, men vi have her tillige et Forhold, hvor
Avtoriteten paa ingen Maade kan være Følge af
Kontrakt. Barnet staar umiddelbart under Moderens
Myndighed, eller som Hobbes paa sit Sprog
% --- p. 28 ---
\setcounter{footnote}{0}%
\opage{28}udtrykker det, en Søn kan der ikke gives i Naturtilstanden.
Mennesket begynder altsaa i Virkeligheden
aldrig helt forfra, og allerede heri ligger et
Fingerpeg om, at heller ikke Ethiken, Læren om
den rette menneskelige Handlen, kan begrundes dybt
nok ved at tage et rent individualistisk Udgangspunkt.
Mennesket gjør som nydende og lidende Alt til
Middel for sig, direkte eller indirekte; men det
ethiske Forhold opstaar først, naar Mennesket føler
sig ikke blot som Formaal, men som Middel, ikke
blot som Herre, men som Tjener, --- Middel og
Tjener for Formaal, der ere mere omfattende end
hans egen individuelle Lyst og Smerte. Dette strider
ikke mod, hvad Bentham med saa stor Ret fremhæver
overfor den asketiske Ethik, at ingen Opofrelse
har Værdi og kan forsvares. som ikke frembringer % PRINTED AS IS: „forsvares. som“ (a point for a comma)
mere Lykke end der ofres. Thi fordi Individets
Lykke ikke kan være Grundlaget for det
ethiske Forhold, er det ikke sagt, at det Ethiske
staar fjendtlig over for den. Det kunde jo være,
at den højeste Lykke, den fuldeste Udfoldelse af
Menneskets Evner og Kræfter, netop forudsætter, at
han staar i en højere Sags Tjeneste og ganske hengiver
sig til den uden ængstelig at agte paa sin
egen individuelle Lyst eller Smerte.

Først naar Spørgsmaalet ikke er om Ethikens
Grundlag men om dens Indhold, viser Benthams
Storhed sig. Det er et stort filanthropisk og reformatorisk
Princip, han har opstillet i Sætningen
om den størst mulige Lykke for Alle. Hans Lære
er snarere en Theori om den rette Lovgivning end
% --- p. 29 ---
\setcounter{footnote}{0}%
\opage{29}en ethisk Theori. Lovgivningen maa stedse stræbe
efter at bringe en Harmoni tilveje mellem de individuelle
Interesser. Den ser paa Resultatet, ikke
paa Motivet. Saaledes lærer Bentham ogsaa, at det
er selve Handlingen og dennes Følger, ikke Motivet,
det kommer an paa. Alle Motiver ere jo efter
hans Opfattelse kun forskjellige Former for en og
samme Grunddrift; Menneskene handle kun for at
opnaa Lyst og undgaa Smerte; heri stemme altsaa
alle Handlinger overens, hvorimod deres Følger ere
meget forskjellige. Alligevel er det heller ikke Følgerne
i sig selv, der efter Benthams Opfattelse gjør
en Handling ethisk eller uethisk; han kalder kun
den Handling dydig, hvor en mere omfattende Erkjendelse
og ikke den øjeblikkelige Tilskyndelse er
den bestemmende Grund; Dyden forudsætter Anstrengelse.
Han gaar altsaa selv tilbage til Motivet (hvoraf
naturligvis Hensynet til Følgerne altid maa være en
Del), og bekræfter den foreløbige Definition af den
ethiske Handlen, vi have opstillet, idet ogsaa han
fordrer en højere Lov som bestemmende for Handlingen,
kun at denne Lov efter hans Opfattelse udspringer
af selve Lyksalighedsdriften, er en direkte
Fortolkning af dennes Krav. Men dette er for snævert
et Standpunkt.\footnote[1]{I det Forsvar, Stuart Mill har ført for „Utilitarianismen“ eller den paa Lyksalighedsdriften grundede Ethik, har han, som jeg et andet Sted (Den engelske Filosophi i vor Tid, p. 51---58) har vist, omdannet denne Lære saaledes, at han næsten ikke længere forsvarer, hvad Modstanderne, have angrebet, Det er især to Punkter, hvori han afviger \opage{30}fra Bentham. Han forlader Benthams og de ældste Utilitarianeres Lære, at alle Lyster ere lige i Art og kun forskjellige i Styrke og Varighed, og skjelner mellem højere og lavere Nydelser. Men da Maalestokken for denne Værdiforskjel ikke kan søges i selve Nydelsen, er Utilitarianismen hermed egenlig forladt. Han lægger dernæst langt større Vægt paa de sympathetiske Følelser og ser i dem det Ethiskes egenlige Grundlag, hvorved han opgiver det rent individualistiske Udgangspunkt, som var det ene herskende hos Bentham.}% the note runs over from p. 29 to p. 30 (joined by draft.py) % PRINTED AS IS: „angrebet, Det“ (a comma for a full stop; image)


% --- p. 30 ---
\setcounter{footnote}{0}%
\opage{30}Det kunde synes, som om Læren om Lyksalighedsdriften
som den ethiske Handlens Grundlag, bekræftedes
ved \emph{Darwins} bekjendte Lære om Kampen
for Tilværelsen. Lige som Planten og Dyret naar
ogsaa Mennesket kun sit Væsens højeste Udvikling
ved Arbejdet for at bevare og opholde Livet. Det
er en Kjendsgjerning, at Kultur og Samfundsliv for
en stor Del kun ere mulige ved Interessernes Vexelvirkning.
Det er den aldrig hvilende Stræben efter
højere og varigere Livsnydelse, som fører til nye
Arbejder og Opdagelser; den Vilde holdes ved sine
faa og simple Fornødenheder tilbage paa et lavere
Kulturtrin. Men hvor højt et Trin i Kultur man
end forudsætter, er Livet dog stedse en Kamp for
at bestaa. Hvert Øjeblik kan det simple Spørgsmaal
om at være eller ikke være stilles, og Selvopholdelsesdriften
kan da ytre sig i hele sin brutale
Nøgenhed. Saaledes f. Ex. ombord paa et Skib, der
brænder i aaben Sø, og hvor Mandskabet bemægtiger sig
Baadene, medens Passagererne slaas paa Liv og Død
om Redningsbælterne. Den Stærkeste og den mest
Hensynsløse vil her lettest undkomme. Hvad kun
en enkelt skrækkelig Begivenhed fremkalder blandt
% --- p. 31 ---
\setcounter{footnote}{0}%
\opage{31}de civiliserede Mennesker, er staaende Vedtægt hos
de Vilde, der lige som Dyrene skille sig af med
syge og affældige Individer.

Hvad der først maa bemærkes, er, at selv om
Egoismen skulde være det eneste Motiv i Kampen
for Tilværelsen, vilde denne Omstændighed alligevel
ikke afgive en tilstrækkelig Grund til at udlede den
ethiske Lov af Egoismen; Spørgsmaalet vilde da
netop være, om der ingen Vej gaves til at komme ud
over Egoismen. Men man maa ikke stirre sig blind
paa Exempler af Brutalitet, tilmed da man overfor
dem vilde kunne stille Exempler paa heltemodig og
uegennyttig Opofrelse.\footnote[1]{Den 23de December 1871 brændte Dampskibet „Amerika“ paa Rejsen fra Buenos Ayres til Montevideo. Kaptajnen var strax sprungen overbord med et Redningsbælte om Livet; Mandskabet reddede sig i en Baad. Passagererne havde kun Valget imellem Døden i Flammerne eller i Bølgerne. Man sloges med Dolk og Revolver om Redningsbælterne. En Usling dræbte en ung, hjælpeløs Kone for at bemægtige sig hendes Redningsbælte, medens en anden Passager, Louis Viala, frivillig afstod sit til en ung Kvinde. Uslingen blev reddet, Viala omkom. (Berlingske Tidende 3dje Febr. 1872).} Kampen for Tilværelsen
vil væsenlig sige Kampen for at bevare den bestemte
Form for Tilværelsen, man har opnaaet. Jo mere
Udviklingen skrider frem, desto mindre er det den
nøgne Existens, det drejer sig om; Menneskene have
deres Liv i visse \emph{Former}, som de ikke kunne eller
ville opgive, Former, der staa som Udtryk for det
Menneskelige, og hvis Bevarelse og Udvikling er
Kampens Formaal. Kampen gjælder Hævdelse af
% --- p. 32 ---
\setcounter{footnote}{0}%
\opage{32}alt det Sande, Skjønne og Gode, Menneskeslægten
under sin Livsgang har frembragt. Den, som opofrer
sig for et idealt Formaal, sin Kunst eller sin Videnskab,
eller Den, der sætter sit Liv til for at redde en
Anden, kæmper lige saa vel for Livet som Egoisten,
der ofrer Alt og Alle for sin egen Fordel. Nationaløkonomerne
tale om en tilvant Levevis, en «standard
of life», som Indbegreb af de Fordringer, Arbejderne
stille til Livet, og i hvis Niveau de kæmpe for at
holde sig. Ogsaa i ethisk Henseende gives der en
«standard of life», om hvis Bevarelse det egentlig
er, at Kampen drejer sig. Mennesket kæmper stedse
for at hævde, hvad der for ham er det Højeste, og
hvortil han har knyttet sin hele Personlighed. Man
har ikke Ret til kun at se hen til de mest elementære
og brutale Former for Livskampen. For at
lære Menneskets Væsen at kjende, maa man ikke
blot studere det paa et enkelt Udviklingstrin, men
følge det paa alle forskjellige Trin, og i den Lov, der
gjælder for Udviklingsgangen fra det laveste af disse
Trin til det højeste, vil man have det fulde Udtryk
for den menneskelige Natur. Kaalormen og Sommerfuglen
er det samme Dyr paa forskjellige Stadier,
og kun i Loven for Udviklingen fra det ene Stadium
til det andet har man dette Dyrs Natur udtrykt.

2. Overfor Forsøgene paa at udlede det Ethiske
af Lyksalighedsdriften stille vi nu de Forsøg, man
har gjort paa at udlede det af den naturlige Sympathi.
I Oldtiden fremdrog \emph{Stoikerne} dette Udgangspunkt,
idet de viste, hvorledes Menneskekjær%
% --- p. 33 ---
\setcounter{footnote}{0}%
\opage{33}ligheden har sin Kilde i Forældrenes Kjærlighed til
Børnene, hvorfra den efterhaanden udbreder sig i
videre og videre Kredse, omfatter hele Folket, ja
tilsidst hele Menneskeslægten. Lige som vi bruge
vore Lemmer førend vi have lært, til hvilken Nytte
vi have faaet dem, saaledes have vi fra selve Naturens
Haand modtaget en Drift til Samfund med
Andre. I nyere Tid er denne Lære bleven videre
udviklet af \emph{Shaftesbury}, Hume og \emph{Adam
Smith}, og i vore Dage af \emph{Auguste Comte}, \emph{Darwin}
og \emph{Spencer}. Hos de to sidst nævnte er den
bragt i Forbindelse med Udviklingshypothesen.

Naar Mennesket umiddelbart og instinktmæssig
deltager i Andres Vel og Ve og strax fra Fødselen
af er Led af et Samfund, saa vil ogsaa Samfundets
Bevarelse staa som Formaal for den Enkelte. Allerede
hos Dyrene finde vi stærke sociale Drifter.
Darwin anfører efter Brehm det Træk, at da en
Flok Bavianer ved Stenkast blev jaget paa Flugt af
Rejsende, vendte nogle af de kraftigste Hanner flere
Gange om for med Livsfare at redde de Unger, der
ikke kunde følge med. Saadanne sociale Drifter udvikles
ved Kvalitetsvalget. De Dyr, der hjælpe
hverandre indbyrdes, faa en Fordel i Livskampen
overfor dem, der kæmpe isolerede. Ved Nedarvning
og stadig Øvelse udvikles de sociale Interesser stedse
videre; og da Slægtens Bestaaen beror paa dem,
ville de i Reglen være de stærkeste; de virke stille,
men uafbrudt, selv om de i enkelte Øjeblikke trænges
tilbage af den egoistiske Selvopholdelses- og Lyksalighedsdrift.
Naar Individet, efter saaledes at være
% --- p. 34 ---
\setcounter{footnote}{0}%
\opage{34}bukket under for Fristelsen, er kommet i sin normale
Tilstand igjen, vil det mærke, at et varigere
og dybere begrundet Instinkt er blevet hindret i sin
Virksomhed. Allerede hos Dyret kan man derfor
tale om Anger; dog kan kun Mennesket kaldes et
moralsk Væsen, fordi han har Evne til at anstille
Sammenligninger mellem sine forbigangne og fremtidige
Handlinger eller Motiver og billige eller misbillige
dem. Man kan ofte hos Dyrene iagttage en
Kamp mellem Instinkter, f. Ex. hos Hunden mellem
Lydighed mod Herren og Kjærlighed til Hvalpene,
hos Svalen mellem Moderkjærligheden og Vandredriften.
«Saa længe Hunsvalen mader sine Unger
eller ruger over dem, er Moderkjærligheden sandsynligvis
stærkere end Vandreinstinktet; men det
Instinkt, som er mest udholdende, vinder Sejr, og
tilsidst, naar hendes Unger et Øjeblik ere ude af Syne,
tager hun Flugten og forlader dem. Hvilke Samvittighedsskrupler
vilde ikke hver enkelt Fugl føle,
dersom den, naar den var kommen til Ende med
sin Rejse, og Vandreinstinktet ikke længere virkede,
var begavet med en saa stor Sindslivlighed, at den
ikke kunde forhindre, at der forbi dens Sjæl stadig
drog et Billede af dens Unger, døende i det vinterkolde
Nord af Hunger og Kulde!»\footnote[1]{\emph{Darwin}: Om Menneskets Oprindelse. Dansk Overs. I. p. 86.} Menneskets
intellektuelle Evner ere nu netop saaledes udviklede,
at tidligere Indtryk og Billeder stadig drage forbi hans
Sjæl, saa at han ikke kan undgaa at anstille Sammen%
% --- p. 35 ---
\setcounter{footnote}{0}%
\opage{35}ligninger; kun han kan derfor føle virkelig Anger;
den vækker hos ham en skærpet Følelse af Pligt,
det vil sige af Lydighed mod det dybest liggende
og stadig virkende Instinkt i hans Natur; at der
er Noget, Mennesket \emph{bør}, vil sige, at han har en
Bevidsthed om Tilværelsen af et vedholdende Instinkt,
der tjener ham som Fører, men rigtig nok er
udsat for ikke at blive adlydt.

Lige som ifølge Benthams Lære ikke enhver Tilfredsstillelse
af Lyksalighedsdriften kunde kaldes
ethisk, men kun den, hvor Mennesket fulgte Forstandens
kloge Beregning trods den øjeblikkelige Indskydelse,
saaledes kunne Sympathimoralens Tilhængere heller
ikke betragte enhver Handlen af sympathetisk Drift
som ethisk. Tacitus fortæller om de gamle Germaner,
at skjønt de glædede sig ved Foræringer,
mente de dog hverken at skylde Taknemmelighed,
naar de havde modtaget dem, eller at have Krav
derpaa, naar de havde givet dem. Og ligeledes fortælles
det om nogle indianske Stammer, at de handlede
af Hengivenhed lige som de handlede af Attraa,
uden Hensyn til Følgerne. Naar de havde vist en
Anden en Venlighed, havde de tilfredsstillet en
Drift; Sagen var endt og forsvandt fra Erindringen;
men de følte omvendt heller ikke Fortrydelse over
at have tilfredsstillet deres selviske Lidenskaber;
baade i deres Sympathi og i deres Egoisme handlede
de som sande Øjeblikkets Børn. Atter her
viser det sig, at det Ethiske forudsætter et Forbillede,
en Lov, som ligger ud over Øjeblikkets Stemning
og peger hen mod en mere omfattende og fast Form
% --- p. 36 ---
\setcounter{footnote}{0}%
\opage{36}for Menneskelivet. Derfor fremhæve Darwin og
hans Forgængere Erindringens og Tankens Betydning.
Den er hos dem ikke i Lyksalighedsdriftens, men i
Sympathiens Tjeneste, fastholder dennes Krav overfor
Selvopholdelsesdriftens og Egoismens Fordringer.
Det personlige Princip er her underordnet det sociale.
Ved Erkjendelsens lutrende og lysende Magt
kunne Sympathierne efterhaanden befries fra deres
Skranker, gaa ud over Familiens og Nationens snævrere
Kredse og blive en uinteresseret Kjærlighed til
alle levende Væsener.

Ogsaa Bentham og Utilitarianerne kunde henvise
til en stor, social Sammenhæng, i hvilken Individet
stod i ethisk Henseende. Alt hvad der hører til
det menneskelige Livs Opholdelse og til Tilfredsstillelsen
af dets Fornødenheder, bringes ikke tilveje
paany indenfor hver enkelt Generation, men hver
ny Slægt arbejder paa Grundlag af faste Institutioner,
Former og Overleveringer, som den arver fra
tidligere Slægter. Kulturen beror paa Tradition.
Ikke blot ved sin Fornuft, sin Evne til at kombinere
og forudse, hæver Mennesket sig over Dyret;
Fremskridtet vilde ikke blive stort, naar Enhver
skulde begynde forfra. Sproget og Traditionen gjøre
Ophoben og Udbredelse af Opfindelser og Opdagelser
mulig. Det System af Interesser, hvortil Utilitarianismen
henviser, omfatter altsaa ikke blot alle
samtidig levende Mennesker, men ogsaa den talløse
Række af tidligere Slægter. Som Kulturvæsener ere
vi solidarisk forbundne med vore Forgængere. Men
ifølge Nyttemoralen er Kulturen kun en ydre Sam%
% --- p. 37 ---
\setcounter{footnote}{0}%
\opage{37}menhæng, i hvilken Individet lever sig ind, fordi
han finder sine egne Interesser tilfredsstillede ved
at virke hen til dens Befæstelse og Udvikling. Bentham
og hans Skole ser ikke hen til den Indflydelse,
Kulturen ifølge Livets almindelige Love har og maa
have paa selve Menneskenaturen. Den menneskelige
Natur udvikles og omformes gjennem det kulturhistoriske
Arbejde. Med de forandrede Livsvilkaar
ville ogsaa efterhaanden de oprindelige Drifter og
Anlæg forandres. Ikke blot gjennem ydre Tradition
faar den Enkelte Del i, hvad Slægten har erhvervet;
i sin inderste Natur har han Eftervirkningerne
af Fortidens Anstrengelser, ifølge den almindelige
fysiologiske Lov, at Funktionen virker tilbage paa
Organisationen. Menneskenes oprindelige Vilkaar
maatte gjøre dem til Egoister; de mildere Former
Kampen for Livet efterhaanden antager, maa ogsaa
formilde Menneskenaturen. Udviklingshypothesens
store Talsmænd, Darwin og Spencer, træde derfor
op mod Bentham, Mill og Bain.\footnote[1]{\emph{Herbert Spencer} skrev i Aaret 1846 en Afhandling om Sympathien i samme Retning som Adam Smiths „Theory of moral sentiments“, som han dengang ikke kjendte. 1850 udkom hans „Social Statics“, hvor det samme Emne behandles. Ogsaa i Slutningen af hans Psykologi kommer han tilbage dertil, men hans „Principles of morality“ skulle først danne Slutningen af hans store Værk.} Medens Utilitarianerne
forudsætte Lyksalighedsdriften som lige
stærk til alle Tider, er Sympathimoralens Udgangspunkt
ikke en konstant Størrelse, men varierer
stadig. Sit dybeste Grundlag har den sympathetiske
% --- p. 38 ---
\setcounter{footnote}{0}%
\opage{38}Drift, som flere Gange antydet, i Moderinstinktet.
Auguste Comte finder derfor det første Spor til
det moralske og sociale Liv paa det Punkt indenfor
Dyreverdenen, hvor Hermafroditismen er hørt op,
og Kjønnene adskilte. Fra Familieforholdet udbreder
det sig nu i videre og videre Kredse. Den
rent empiriske Opfattelse kan ikke tilstrækkelig
begrunde denne Sympathiens stille Væxt. Vel
har Hartley, og efter ham i vore Dage Stuart Mill,
med Rette fremhævet, at i Kraft af Ideassociationens
Love vil Mennesket komme til at elske det for dets
egen Skyld, som han først kun interesserede sig for,
fordi det tilfredsstillede hans Egenkjærlighed, lige
som vi uvilkaarlig tilskrive Penge en Betydning og
et Værd, som egenlig kun tilkommer det, vi kunne
opnaa for Pengene. Den, som stedse nødes til at
arbejde sammen med Andre for at opnaa sine egne
Formaal, vil efterhaanden glemme sit oprindelige
Motiv og føle en uegennyttig Interesse for Andres
Vel og det fælles Maal. Opdragelsens Kunst bestaar
netop i at knytte den Enkeltes Lyst og Smerte til
de rette Gjenstande, binde hans Interesse saaledes
til dem, at han tilsidst ganske kan gaa op i dem.
Men Udviklingshypothesen lærer os, at det ikke er
Opdragelsen og den individuelle Erfaring, som her
skal gjøre Alt. Den sympathetiske Drift er ikke
blot Produkt af Ideassociationerne. Den voxer gjennem
Slægterne, idet den stedse hyppigere Lejlighed
til at ytre sig og den stedse mere omfattende Kreds,
indenfor hvilken den rører sig, nødvendigvis ogsaa
maa forstærke Anlæget hos Individerne til at lade
% --- p. 39 ---
\setcounter{footnote}{0}%
\opage{39}sig beherske af den. Stille og umærkelig foregaar
der altsaa ifølge denne Hypothese en Metamorfose
af den menneskelige Natur. Enhver begejstret
Stemning, enhver ædel og menneskekjærlig Følelse
føjer et Atom til Bygningen af den fuldkomne
Menneskehed. Intet gaar til Spilde lige saa lidt
i den aandelige som i den legemlige Natur.

Den ethiske Lov, som den former sig paa Sympathimoralens
Standpunkt, vilde have til Indhold
Udviklingen af den fuldendte, omfattende Menneskekjærlighed.
Men det Spørgsmaal kan ikke holdes
borte, hvori det Ethiske i Sympathien egenlig ligger.
Som umiddelbar Drift kan denne lige saa vel knytte
sig til det Værdiløse eller endog Fordærvelige, som
til det Gode og Rette. Adam Smith bemærker selv
(hvad der maaske ikke længere vilde passe i vore
Dage), at den store Masse er uinteresseret Beundrer
af alt Usædvanligt, især stor Magt og Rigdom.
Kjærlighed er jo i det Hele blind. For at gjøre
Andre lykkelige maa den have et Forbillede, et Ideal
af Lykke, af fuldendt menneskeligt Liv, til hvilket
den saa kan søge at føre Andre. Men dette kan
den kun have, naar Tanken er mere end Sympathiens
Stedfortræder og virkelig kan optræde som ordnende
og ledende Princip. Kjærligheden maa underordne
sig Retfærdigheden (der i Følge Leibnitz er caritas
sapientis, \emph{den Vises} Kjærlighed), som det ledende,
fordelende og harmoniserende Princip i det menneskelige
Samfund.

Adam Smith har gjort et interessant Forsøg paa
at vise, hvorledes Sympathiens Indflydelse fører til
% --- p. 40 ---
\setcounter{footnote}{0}%
\opage{40}Dannelsen af en ethisk Maalestok. Idet vi umiddelbart
føle Lyst eller Smerte ved Andres Handlinger,
selv om de ikke angaa os, komme vi ogsaa til at udtale
vor Billigelse eller Misbilligelse af dem. Vi
lide med Ofret for en grusom eller troløs Handling
og føle Antipathi mod den, der begaar en saadan
Handling. Paa denne Sympathi og Antipathi beror
vor Billigelse og Misbilligelse, som altsaa opstaar
aldeles instinktmæssig, idet vi uvilkaarlig
sætte os i de Andres, den Handlendes og den
Lidendes Sted. Vi agere Tilskuere, som føle
med de i Dramaet optrædende Personer. Vore første
moralske Domme fælde vi saaledes om Andre. Dog
kommer der et Tidspunkt, da vi se, at vi selv ere
handlende Personer i Livets Drama lige saa vel
som de Andre, hvorfor vi maa antage, at Andre ikke
mindre bedømme os, end vi dem. Og selve Sympathien,
som søger at udjævne Forskjellen mellem
vor og Andres Følelsestilstand, idet enhver Disharmoni
her er forbunden med en pinlig Følelse, vil
føre os til at undersøge vort Væsen og vore Handlinger
for at se, med hvilke Øjne vi kunne vente,
at Andre betragte os. Vi sætte os altsaa atter i
Andres Sted, ikke længere blot for at glæde os eller
lide med dem, men for fra deres Standpunkt at
beskue og vurdere os selv. Enhver deler sig saaledes
i to Personer, en dømmende Tilskuer og et
handlende Væsen. Det naturlige Stade forlades; Individet
ser bort fra sig selv, netop for at lære sig selv at
kjende. Det er en Erfaringskjendsgjerning, at den,
som i Andres Nærværelse kommer i lidenskabelig
% --- p. 41 ---
\setcounter{footnote}{0}%
\opage{41}Bevægelse, uvilkaarlig nedstemmer denne, naar han
føler, at de Andre, som upartiske Tilskuere, ikke
dele den. Herpaa beror det opdragende Element
i Samlivet med Andre. Enhver føler sig som Gjenstand
for en stadig og opmærksom Kritik og indretter
uvilkaarlig sin Handlen saaledes, at han kan
bestaa for dens Domstol. Barnet, Ynglingen og
Manden lære saaledes i de forskjellige Livsforhold
at se bort fra det individuelle Standpunkt. Tilsidst
behøver den indre Tilskuer og Dommer i vort Bryst
ikke at styrkes ved Nærværelsen af ydre Tilskuere;
vi have gjennemgaaet Selvbeherskelsens store Skole
og have vor Regel, vort Forbillede i os selv. Vi
se nu ogsaa bort fra den endelige og erfaringsmæssige
Maalestok, som er given i vore Omgivelsers
Meninger og Følelser, og vende vort Blik mod et
Fuldkommenhedens Ideal, som efterhaanden har
dannet sig i vor Forestilling. Den ydre Tilskuer
bliver altsaa først til en indre Tilskuer, og denne
igjen til et idealt Forbillede.

Naar Udviklingen har naaet dette Punkt, træder
ifølge Adam Smith det oprindelige Udgangspunkt
helt i Skygge. Baade de sympathetiske og de egoistiske
Drifter træde til Side for en højere Følelse.
«Naar vi stedse berøres langt stærkere af, hvad der
angaar os selv, end af hvad der angaar Andre,
hvad er det da, som ved enhver Lejlighed faar den
ædle Mand og undertiden endog den lave Karakter
til at ofre deres egne Interesser for Andres større
Interesser? Det er ikke Menneskekjærlighedens
milde Magt, ikke den svage Gnist af Velvillie, som
% --- p. 42 ---
\setcounter{footnote}{0}%
\opage{42}Naturen har tændt i Menneskets Hjerte, der saaledes
er i Stand til at modvirke Egenkjærlighedens
stærkeste Impulser. Det er en stærkere Magt, et
kraftigere Motiv, der virker ved saadanne Lejligheder.
Det er Fornuften, Tanken, Samvittigheden,
det indre Menneske, som bor i vort Bryst, den
store Dommer og Vurderer af al vor Adfærd. Den
er det, der, naar som helst vi ere i Færd med at
handle saaledes, at vi hindre Andres Lykke, tilraaber
os med en Stemme, som kan overdøve vore
mest fordringsfulde Lidenskaber, at vi kun ere En
af Mange, i ingen Henseende bedre end enhver Anden,
og at naar vi saa frækt og blindt foretrække os
selv for Andre, blive vi med Rette Gjenstand for
Harme, Afsky og Forbandelse. Af den lære vi
vor virkelige Lidenhed og Lidenheden af Alt, der
angaar os, og kun ved denne upartiske Tilskuers
Øje kunne Egenkjærlighedens naturlige Illusioner
berigtiges. Den viser os Højsindets Skjønhed og
Uretfærdighedens Hæslighed, det Rigtige i at opgive
vore egne højeste Interesser for Andres endnu højere
Interesser og det Hæslige i at paaføre Andre
den mindste Uret for at opnaa den højeste Fordel for
os selv. Det er ikke Kjærligheden til vor Næste eller
til Menneskeslægten som ved saa mange Lejligheder
ansporer os til at udøve disse guddommelige Dyder.
Det er en stærkere Kjærlighed, en mægtigere Følelse,
som i Almindelighed ytrer sig ved saadanne Lejligheder:
Kjærligheden til hvad der er ædelt og hæderligt,
til vor Karakters Storhed, Værdighed og Højhed.»\footnote[1]{Theory of moral sentiments. Vol. 1. part. 3. chap. 3.}

% --- p. 43 ---
\setcounter{footnote}{0}%
\opage{43}Ved denne Fremstilling har Adam Smith anerkjendt,
at den sympathetiske Drift i sig selv kun er
et Element med ved Dannelsen af den ethiske Lov,
men ikke indeholder den fulde Forklaring af denne,
og navnlig ikke fuldelig udtrykker den Følelse, det
Sindelag, som det ethiske Forbillede vækker. Man
vil tillige let se, at Sympathi er for svagt og lidet omfattende
et Udtryk om Forholdet til de opdragende
Magter i Menneskets og i Slægtens Historie. Sympathi
forudsætter et Forhold mellem Ligemænd,
men den opdragende Magt staar i det Mindste oprindelig
som Overmand. For saa vidt have Hobbes
og Bentham mere Ret, naar de henvise til Frygt
og Haab som de Følelser, der røre sig overfor Avtoriteterne.
Dog, denne Betragtning maa vi holde
tilbage, til vi komme til den særlige Undersøgelse
af Avtoritetens Begreb. Derimod opfordrer Smiths
Lære om «den upartiske Tilskuer» til at undersøge
Erkjendelsens, Fornuftens Betydning i ethisk Henseende.
Den spiller, som vi have set, en stor Rolle
baade i Nyttemoralen og i Sympathimoralen, men i
dem begge som underordnet Princip. For bedre at
opfatte dens Betydning vende vi os nu til de Forsøg,
man har gjort paa at gjøre den selv til det
Ethiskes Grundlag.

3. Baade Nyttemoralen og Sympathimoralen indrømme,
at de umiddelbare Drifter ikke kunne kaldes
ethiske, men at der først bliver Tale om ethisk
Handlen, naar der har dannet sig en Bevidsthed om
en Lov og Regel, som viser ud over Øjeblikkets
% --- p. 44 ---
\setcounter{footnote}{0}%
\opage{44}umiddelbare Tilskyndelser. En saadan Bevidsthed
forudsætter Erkjendelse, Fornuft, Tænkning. Derfor
er det kun i Menneskeverdenen, at det Ethiske hører
hjemme; den Udviklingsgrad af Bevidsthedslivet,
som det forudsætter, findes endnu ikke i Dyreverdenen,
og er sikkert ogsaa først efter lange Tiders
Forløb opnaaet af Menneskene. Det har varet længe,
inden Menneskelivet er begyndt at ledes af Tanker,
og vi ere endnu ikke komne ud over Begyndelsen.
--- Naar vi her stille Tanken eller Fornuften som
ethisk Udgangspunkt ved Siden af Selvopholdelsesdriften
og Sympathien, opfatte vi den ikke længere
som Tjener, men som Herre, og gaa ud fra den
Kjendsgjerning, at den bevidste Erkjendelse under
visse Betingelser kan bestemme Menneskets Handlen.
Spørgsmaalet er da, om dette er en udtømmende
Forklaring af det Ethiske.

Men vi standses her strax af en Betænkelighed.
Ifølge vore almindelige psykologiske Overvejelser
kan Erkjendelsen i sig selv ikke sætte Villien i
Bevægelse. Tanken bliver kun Motiv ved at fremkalde
Følelse og Drift. Fordi jeg aldeles tydelig
og klart har indset Noget, følger ikke, at jeg handler
derefter; det gjør jeg kun, naar Indsigten vækker
levende Følelse og Attraa hos mig. Der vil ganske
vist være en Tendens til at bringe Harmoni mellem
de forskjellige Elementer i Bevidsthedslivet, men
denne Tendens kan i sig selv være lige saa instinktmæssig
som Lyksalighedssdriften og Sympathien. % PRINTED AS IS: „Lyksalighedssdriften“ (ss)
Tænke vi os en klart og alsidig udviklet Fornuft,
i Overensstemmelse med hvilken Villien umiddelbart
% --- p. 45 ---
\setcounter{footnote}{0}%
\opage{45}virker, vil et saadant Fænomen dog ikke henføres
til det Ethiske. Dette viser sig atter her at forudsætte
en Afstand, fordi det stiller en Opgave, som
skal løses, en Afstand, som udfyldes ved Forestillingen
om en Lov, hvis Gyldighed Mennesket anerkjender.
Det vil derfor ogsaa vise sig, at naar man
opstiller Fornuften som ethisk Princip, er det egenlig
ikke Fornuften i og for sig, man mener, men Fornuften
som Udtryk for Menneskets sande Væsen,
saa at Fordringen gaar ud paa, at Menneskenaturens
Adel og Værdighed sker Fyldest.

Det almindelige Thema for dette Standpunkt har
\emph{Pascal} angivet i en berømt Udtalelse. «Mennesket
er det svageste Siv i Naturen, men han er et tænkende
Siv; Universet behøver ikke at ruste sig for
at knuse ham; en Smule Røg, en Draabe Vand er
nok til at gjøre det af med ham. Men selv om Universet
knuser ham, er han dog ædlere end sin Morder,
fordi han véd, at han dør, medens Universet selv
ikke kjender den Magt, det har over ham. Hele
vor Værdighed bestaar altsaa i Tanken. I Tanken
maa vi sætte vor Stolthed, ikke i Tiden eller i
Rummet, som vi ikke kunne udfylde. Lad os da
arbejde paa at tænke rigtig; det er Moralens Princip.»
Hvad Pascal her udtrykker, er Begejstringen over
at høre til Aandens Verden. Tanken giver Mennesket
en Værdighed, som ubevidste Væsener ikke
have.

Allerede hos \emph{Sokrates} finde vi dette Princip.
Han lærte, at Dyd er Viden, fordi kun den Handlen
er ret, som udspringer af virkelig Indsigt. Fremfor
% --- p. 46 ---
\setcounter{footnote}{0}%
\opage{46}Alt kommer det an paa at kjende sig selv; thi kun
Den, som kjender sine egne Evner, Kræfter og Drifter,
kan optræde paa rette Maade i Livet. Den
nærmere Begrundelse, Sokrates giver heraf, gaar i
Retning af Nyttelæren, og Stuart Mill har saaledes
kunnet tage ham til Indtægt for denne. Vigtigheden
af ret Selverkjendelse indskærper Sokrates nemlig
ved at henvise til Uvidenhedens og Forblindelsens
uheldige Følger. Men selv gjennem den mest prosaiske
Gjengivelse af Sokrates' Lære fremgaar det
dog, at han egenlig har lagt Vægten paa et andet
Punkt. Han hævdede, at «Intet var stærkere end
Indsigt», at den var «noget Skjønt, som er i Stand
til at herske over Mennesket, stærk nok til at hjælpe
ham gjennem Livet.» Den, som ikke lader sig lede
af Erkjendelsen, men af blinde Lyster og Drifter,
er slavisk eller dyrisk. En Slave er den, som ikke
er delagtig i det Skjønne og Gode. «Synes det Dig
ikke skjændigt», siger han til Aristip, «at et Menneske
vil bære sig ad lige som de mest ufornuftige
Dyr, der drives af den sanselige Lyst til Lokkemaden
og blindt gaa i Fælden?» I sin Forsvarstale
henviser han ligeledes til Tanken og dens Magt som
Udtrykket for Menneskets sande Væsen, hvis Pleje
maa ligge ham nærmere end Omsorgen for det Legemlige
og de ydre Goder. «Min Ven», saaledes
tiltaler han sin Landsmand, «Du som er en Athenienser,
Borger i en stor Stad, den for Dannelse og
Kraft mest berømte af alle, skammer Du Dig ikke
over kun at bryde Dig om at samle Dig Rigdom
og erhverve Dig Ære, medens Du ikke bekymrer
% --- p. 47 ---
\setcounter{footnote}{0}%
\opage{47}Dig om Sandhed og Visdom, om din Sjæl og dens
Opdragelse!»

Allerede her, hos selve Ethikens Grundlægger,
se vi altsaa, at naar Fornuften opstilles som ethisk
Princip, menes derved egenlig den indre Frihed og
Harmoni i Menneskets Væsen, som kun er mulig,
naar hans Handlen er i Overensstemmelse med hans
højeste Erkjendelse. Mennesket drages da ikke til
forskjellige Sider af modstridende Lyster og Drifter;
han ledes heller ikke blindt af Indtrykkene fra Omverdenen;
der kommer Enhed, Sammenhæng og
Selvstændighed i hans Væsen. Han bliver en lille
Verden for sig selv, uafhængig overfor den ydre
Verden, og indeslutter i sig selv en harmonisk Fylde.

Hvad Sokrates i al Simpelhed havde antydet,
vedblev at være Hovedtanken i den græske Ethik.
Platon, Aristoteles og Stoikerne opfattede hver paa
sin Maade Menneskets ethiske Handlen som bestaaende
deri, at han virker efter sit Væsens højeste
Lov, idet den fornuftige Indsigt raader. De finde Friheden
i Overensstemmelse med den rette Indsigt,
og i denne Frihed den sande Menneskelighed. Netop
denne umiddelbare Overbevisning om Tankens Herredømme,
som bragte Sokrates til med Indignation at
forkaste den Antagelse, at «Noget kunde være stærkere
end Indsigt», udgjorde den græske Ethiks Ensidighed.
Den oversaa det mangfoldig Betingede
og Begrænsede i det menneskelige Aandsliv og alle
de Modsætninger og Modsigelser, Kampe og Nederlag,
som heraf følge. For den smeltede Villien dels
% --- p. 48 ---
\setcounter{footnote}{0}%
\opage{48}sammen med Erkjendelsen, dels med den umiddelbare
Drift.

I nyere Tid drog Kant atter den indre Friheds
Ide frem, idet han pegede mod Fornuftens Lov
som den, der ubetinget forpligter Mennesket. Han
viste klart, at den Lov, det Ethiske forudsætter,
ikke kan være en i Naturen umiddelbart given, men
maa skyldes en anden Kilde end den, af hvilken
Menneskets uvilkaarlige Handlinger udspringe. Den
ethiske Lov kan ifølge Kant ikke begrundes ---
hverken psykologisk, fysisk eller theologisk begrundes.
Den viser ud over alt Givet, over al Erfaring
og Virkelighed, idet den foreskriver, hvad der skal
ske. Kun i Fornuften have vi en saadan ud over
al Erfaring pegende Magt. Fornuftens ubetingede
og ubegrundelige Bud staar som en Kjendsgjerning
for Bevidstheden, en umiddelbar Fordring, hvorunder
Mennesket maa bøje sig. Det vilde miste Karakteren
af ubetinget Bud (kategorisk Imperativ), dersom
det støttedes paa anden Maade end ved Fornuftens
umiddelbare Fordring.

Fremskridtet fra Sokrates til Kant ses let. Sokrates
formaar endnu kun at begrunde Indsigtens
Værd ved at henvise til dens ydre Nytte; denne
Begrundelse staar hos ham Side om Side med Appellen
til Menneskets Værdighed som frit Fornuftvæsen.
Desuden er Tanken om en Modsætning, en
Dobbelthed indenfor Menneskets Væsen Sokrates
fremmed. Med den rette Indsigt er Alt givet; Tanken
optræder ikke som Lov, fordi der forudsættes
en umiddelbar Harmoni mellem den og de andre
% --- p. 49 ---
\setcounter{footnote}{0}%
\opage{49}sjælelige Elementer. Kant derimod hæver Fornuften
som ethisk Princip over alt umiddelbart Givet og % PRINTED AS IS: „og|og“ (doubled over the line end)
og afviser enhver erfaringsmæssig Begrundelse. Og
han bestemmer Fornuften som lovgivende, hvorved
der forudsættes en Afstand mellem ledende og lydende
Elementer i Menneskets Væsen.

Man vilde ikke ganske have Uret i at udlede
denne Modsætning mellem Sokrates og Kant af den
store religiøse Bevægelse, som ligger mellem disse
to største Skikkelser i den humane Ethiks Historie.
Kristendommen opstillede et fra en overnaturlig
Verden hentet Ideal, ophøjet over Alt, hvad menneskelig
Tanke skulde kunne naa ved egne Kræfter.
Synes vi ikke at høre en Gjenklang heraf,
naar ifølge Kants Lære Fornuftloven forsmaar al
fra Virkeligheden hentet Begrundelse? Kristendommen
erklærede den menneskelige Natur for syndig
og fordærvet. Mindes vi ikke herom ved Kants
Forudsætning om Afstanden mellem Fornuftens Bud
og Menneskets virkelige Handlen?

Dog vil den fuldstændige Forklaring ikke hermed
være given. Saa indgribende hin religiøse Bevægelse
end har virket paa det evropæiske Aandsliv,
saa maa man dog ikke overse, at den selv kun er en af
Formerne for den almindelige Modsætning mellem
Oldtiden og den nyere Tid. For Grækeren stod
Livet som umiddelbart givet og afsluttet. Historien
fortalte ham endnu ikke tydelig og indtrængende
om store Konflikter mellem forskjellige Trin af
menneskelig Udvikling; han havde ikke den Opfattelse
af Livet som en Kamp for Tilværelse og
% --- p. 50 ---
\setcounter{footnote}{0}%
\opage{50}Fremgang, som den nyere Tid har vundet gjennem
sine Erfaringer, sin Handlen og Forsken. Naar
Livet i Stedet for at staa for Forestillingen som et
Kunstværk i en begrænset Ramme, viser sig at
være Bevægelse og Udvikling, maa Forestillingen
med en ganske anden Energi vende sig frem ad mod
det, der endnu ikke er, men skal være. Tanken om
en Opgave bliver nu ganske anderledes levende.
Tanken iler forud, foregriber paa ideal Maade det
nye Livstrin og fælder Dommen over det tidligere
Trin. Den ideale Anskuelse, som saaledes danner
sig, har ikke sin fulde Forklaring i det umiddelbart
Givne og Foreliggende. Vi vises her hen til Bevidsthedslivets
oprindelige Energi, hvorved selv den
simpleste Sansefornemmelse bliver mere end en
passiv Modtagen, og hvorved ethvert friere Forhold
til Sansningens og Erfaringens Indhold og enhver
Bearbejdelse deraf bliver mulig. Det er denne
Energi, som gjør det muligt at danne fuldstændige
og afsluttede Billeder, Idealer, af de altid ufuldstændige
Elementer i den virkelige Verden. Her
er Kilden til al Poesi og Kunst, til alle højere
Ideer. Vil man kalde denne Evne hos Mennesket
til at danne ideale Forestillinger Fornuft, saa er
det fuldstændig rigtigt, at kun ved Fornuften som
den, der gaar ud over det Givne, bliver den ethiske
Lov mulig, og at denne Lov altid staar ophøjet over
den virkelige Handlen som et Forbillede, hvortil
den stræber at nærme sig.

Men naar vi tage Fornuft og Fornuftlov i denne
Betydning, kunne vi ikke som Kant blive staaende
% --- p. 51 ---
\setcounter{footnote}{0}%
\opage{51}ved Fornuftens Bud som blot Kjendsgjerning; vi
maa undersøge, hvorledes Forestillingen om en saadan
Lov psykologisk og historisk har udviklet sig.
Denne Opgave kunne vi ikke unddrage os uden at
krænke den videnskabelige Forskens simpleste Regler.
«Kjendsgjerningerne» paa det aandelige og
ethiske Omraade maa lige saa vel underkastes Prøvelse
og Begrundelse som alle andre Kjendsgjerninger.
Vi maa her paakalde den kritiske Kant mod
den dogmatiske Kant; og han kommer os her selv
i Møde ved sin Fremstilling af, hvorledes Fornuftloven
bestemmer Menneskets Handlen. Kun gjennem
Følelsen kan, som vi vide, Fornuften vække
Villien. Fornuftloven staar ved sin ubetingede
Myndighed som en over Individet ophøjet Magt.
De sanselige Drifter trænges tilbage, den umiddelbare
Livsfølelse og Selvtilliden hæmmes og ydmyges
ved Tanken om denne ubetingede Lov. Mennesket
føler overfor Loven i sit Indre lige som overfor den
klare Stjernehimmel sin Lidenhed paa samme Tid,
som han fyldes med Beundring og Ærefrygt; thi
den ethiske Lov virker kun frastødende paa den
sanselige og egoistiske Drift for at øve saa meget
større Tiltrækning paa Menneskets egenlige Jeg,
som føler sig Et med Loven. I denne finde vi
derfor et Udtryk for vor sande Villie. Men hvad
man end mener om Fornuften, --- en Følelse som
Beundringen og Ærefrygten maa jo dog have sin
Naturhistorie, og vi have ikke ret forstaaet det
Ethiske, førend vi have udforsket den.

Da Kant betragter den ethiske Fornuftlov som
% --- p. 52 ---
\setcounter{footnote}{0}%
\opage{52}hævet over al Erfaring, alt Givet og Foreliggende,
kan han ikke give den noget bestemt Indhold. Men
i selve Begrebet om en Lov ligger, at den er almengyldig,
ikke retter og læmper sig efter det enkelte
Individ, men er fælles for ham og alle Andre. Den
ethiske Lovs almindelige Indhold er da: Enhver
skal handle saaledes, at den Grundsætning, han følger
i sin Handlen, kan tjene som almindelig Regel
for Alle. Mennesket skal betragte enhver af sine
Handlinger som Led i en højere Natursammenhæng,
der skal virkeliggjøres gjennem hans Stræben.
--- Dette Princip betragter Kant selv som rent formelt
og uafhængigt af Erfaringen, og hans Kritikere give
ham i Regelen Ret heri, kun at de anse det for
en Fejl, han for en Dyd. Men hvorfra, maa man
spørge, faar den rene Fornuft Forestillingen om en
almindelig Lovgivning uden gjennem Erfaringen?
I det virkelige Liv se vi det enkelte Menneske som
Led i Samfundet, i Slægten; vi se bestemte, faktiske
Love gjælde for Samlivet mellem Menneskene, og
vi danne saa ved Abstraktion og Idealisation Forestillingen
om et fuldkomment Samfund, hvor enhver
Enkelt indordner sig med fuldt Hensyn til alle
Andre, saa at en og samme Lov gjælder for Alle.
I det ethiske Princip, Kant opstiller, kunne vi derfor
ikke se noget rent Apriorisk, men en historisk
Ide ført tilbage til ideal Simpelhed. Som vi tidligere
have set (ovenfor S. 4---5) ligger det i vor
Erkjendelses Natur at operere med ideale Begreber,
men disse anvendes i Ethiken paa en anden Maade
end i Naturvidenskaben. Paa Naturens Omraade
% --- p. 53 ---
\setcounter{footnote}{0}%
\opage{53}betragte vi de ideale Begreber som Tilnærmelser til
den virkelige Sammenhæng, paa det ethiske Omraade
derimod det Virkelige som Tilnærmelse til det forestillede
Ideal. Stedse ligger Erfaringen til Grund;
kun fra den kunne Elementerne hentes til de ideale
Konstruktioner.

Naar vi se Sagen paa denne Maade, have vi tillige
Tilknytningspunkter for Nyttelærens og Sympathilærens
Ideer. Kulturen, betragtet som hvilende paa
et System af Interesser, der gjensidig støtte hverandre,
faar sin Plads i Forestillingen om det ideale
Samfund som dettes reale Grundlag. Uden at have
sin faste Grund i Livets primitive Krav svæver
dette i Luften. Overfor enhver asketisk Betragtningsmaade
er det tillige af Betydning at hævde Individets Ret til Selvtilfredsstillelse. Hvad Utilitarianismen
overser, er, at denne Ret ethisk set netop maa
have sin Grund i Individets Værdighed som Medlem
af de menneskelige Personligheders Rige, lige som
jo enhver juridisk Ret peger tilbage til Samfundet
og dets Lov som Grundlag. Sympathilæren har havt
en dybere Forstaaelse af Individets indre Forhold til
Slægten; og den fremhæver med Rette Velvillien
og Menneskekjærligheden som det levende Baand,
der bevarer Samfundet og gjør det til mere end en
mekanisk sammenføjet Dynge af Atomer; men da
vor Anerkjendelse af Loven er et Hovedelement i
den ethiske Handlen, kunne vi ikke blive staaende
ved instinktmæssige Følelser. Saavel Sympathien
som den individuelle Interesse maa underordnes den
højeste Forstaaelse, vi ere i Stand til at vinde af
% --- p. 54 ---
\setcounter{footnote}{0}%
\opage{54}Menneskelivets Grundlov, Formaal og Opgaver.
Virkeliggjørelsen af Humanitetens Rige maa staa
som Formaal og Grundlov for vor Handlen.

Men her har det da atter vist sig, at vi ikke
kunne forstaa det Ethiske, saa længe vi blive staaende
ved Betragtningen af det enkelte Individ. Vi
vises ogsaa her tilbage til den store Sammenhæng,
hvori den Enkelte opstaar og udvikler sig, og som
han skylder hvad han er og har. Hvad enten vi
gaa ud fra Fornuften, fra Sympathien eller fra Lyksalighedsdriften,
føres vi med Nødvendighed ud over
Individualitetens Omraade. Af Samfundets Moderskjød
udspringe de enkelte Individualiteter, og det
er paa den anden Side som Repræsentanter for og
Medarbejdere paa det ideale Samfund af harmoniske
Personligheder, at de naa den højeste Udvikling. Vi
maa derfor begynde en ny Undersøgelse af vort
Spørgsmaal, saaledes at vi ikke længere gaa ud fra
det enkelte Individ men fra den højere Magt, der
fra først af omslutter alle Individer.

% --- p. 55 ---
\setcounter{footnote}{0}%
\chap{Avtoriteten.}
\opage{55}


Jo mere indviklet og sammensat et Emne er,
desto nødvendigere er det at anvende den beskrivende
Methode. Saa længe vi holdt os til det enkelte
Individ og de almindelige psykologiske Love,
som bestemme dets Handlen, vare Forholdene lette
at overskue og analysere, men saa uundværlige de
individuelle Udgangspunkter end ere, har det dog
vist sig, at der ved Analysen fremkom væsenlige
Elementer, som kun finde deres Forklaring, naar vi
tage et socialt Udgangspunkt. Det enkelte Individ
kunde vi i det Væsenlige betragte som en fast Type;
Samfundet, og i det Hele de Magter, i Ly af hvilke
Individet bliver til og udvikler sig, hører derimod
til Historiens rige og mangfoldige Verden, hvor Alt
er Bevægelse og Udvikling underlagt. Ethiken viser
her tilbage til Historien. Det har ikke været den
ringeste Aarsag til Ethikens Ufuldkommenhed, at
man ikke har agtet paa Historiens store Betydning
som Grundlag for den. Navnlig kunde det 18de
% --- p. 56 ---
\setcounter{footnote}{0}%
\opage{56}Aarhundrede, hvis Sønner de største moderne Ethikere
(Ad. Smith, Bentham og Kant) vare, paa Grund
af sin hele abstrakte og individualistiske Retning
ikke tilbørlig vurdere denne Betydning. Det er
først vort Aarhundrede, hvis opladte historiske Sans
ogsaa har ført til Overbevisningen om Vigtigheden
af det historiske Grundlag for Ethiken. Schleiermacher
og Hegel, Comte og Spencer vise sig hver
paa sin Maade gjennemtrængte af denne Overbevisning.

Det har i det Foregaaende vist sig, at man ikke
kan bestemme den ethiske Handlens Natur uden ved
at henvise til Forestillingen om en Lov; men Lovens
Begreb hente vi oprindelig fra det menneskelige
Samfund. Det gaar paa det ethiske Omraade som
paa det aandelige Omraade i det Hele; vi bruge
Udtryk og Forestillinger, som ere hentede fra den
ydre Verden, for at gjøre os den indre Verdens
Væsen og Sammenhæng klar. Ogsaa Forpligtelse
er egenlig et juridisk Udtryk. Baade Lov og Forpligtelse
forudsætte en Myndighed, der ikke blot
befaler, men hævder Befalingens Gjennemførelse.
Uden Magt ingen Lov og ingen Forpligtelse. Men
Magten alene vilde kun fremkalde Frygt. Den er
kun den sidste Udvej (ultima ratio) til Lovens Hævdelse.
Myndigheden, hvorpaa Loven hviler, appellerer
til andre Elementer i Menneskets Væsen end
Frygten. Det Afskrækkende spiller sin største
Rolle paa de laveste Udviklingstrin, hvor Loven
væsenlig er et Hegn mod Fare og Fordærv. Efter
den gamle Mythe var den første Lov, der udgik
til Mennesket, Forbudet mod at spise af et vist
% --- p. 57 ---
\setcounter{footnote}{0}%
\opage{57}Træs Frugt, da det vilde medføre Døden. Barnet
lærer første Gang Loven at kjende, naar det skal
opdrages til Renlighed, og herved spiller det Afværgende
og Afskrækkende den største Rolle. En
mere positiv Betydning faar Loven, naar Myndigheden
ikke blot staar som truende Magt, men ved
sin hele Overlegenhed kommer til at staa som ledende
Forbillede for Mennesket. Naar Mennesket
hos den Myndighed, han er underlagt, ser Evner og
Kræfter, som langt overgaa dem, han føler hos sig
selv, føler han Sympathi med den, og Efterlignelsesdriften
vækkes. Allerede hos Dyrene bliver den
hurtigste og kraftigste Flokkens Fører. Saa snart
Intelligensen er nogenlunde udviklet, vækkes Beundringen
overfor alt Stort, Nyt og Paafaldende; og
med den forbinder sig atter Ærefrygt og Agtelse,
jo mere Mennesket under sine forøgede Erfaringer
føler sin Lidenhed og Afmagt, jo mere altsaa hans
Selvfølelse nedstemmes ved Betragtningen af Forbilledets
Overlegenhed. De Velgjerninger, han modtager
af den højere Magt, vække Taknemmelighed,
og tilsidst knytter han sit hele Ve og Vel til den
med fri Tillid og Hengivenhed. Overfor den blotte
Magt vil Mennesket altid kunne tænke paa Modstand,
indtil han aldeles overvældes og lammes af
Frygt; men naar Følelsen af Afhængighed lutres
ved at forbindes med Sympathiens, Efterlignelsesdriftens,
Beundringens, Ærefrygtens og Taknemmelighedens
Følelser og derved udvikler sig til Anerkjendelse
og Hengivelse, have vi Forpligtelsens Væsen
% --- p. 58 ---
\setcounter{footnote}{0}%
\opage{58}i ethisk Forstand. Tyngdepunktet begynder at flytte
sig fra det Ydre til det Indre.

Dette Forhold til en forpligtende Myndighed
kan spores gjennem hele Menneskehedens Udviklingsgang
under forskjellige Former. En absolut Naturtilstand
uden Avtoriteter findes ikke. Allerede
Hobbes bemærker, som vi tidligere have anført, at
Ingen fødes i Naturtilstanden. Aldeles uden Avtoritet
vilde kun saadanne Væsener være, der voxede
op af Jorden som i de Gamles Myther eller efter
det homeriske Udtryk vare Afkom af en Sten eller
et Træ; og kun de vilde kunne være fuldstændige
Individualister. Familieavtoriteten er altsaa den mest
primitive. Paa de laveste Trin, paa hvilke man
finder Mennesker leve (ifølge forskjellige nyere Rejsende
saa vel som efter Herodots Beretning om flere
vilde Nationer), existerer ganske vist Familien endnu
ikke som afsluttet Gruppe; og Moderen tager sig kun
af Barnet saa længe dette er aldeles hjælpeløst; men
allerede ved den Beskyttelse og de Velgjerninger,
som her blive Barnet til Del, og uden hvilke det
ikke kunde leve, lægges Grunden til et Pietetsforhold,
som kan blive Spire til videre Udvikling. Naar det
omstrejfende Liv afløses af mere faste og regelmæssige
Forhold, kommer Familielivet til at spille en større
Rolle. Faderens Avtoritet afløser nu Moderens.
Paa dette Trin er Ægteskabet grundet paa Magt,
ikke paa Kjærlighed; Hustru, Børn og Slaver maa
i lige høj Grad bøje sig for Familiehovedets Villie.\footnote[1]{\emph{Lubbock}: Origin of civilisation.}
% --- p. 59 ---
\setcounter{footnote}{0}%
\opage{59}Den Beskyttelse, de Svage nyde, og det stadige Samliv,
til Dels ogsaa Samarbejde, fremkalder ogsaa
andre Følelser end Frygten, men denne Følelse ligger
dog stedse til Grund. Familieoverhovedets ubetingede
Myndighed indenfor sin Kreds bevarede sig
endog, efter at Familien var bleven Del af et mere
omfattende Samfund. I det gamle Rom var Familiefaderens
og Ægtemandens Magt slet ingen Indskrænkninger
underlagt, han var ikke ansvarlig for
Nogen paa Jorden. Hvor indgribende en Indflydelse
Forholdet til denne Avtoritet maa have øvet paa
de opvoxende Slægters Tænke- og Handlemaade,
vil let ses; og hvor forandrede disse Forhold end ere i
de moderne Samfund, saa have dog stadig de fleste
Mennesker inden for Familien de første Gjenstande
for Frygt og Kjærlighed, Beundring og Taknemmelighed,
Agtelse og Ærefrygt. Det er under disse
levende Kræfters Indflydelse, at Forestillingskredsen
dannes og de første Regler for Handlingen, de første
moralske Love opstaa. Den Enkelte faar her sin
praktiske Verdensanskuelse grundlagt, og denne
Verdensanskuelses Gyldighed føler han som Et med
selve den Myndighed, hvorfra den stammer. Krænkelse
eller Modsigelse af den føles saa stærkt, fordi
den rammer selve den inderligste og dybeste Følelse
hos Individet.

Blodets Baand er det første, som forener Menneskene
og derved skaber en Avtoritet. I lange
Tider har det været det eneste eller i hvert Tilfælde
det mægtigste; men Familien viser ud over sig
selv, idet den ene Families Hoved staar selvstændig
% --- p. 60 ---
\setcounter{footnote}{0}%
\opage{60}overfor den andens. Ved Sammensmeltning af de enkelte
Familier og Slægter opstaar et borgerligt Samfund.
En absolut Grænse mellem de to Arter af
Samfund kan, historisk set, være vanskelig at drage,
da Familiens Hoved jo tidligere havde meget af den
Myndighed, som senere Statsoverhovedet fik. I den
romerske Stat kan man spore den oprindelige Sammenslutning
af de enkelte Slægter. «Den romerske
Stat», siger Mommsen,\footnote[1]{Römische Geschichte. I. S. 59.} «beror paa det romerske
Hus, saa vel med Hensyn til Elementerne som med
Hensyn til Formen. Folket opstod, idet de gamle
Slægter paa en eller anden Maade sluttede sig sammen;
det romerske Landomraade opstod ved disse Slægters
forenede Distrikter; romersk Borger var den, der
hørte til en af disse Slægter.» Blodets Baand ligger
her tydelig til Grund. Et saadant primitivt
Samfund kan kun voxe ved Adoption; endog Venskabsforholdet
faar hos mange vilde og barbariske
Folk først sin Betydning, naar Vennerne have blandet
Blod og saaledes i det Mindste symbolsk bragt
et naturligt Baand tilveje.\footnote[2]{Jfr. \emph{Maine}: Early history of institutions. S. 228.} Enhver, som ikke var
knyttet til Stammen ved dette Baand, var Fjende
eller Slave. Da Hovederne for de enkelte Slægter
staa som Ligemænd, vil der ikke være noget stadigt
Overhoved for hele Flokken. Lige som Sammenholdet
i det Hele endnu er saa svagt, at Flokken
adspredes og hver sørger for sig, naar der er Nød
paa Færde, saaledes har man ogsaa kun en Fører,
% --- p. 61 ---
\setcounter{footnote}{0}%
\opage{61}saa længe man behøver ham. Dog ville de Gamle
og Kloge, de Stærke og Modige staa i større Anseelse
end de Andre. Og jo mere Kampen med
andre Stammer gjør Sammenhold nødvendig, desto
mere vil Avtoriteten ogsaa befæstes. Krigen synes
derfor, som navnlig Bagehot har fremhævet, at have
virket som forbindende Middel indenfor Stammen.
Førerne i Kampen betragtedes som Alles Beskyttere.
Under meget simple Kulturforhold kan der ikke opstaa
nogen stor ydre Afstand mellem Bydende og % PRINTED AS IS: „og|og“ (doubled over the line end)
og Lydende; derfor faa vi først en egenlig politisk
Avtoritet, naar Ulighed i Ejendommens Fordeling
fjerner de Herskende fra dem, der adlyde dem, og
naar tillige Samfundet bliver tilstrækkelig stort til,
at Herskeren unddrages de Enkeltes stadige Blikke.
Herskeren bliver nu Gjenstand for Ærefrygt; Misundelsen
er udelukket ved selve den store Afstand.
Hans Villie er ubetinget Lov.\footnote[1]{Smlgn. \emph{Adam Ferguson}: An essay on the history of the civil society. Part II. Sect. 3, og \emph{Herodots} Fortælling om Dejokes (Herod. I. 96).} Skjønt Vane og
Instinkt, Vedtægt og Skik danne en Slags Grundlov,
som selv Despoten ikke kan eller vil bryde, er
det dog et Tegn paa Fremskridt, naar en personlig,
bevidst Villie gjør sig gjældende; kun derved bliver
Lov i Betydning af almindelig, upartisk og ubøjelig
Regel mulig.\footnote[2]{\emph{Maine}: Early history. S. 392.} Og dette er ikke blot et Fremskridt
i ren kulturhistorisk, men ogsaa i ethisk Henseende,
for saa vidt den blinde Vane afløses af bevidst
Underkastelse.

% --- p. 62 ---
\setcounter{footnote}{0}%
\opage{62}Magten, den beskyttende og den bydende Magt,
er det første Grundlag for den politiske Avtoritet,
lige som for Avtoriteten indenfor Familien. Støttet
til dette Grundlag kunne de højere og friere Former
for Samfundslivet udvikle sig. Man tænker sig
ofte de Vildes Liv som frit og ubundet; men det
føres tvertimod under et stærkt Avtoritetstryk. «De
styres,» siger Lubbock, «af Regler og Vedtægter,
som danne et af de grusomste Tyrannier, der maaske
nogensinde har existeret paa Jordens Overflade.
De Svages Villie, deres Ejendom, deres Liv ere de
Stærkeres Herredømme underkastede. Hele Systemets
Tendens er at give Alt til de Stærke og Gamle
til Skade for de Unge og Svage, men især for Kvinderne.»\footnote[1]{Origin of civilisation, S. 434.} Naar dette stærke Tryk tages bort, indtræder
ofte det fuldstændigste Anarki, den vildeste
Tøjlesløshed, som hos nogle Negerstammer, hvor Alt
er i Opløsning ved Høvdingens Død. Det er en
streng Skole, Menneskeslægten paa dette Trin maa
gjennemgaa, men den har været nødvendig. Og
selv den mest despotiske Statsavtoritet maa dog have
vakt andre Følelser end Frygten; den optræder jo
ikke blot bydende, men ogsaa beskyttende, yder et
Værn, i Ly af hvilket Livet kan gaa sin Gang.
De Privilegier, som Herskeren tiltog sig, vare i
Virkeligheden kjøbte med Tjenester. Et stort Exempel
herpaa have vi i Middelalderens Institutioner,
Grundlaget for «l'ancien régime». Kongen, Adelen
og Gejstligheden have i raa og vilde Tider ved
% --- p. 63 ---
\setcounter{footnote}{0}%
\opage{63}deres fysiske og aandelige Myndighed gjort menneskeligt
Liv og menneskelig Udvikling mulig.\footnote[1]{Kfr. H. \emph{Taine}: L'ancien régime. Livre I. Chap. 1.: Origine des priviléges.} I Tilliden
og Taknemmeligheden, som den beskyttende
Magt vækker, ligger Spiren til Avtoritetens ethiske
Karakter, som oprindelig trænges tilbage, saa længe
Frygten og den instinktmæssige Sædvane herske.

Vi have fremdraget nogle Grundtræk af Avtoritetens
ældste Historie, saaledes som den er os tilgængelig
især gjennem de senere Aars sammenliglignende % PRINTED AS IS: „sammenlig-|lignende“ (a syllable set twice over the line end)
Studium af de primitive Samfundsformer.
Her ere Forholdene simplest og lettest at overskue;
paa senere Udviklingstrin spille saa mange forskjellige
Elementer en Rolle, at det er forbundet med
Vanskelighed at betragte ethvert af dem for sig. Lige
som Studiet af Fostertilstanden yder Fysiologen værdifulde
Bidrag til Forstaaelse af den fuldendte Organisme
og dens Livsvilkaar, saaledes er Studiet af
hine simple og primitive Livsformer ogsaa af Vigtighed
for Psykologen, Ethikeren og Sociologen.\footnote[2]{Vi have i vor egen Literatur et Værk om dette Emne fra Begyndelsen af dette Aarhundrede af C. \emph{Bastholm} (Historiske Efterretninger til Kundskab om Mennesket i dets vilde og raa Tilstand. 4 Bind. Kjøbenhavn 1803---4). Under Titel af „Descriptive Sociology“ har \emph{Herbert Spencer} med flere Medarbejdere begyndt at udgive tabellariske Oversigter over de forskjellige Folkeslags kulturhistoriske Forhold.}
De ere som Experimenter, Naturen og Historien
have gjort for os paa et Omraade, hvor vi ikke
% --- p. 64 ---
\setcounter{footnote}{0}%
\opage{64}kunne experimentere; og dog have vi i det Foregaaende
gjort Forholdene endnu simplere, end de i
Virkeligheden ere; thi saa snart Mennesket har
hævet sig til den virkelige Begyndelse af et Samfundsliv,
føler han sig ikke blot i Forhold til de Myndigheder,
der raade i Familie, Stamme og Stat,
men til en endnu højere Myndighed, som ligger ud
over disse Samfunds Grænser, og hvori tilsidst hine
Myndigheder selv søge deres Støtte: den religiøse
Avtoritet. Først ved denne sidste Støtte kan Avtoritetsprincipet
blive absolut. Som En af dette
Princips mest konsekvente Forsvarere, Bonald, har
sagt: «Religionen, der er det almindelige Baand i
ethvert Samfund, strammer fortrinsvis det politiske
Samfunds Knude; selve Ordet Religion (religare)
antyder tilstrækkelig, at den er de menneskelige
Samfunds, Familiers og Staters naturlige og nødvendige
Baand. --- Religionen bringer Orden ind i
Samfundet, fordi den lærer Menneskene, hvorfra
Magten og Pligterne stamme. --- Al Magt er skabt
i Guds Billede og stammer fra Gud.»\footnote[1]{G. \emph{Brandes}: Reaktionen i Frankrig. S. 110; 118 f.}

Hvis man som det Væsenlige i Religionen betragter
Følelsen af Afhængighed overfor en højere
Tingenes Orden, maa man i Virkeligheden give dem
Ret, der paastaa, at der gives Folkeslag uden Religion.
En saadan Følelse forudsætter altid en vis
aandelig Udvikling, som ikke findes hos Mennesker,
der mangle Evne til at danne almindelige Forestillinger,
eller som endnu ikke have tænkt over,
% --- p. 65 ---
\setcounter{footnote}{0}%
\opage{65}hvorvidt den Sol, der staar op om Morgenen, er den
samme, der gaar ned om Aftenen.\footnote[1]{\emph{Lubbock}: Origin of civilisation, p. 10; 475. (Mange vilde Folks Sprog have Navne for hver enkelt Farve, hvert enkelt Slags Træ, men intet til at betegne Farve eller Træ i Almindelighed).} Hvis man derimod
til Religion kun forlanger Tro paa aandelige
Væsener,\footnote[2]{\emph{Tylor}: Primitive culture. I. p. 382.} saa kan man ikke nægte, at den findes
hos den raaeste Fetischdyrker, der ikke engang
nærer Frygt for sin Fetisch, men kun betragter den
som nødvendigt Middel og Redskab. En Spore til
et moralsk Forhold viser sig allerede indenfor Fetischismen,
idet Negeren tilhyller Fetischen (ligesom
den russiske Bonde sin Helgen), naar han gjør Noget,
han skammer sig ved. Selv hvor større og mægtigere
Væsener tilbedes, er Frygt længe den herskende
Følelse: man tilbeder kun de onde Væsener;
de gode Væsener, mener man, behøve ingen Opmærksomhed,
da det er deres Natur at være gode.
Denne Dualisme imellem gode og onde Væsener er
derfor heller ikke af ethisk Natur, men udspringer
af den Nytte eller Skade, Mennesket har erfaret.
Solen er saaledes i kolde Lande et godt Væsen, i
varme Lande et ondt Væsen. Det er alt andet end
venlige Følelser, Menneskene paa dette Trin nære
overfor de Guder, de dyrke.

Først i den egentlige Afgudsdyrkelse, hvor Anthropomorphismen
er trængt igjennem, bliver Religionen
virkelig Underkastelse. De ethiske Forestillinger
overføres fra det politiske Samfund, hvor deres
% --- p. 66 ---
\setcounter{footnote}{0}%
\opage{66}første Udvikling har fundet Sted, til Gudeverdenen.
Menneskeslægten synes at have naaet Monarkiets
Standpunkt i Politiken, førend den naaede Afgudsdyrkelsens
i Religionen. Guddommene blive nu
efterhaanden fra blotte Naturmagter til ethiske Magter,
en Overgang, som særlig kan forfølges indenfor
den græske Mythologi. Og lige som Gudetroen ændres
ogsaa Troen paa Udødeligheden; Menneskene
tro nu ikke mere blot paa en i sig ligegyldig Fortsættelse
af Existensen, men paa en Gjengjældelse
efter Døden. Maalestokken for denne Gjengjældelse
er given med, hvad Folket anpriser som Dyder eller
Laster.\footnote[1]{Kfr. \emph{Tylors} og \emph{Lubbocks} ovennævnte Værker og \emph{Prellers} „Griechische Mythologie“.}

Det betegner et højere Trin, naar ogsaa det
Gode erkjendes og tilbedes. Med Taknemmeligheden
forbinder sig Beundringen over det Store og Ophøjede.
Fidjiboernes Ord for Gud (kalon) betyder
i det Hele, hvad der er stort og vidunderligt, er et
superlativisk Udtryk. Og paa det højeste Trin af
Folkereligion, vi kjende, hæver den religiøse Forestilling
sin Gjenstand helt over Naturen, idet Alt
i Verden ses som frembragt af Guddommens Villie.
Følelsen af Afhængighed bliver først her fuldstændig,
idet Mennesket overfor Almagten staar som
absolut afmægtig. Hvad den almægtige Guddomsvillie
byder, er den højeste Lov. Et fastere og
dybere Grundlag synes ikke at kunne findes for den
ethiske Lov. Her har Avtoriteten naaet sin højeste
% --- p. 67 ---
\setcounter{footnote}{0}%
\opage{67}Form, fremtræder som absolut, ubetinget. Avtoriteten
i Familie og Stat er kun afledet i Forhold
til den. Og uagtet den guddommelige Myndighed
for den Troende har iført sig Visdommens, Kjærlighedens
og Retfærdighedens Attributer, saa han i
den ikke blot tilbeder den uendelige Magt men
ogsaa Hellighedens Ideal, er det dog Villiens Almagt
og den Frygt og Bæven, Forestillingen om den
maa vække, som er det sidste Grundlag. De ældre
Theologer udtale uforbeholdent dette. Det er Attraaen
efter den evige Salighed, som fører Mennesket
til Gud; «Menneskene,» siger Augustinus, «bør ingen
Gud dyrke, uden at han kan gjøre dem salige;
hvis Saligheden var en Gudinde, var hun den eneste,
der burde dyrkes.» Lydigheden mod Guds Villie
er derfor, som Augustinus ligeledes viser, den kristelige
Grunddyd.\footnote[1]{De civitate dei. lib. V. præf. --- XIV. c. 12.} Den konsekvente Gjennemførelse
af dette Standpunkt kan derfor ikke opgive Læren
om Helvedstraffene, thi uden dem vilde Myndigheden
svækkes. De ere Slutstenen paa den absolute Avtoritets
Bygning. ---

Disse Træk af Avtoritetens Historie skulde kun
nogenlunde anskueliggjøre et Princip, som har grebet
saare mægtig ind i Menneskeslægtens ethiske Udvikling.
Der var en Tid, hvor man i sin revolutionære
Iver kun havde Øje for de mindre heldige
Steder ved Avtoritetsprincipet og ikke kunde nævne
det uden for at deklamere derimod. Nu ere vel de
Fleste enige om dets store historiske Mission, som
% --- p. 68 ---
\setcounter{footnote}{0}%
\opage{68}endnu ikke er endt, og man er maaske ofte af
kontrarevolutionær Iver tilbøjelig til at anse dets
kulturhistoriske Betydning som et Bevis for dets
principielle Sandhed. Mange moderne Apologeters
Optræden minder, som \emph{Vacherot} træffende har sagt,
om Scipio Afrikanus, der besvarede Anklagen for
Underslæb med de Ord: «Lader os gaa til Kapitolium
og takke Guderne, fordi jeg slog Hannibal ved
Zama!» Mængden fulgte ham, men Regnskabet blev
vel derfor ikke rigtigere. Vi kunne ikke undgaa
at opgjøre Regnskabet. Vi førtes til Avtoriteten
som til den historiske Magt, der oprindelig har vakt
Forestillingen om en højere Lov i Menneskets Bryst,
og hvortil denne Lov hidtil har støttet sig. Lige
saa nødvendigt som det var for os at fremdrage Avtoritetens
Begreb, lige saa nødvendigt er det nu at
undersøge, om det er nødvendig knyttet til Lovens
Begreb, og hvilket af disse to Begreber der ethisk
set er det væsentlige. Da det gjælder et Principspørgsmaal,
maa vi tage Avtoriteten for os i dens
fuldendte, dens ubetingede Form, og saaledes have
vi den, som paavist, i den religiøse Avtoritet. Det
Principspørgsmaal, vi opkaste, er det gamle Stridspunkt
mellem theologisk og filosofisk, kristelig og
human Ethik.

Loven forudsætter en Myndighed, for hvilken
Mennesket bøjer sig. Men hermed er det endnu ikke
givet, at Myndigheden selv er Et med Loven.
Bøjer jeg mig for Loven, fordi den er Udtryk
for Avtoritetens Villie, saa ere mine Bevæggrunde
ikke udsprungne af selve Forestillingen om Loven
% --- p. 69 ---
\setcounter{footnote}{0}%
\opage{69}og dens Indhold, men maa søges i det Forhold,
hvori jeg, afset fra Loven, staar til den Bydende,
dette Forhold være nu et Frygtens, et Pietetens
eller et Ærefrygtens Forhold. Altsaa er det kun
ad en Omvej, at jeg naar til Anerkjendelse af Loven.
Hvis det var en Anden, der udtalte den, vilde
jeg maaske ikke lyde den; men jeg vilde ogsaa lyde
et modsat Lovbud, naar det udgik fra den samme
Avtoritet. Lovens Indhold bliver altsaa en Tilfældighed,
og Mennesket er hildet i den Modsigelse, at
hans mest energiske Handlen, det, hvorpaa han
sætter Alt ind, staar i et rent vilkaarligt Forhold
til hans eget Væsen. Der skyder sig stedse en
tredie Magt ind imellem ham og det, han arbejder
for, og det en ubetinget Magt, over for hvilken alt
Andet bliver Intet. Den ubetingede Avtoritet er
en «nidkjær Gud», som Intet taaler ved Siden
af sig. Og selv om Avtoriteten ikke er absolut, ---
naar den ikke kan gjøre sig til Et med Lovens Indhold,
træder den forstyrrende ind imellem Mennesket
og hans Opgave, hindrer det direkte, inderlige
Forhold.

I den (oprindelige, konsekvente) kristelige Ethik
træder dette klart frem. Det hedder i Almindelighed,
at Kjærlighedsbuddet er den kristelige Ethiks
Hovedlov. I dette Bud, saa stort og ophøjet det
ogsaa er, ligger der i sig selv Intet, der peger ud
over det Humane. At elske Andre, afset fra alle
Forskjelligheder, og elske dem ubetinget, er et Bud,
som allerede Buddha havde udtalt. De stoiske Filosofer
havde ligeledes tidligere opstillet Kjærligheden
% --- p. 70 ---
\setcounter{footnote}{0}%
\opage{70}til hele Menneskeslægten (caritas generis humani)
som ethisk Grundlov, skjønt de ikke formaaede at
skaffe den en saa dyb og inderlig Anerkjendelse
selv hos de Enfoldige, som Evangeliet, der netop
herved godtgjorde sin historiske Mission. Her,
hvor vi kun have med Principerne at gjøre, maa
vi se bort fra den praktiske Betydning, de kunne
have under visse historiske Forhold. Det viser sig
da ved nærmere Betragtning, at det ikke er for
Menneskenes egen Skyld, vi skulle elske dem, men
for Guds Skyld. Kjærligheden er ikke den kristelige
Grunddyd, men Lydigheden, Underkastelsen
under den guddommelige Villie. Uden dette vilde
Kristendommen ikke være nogen positiv Religion.
Positiv Religion vil sige Religion grundet paa Avtoritet,
og Lydighed er Grundforholdet overfor Avtoriteten.
I Kristendommen finder dette nærmere sit
Udtryk derved, at Kjærligheden underordnes Troen.
Vel lærer Apostelen Paulus, at Kjærligheden er
større end Troen, men Kristendommen anerkjender
kun den Kjærlighed, som er grundlagt paa Troen,
Troen paa den overnaturlige Avtoritet. Derfor blive
Hedningenes Dyder i alvorlige Kristnes Øjne kun
glimrende Laster. I sit Kjærlighedsbud har Kristendommen,
ligesom før den Buddha og Stoikerne, peget
ud over de ydre Forskjelligheder, der splittede den
gamle Verden: der skal ikke længere være Forskjel
mellem Græker og Barbar, mellem Træl og Fri;
men den har selv til Gjengjæld etableret en Modsætning,
der er vel saa stærk, som nogen af hine:
den imellem Troende og Vantro, Salige og For%
% --- p. 71 ---
\setcounter{footnote}{0}%
\opage{71}dømte. Troen og Kjærligheden arbejde derfor hinanden
imod,\footnote[1]{Se \emph{Ludwig Feuerbach}: Wesen des Christenthums (Kfr. min Bog „Filosofien i Tyskland efter Hegel“ p. 65---67).} og Kjærlighedsbuddet kommer først
til sin fulde Ret, naar man gaar ud over Troens
Skranker. Men dette kan Kristendommen som
\emph{positiv} Religion ikke gjøre. Derfor har Kjærlighedsbuddet
først faaet sin Fuldendelse ved det moderne
Toleranceprincip, der behandler Troesforskjellene
paa samme Maade, som Kristendommen behandlede
de ydre og nationale Forskjelle. Spinoza er
den Første, som har draget Kjærlighedsbuddets Konsekvenser
og praktisk gjennemført den Sætning, at
Kjærligheden er større end Troen. Det humane
Element i Kristendommen er herved blevet befriet
for de Skranker, der hæmmede det.

Det er ikke blot Kjærligheden, som kun faar
indirekte Betydning under Forudsætning af Avtoritetsprincipet;
det Samme gjælder om ethvert humant
og kulturhistorisk Forhold. Ægteskabet kan
efter den strenge kristelige Betragtning kun blive
det mindre af to Onder. Paulus raader Trællen fra
at søge at opnaa Frihed, og Luther siger paa sit
djærve Sprog: «Livegenskabet strider ikke imod det
kristelige Væsen, og hvo som siger det, han lyver.»
Overfor den ene uendelige Gjenstand, som skal opfylde
alle Troendes Sind, maa Alt skydes til Side,
som kun har menneskelig Betydning.

Man vil maaske indvende, at vi her stille
% --- p. 72 ---
\setcounter{footnote}{0}%
\opage{72}Sagen paa Spidsen og opfatte Avtoriteten som den
rent vilkaarlige. «Ganske vist,» vil man sige, «antage
vi kun det Gode for godt, fordi det er Guds
Villie; men da det Gode er Et med Guds Væsen,
vil han kun det, som er godt, og det beror altsaa
ikke paa Vilkaarlighed og Tilfældighed, hvad der
bliver Lovens Indhold.» Men dette er kun at skyde
hele Spørgsmaalet længere tilbage, med mindre man
tillægger sig selv et saa intimt Bekjendtskab med
Guddommens Væsen, at man kjender de Grunde,
der bestemme dens Villie. Og kjender man disse
Grunde, kan man indse deres Rigtighed, hvorfor
saa gaa den Omvej at appellere til den overnaturlige
Villie, i Stedet for direkte at anvende dem? Saasnart
Avtoriteten ikke kuer os ganske ved et ubetinget
Magtbud, kan den ikke undgaa at blive
krævet til Regnskab, men et Regnskab forudsætter,
at der er Noget, der staar over Avtoriteten, og hvorefter
denne kan bedømmes.

Skjøndt \emph{Martensen} i sin «kristelige Ethik« % PRINTED AS IS: the closing mark is «, the opening sort (image)
energisk hævder den overnaturlige, ubetingede Avtoritets
Nødvendighed som Grundlag for det Ethiske,
har han dog klart indset, at man maa godtgjøre
Avtoritetens Berettigelse. Han løser Spørgsmaalet
ved at sige: «Avtoritetens Ret grunder i det Ethiske,
hvilket Magten er underordnet» (p. 449). Denne
Sætning kunne vi ganske underskrive. Men hvad
forstaar Martensen saa her ved «det Ethiske»? Skal
det være af samme Natur, som det Ethiske, hvis
Forklaring vi søgte, saa maa vi jo i Følge Martensens
Lære atter søge en højere Avtoritet, hvis Villie
% --- p. 73 ---
\setcounter{footnote}{0}%
\opage{73}kan give det Gyldighed --- en Magt altsaa, der kan
forpligte Gud! Og skal den Sætning, at det Ethiske
forudsætter en forpligtende Magt, der staar over den,
som forpligtes, kun gjælde for det menneskelig Ethiske,
ikke for det Ethiske i Guddommens Væsen, saa
bliver dette sidste «Ethiske» for os et tomt Ord,
et Tegn paa, at man har anvendt Begrebet, hvor det
ikke længere har nogen Gyldighed, end sige at det
kan tjene til nogensomhelst Forklaring. Ud over dette
Dilemma kan den theologiske Ethik ikke komme.

Hvor Avtoriteten er absolut, d. v. s. hvor den
hæves over al Begrundelse, maa den kundgjøre sig
ved et Under. For den theologiske Opfattelse er
derfor Natursammenhængen brudt paa det Punkt,
hvor det Ethiske begynder. «Det ethiske «Bør»,
siger Martensen, «bliver i al Evighed uforklarligt af
Naturen. Det kommer ikke nedenfra, men ovenfra»
(p. 441). Og senere hævder han, at det majestætiske:
«Du skal!» ikke engang finder sin Forklaring i vort
ideale Væsen (p. 456). Men for Martensen er det
aandelige Liv i det Hele «et Under for Naturen,
et Transcendent for hele Naturbegrebet» (p. 16).
Hvad vi her maa indvende, er, at Theologien
stadig forvexler det (endnu eller i sig selv) Uforklarlige
med Underet. Underet er et Magtsprog, et
Uvidenhedsasyl (asylum ignorantiæ), hvortil Theologien
tager sin Tilflugt for at udfylde de Huller,
den mener at have opdaget i den videnskabelige
Erkjendelse. Den humane Ethik, der maa anvende
samme videnskabelige Fremgangsmaade, som ellers
gjælder overalt i Tankens og Naturens Verden, bil%
% --- p. 74 ---
\setcounter{footnote}{0}%
\opage{74}der sig ikke ind at kunne løse alle Gaader; den søger
at komme saa langt den kan ad de naturlige Veje;
Resten lader den staa hen. Men den kan ikke gaa
ud fra saa snevert et Naturbegreb som Theologien
gjør. Skjønt Bevidstheden ikke kan udledes af
det ubevidste Naturliv, kan det dog vises, at den,
hvad Betingelserne for dens Opstaaen og Udvikling
angaar, staar i nøjeste Sammenhæng med det og
ikke kan forstaas uden denne Sammenhæng. De
psykologiske og ethiske Love ere derfor Naturlove,
lige som Aandslivet i det Hele er Led i selve det
uendelige Naturliv og fremkommer paa et vist bestemt
Udviklingstrin af dette. For den populære
Opfattelse staar det ganske vist saaledes, som om
visse aandelige Tilstande, Tanker og Følelser havde
en saa ophøjet Karakter, at de umulig kunde bringes
i Sammenhæng med de almindelige psykologiske
Love. Saaledes stode Himmellegemernes Bevægelser
for de Gamle som guddommelige og ideale, saa at
det var en Formastelse at forklare dem paa naturlig
Maade; og dog have de i den nyere Tid fundet
deres mekaniske Forklaring. Saa langt end Psykologien
og Ethiken staa tilbage for Astronomien, saa
hvile de dog paa samme Princip; ogsaa de stræbe
efter en «mécanique céleste», en naturlig Forklaring
af «de himmelske Ting» paa Aandslivets Omraade.
Troen, Tilbedelsen og Kjærligheden ville finde deres
psykologiske og historiske Forklaring, lige som de
mere elementære Bevidsthedsfænomener allerede
have fundet deres.

Det er den historiske Videnskabs Opgave at
% --- p. 75 ---
\setcounter{footnote}{0}%
\opage{75}skildre os Avtoriteternes Oprindelse og Udvikling.
Disse ere at betragte som Naturvæsener, der udfolde
sig ifølge bestemte Love. Den historiske Udvikling
og dens Love staa over alle Avtoriteter og omfatter
dem alle. Naar Betingelserne for en Avtoritets
Gyldighed ere forbi, giver den, som Historien viser,
Plads for en anden. Historien er Avtoriteternes
Kampplads. Men netop deri, at de kæmpe, ligger,
at de dog ere endelige Væsener. De ere som Olympens
Guder, der steg ned paa Valpladsen og kunde
saares af Mennesker.

Heri finde vi Muligheden af at anerkjende Avtoriteternes
store Betydning for det Ethiske og dog
at hævde, at Avtoritetsprincipet i sig selv ikke er
den fulde Grund til det Ethiske. De ere de opdragende
Magter i Slægtens Historie. Men den,
som opdrager en Anden, stræber jo netop at frigjøre
denne, at føre ham til at staa paa egne Ben
og se med egne Øjne. Han stræber at gjøre sig
selv overflødig, at forsvinde selv, for at Disciplen
kan komme Ansigt til Ansigt med Sagen selv. Denne
Resignation fordres af enhver Opdrager, men den
glemme Avtoriteterne paa ethvert Omraade kun altfor
ofte. Netop herved godtgjøre de deres Endelighed.
Thi det, hvorved Avtoriteten er forskjellig fra
den Lov, som skal hævdes, det Formaal, som skal
opnaas, er den Egoisme, den Egenvillie, som den
gjør gjældende paa Bekostning af den Sag, som
skulde fremmes. Hermed stemmer det godt, at Avtoriteten
altid i sidste Instans kommer til at appellere
til egoistiske Motiver, navnlig til Frygten.
% --- p. 76 ---
\setcounter{footnote}{0}%
\opage{76}Ethisk Gyldighed har den kun som Bærer af et
Indhold, der kan godtgjøre sin Berettigelse ganske
afset fra Avtoriteten, og kun derved bliver den Gjenstand
for Beundring, Ærefrygt og Kjærlighed; virker
den derimod udover den Berettigelse, dette Indhold
giver, saa synker den tilbage til dens mest primitive
Former. Saa længe Avtoriteten staar som
et Tredie mellem den Handlende paa den ene Side,
Handlingens Lov og Formaal paa den anden Side,
saa længe handles der kun indirekte ethisk, fordi
andre Motiver raade end selve Anerkjendelsen af
Loven og Ærefrygten for den. Ethiken nedsætter
altsaa alle Avtoriteter til at være relative i Stedet
for absolute, Midler i Stedet for Formaal. Avtoriteten
er til for Menneskets Skyld, Mennesket ikke
for Avtoritetens. Naar Personligheden er bleven udviklet
saaledes, at den anerkjendte Lov er hans
eneste eller dog dybeste Motiv, da er der i ham
selv en Myndighed, overfor hvilken alle ydre Magter
ere Intet. Tyngdepunktet er forlagt fra det
Ydre til det Indre, ---
al Opdragelses store Maal!

Det er den nyere Tids Opgave at forvandle Avtoriteterne
fra absolute til relative. I Stedet for
den Tynge, hvorunder Livet stønnede, saa længe
en overnaturlig Avtoritet dels umiddelbart, dels gjennem
sine Repræsentanter og Tjenere, omfattede Alt,
Stort og Smaat, med sine Bud og sine Forbud,
arbejder den nyere Tid paa at frembringe en friere
og højere Udvikling ved at gjøre det Tunge til
Underlag i Stedet for Tag og Spir. Herfor kæmpede
Reformationen og Revolutionen, hver paa sin
% --- p. 77 ---
\setcounter{footnote}{0}%
\opage{77}Maade. Avtoriteten er en Livsbetingelse for Samfundslivet
og derved ogsaa for det ethiske Liv, men
Livet naar først sin højere Udvikling, naar Avtoriteten
underordner sig. Den Overbevisning vil trænge
igjennem, at de ethiske Ideer have deres Gyldighed
og deres Værd i sig selv, ved deres inderlige
Sammenhæng med Menneskelivets Væsen og Vilkaar,
skjønt de historisk udvikle sig i Ly af Avtoriteten.
Ved at arbejde paa at gjøre denne fra
absolut til relativ, arbejder man derfor i det Ethiskes
Tjeneste. For de Flestes Bevidsthed have de ethiske
Ideer rigtig nok indgaaet en saa nøje Forbindelse
med Forestillingen om den absolute Avtoritet, at
et Angreb paa denne ogsaa forekommer dem at
være et Angreb paa det Ethiske; heraf den stadige
Beskyldning mod den frie Forsken, at den fører til
uethiske Konsekvenser. Lige som Hedningene i
Alexandria ventede, at Himmel og Jord vilde falde
tilbage i det gamle Kaos, da de Kristne omstyrtede
Serapis' Billedstøtte,\footnote[1]{\emph{Gibbon}: The decline and fall of the Roman Empire. Chap. 28.} saaledes mene Mange nu,
at Menneskelivet i ethisk Henseende vil opløse sig
i et Kaos, naar de gamle Dogmer styrtes fra den
Plads, de hidtil have indtaget som styrende Magter.
For den, der har et friere Blik og en frejdigere Tro,
stiller Sagen sig anderledes. Han vil uden fantastiske
Illusioner og uden at skræmmes af de forbigaaende
Udskejelser, som aldrig udeblive, hvor store
og nye Ideer ere i Gjennembrud, overlade det til
% --- p. 78 ---
\setcounter{footnote}{0}%
\opage{78}den historiske Udviklings stille, men sikre Gang
at bringe Harmoni mellem Livets praktiske Former
og de nye Ideer. Vi have i en tidligere Sammenhæng
set, af hvor væsenlig Betydning det er for
det enkelte Menneske at være i Harmoni med sin
egen højeste Erkjendelse. Det Samme maa gjælde
for hele Samfundet, hele Slægten: er en Ide altsaa
virkelig begrundet i den menneskelige Natur, saa vil
den ogsaa arbejde sig frem til praktisk Anerkjendelse,
hvor lidet vi end formaa at overskue alle de
Mellemtrin, de Veje og Omveje, som her først maa
være at tilbagelægge. Her faar da \emph{Pascal} Ret, naar
han siger: «Lad os arbejde paa at tænke rigtig;
det er Moralens Princip!» Det er i det Mindste
det Bidrag, den theoretiske Ethik burde yde til
det ethiske Livs Udvikling.

Kun de Avtoriteter ere altsaa ethisk berettigede,
som indse deres egen relative Betydning. Den overnaturlige
Avtoritet maa stedse besvare enhver Opfordring
til at underordne sig paa samme Maade som
dens mest konsekvente Repræsentant, Romerkirken,
besvarer Opfordringen til at anerkjende Historiens
Fremskridt: «Umuligt!» (Non possumus). De humane
Avtoriteters Væsen krænkes derimod ikke ved en
saadan Opfordring. Myndigheden er her væsenlig
et Udtryk for Traditionen, som den ene Slægt forplanter
til den anden, giver denne med paa Vejen
som uundværlig Støtte, indtil den kan staa paa sine
egne Ben. Kun derved bliver der Sammenhæng i
Menneskeslægten, og kun derved bliver Fremskridt
og Udvikling mulig, at hver ny Slægt ikke begynder
% --- p. 79 ---
\setcounter{footnote}{0}%
\opage{79}forfra, men staar paa de foregaaendes Skuldre. Brydningen
mellem den ældre og yngre Slægt beror dels
paa, at de overleverede Former ikke altid vige, efter
at deres egentlige Rolle er udspillet, dels ogsaa paa
de nye Formers ubændige og utaalmodige Frembrud.
Men selve de haarde Kampe, som herved kunne opstaa,
have deres store Betydning: Friheden maa
vindes gjennem Kamp; Anerkjendelsen af den indre,
højere Lov kan ikke indpodes udefra. Den Modstand,
som selv de relative Avtoriteter saaledes ere
tilbøjelige til at gjøre mod deres egen Underordning,
tjener netop til inderligere, friere og kraftigere Udvikling
af de nye Former. Dog maa det stedse
staa som en væsenlig Opgave at gjøre denne Overgang
saa kort og saa let som muligt, og dette vil
blive muligt, jo mere den Overbevisning udbreder sig,
at «Avtoritetens Ret grunder i det Ethiske»: jo mere
Myndighederne virkelig handle som opdragende Magter,
og jo inderligere og dybere derfor ogsaa deres Betydning
anerkjendes af Alle. Den ethiske Lov skal
ikke blot udgaa fra Avtoriteten, men ogsaa anvendes
paa den. En absolut Avtoritet er hævet over alle
Love, selv over dem, den selv har givet; men de relative
Avtoriteter kan man maale med deres egen Alen.

Som Led af Slægten finder Enhver sig sat til
at løse sin ejendommelige Opgave saaledes som den
fremgaar af alle historiske Betingelser. Han er
Barn af Fortiden og hjælper med at skabe Fremtiden.
Han kan ikke indse sin Afhængighed af det
Forbigangne uden tillige at indse, at det Kommende
til Dels bestemmes ved hans Færd. Saa begrænset
% --- p. 80 ---
\setcounter{footnote}{0}%
\opage{80}hans Synskreds og hans Virkekreds end maatte være,
er han dog aldrig udelukket fra at betragte sig som
Medarbejder i Menneskehedens store, gjennem Tiderne
fremadskridende Samfund. Bevidstheden herom,
under hvilken Form den end lægger sig for Dagen,
er Et med Følelsen af Forpligtelse, med Anerkjendelse
af en højere Lov. Her have vi det Ethiske
i dets simpleste og reneste Form. \emph{Comte} har smukt
og træffende vist, hvorledes hin Bevidsthed opstaar
indenfor Familielivet og stedse viser tilbage til dette.
Gjennem Pietetsforholdet til Forældrene træder Barnet
i Forhold til Slægtens Fortid og faar strax et
dybt Indtryk af, hvor nøje hele dets Tilværelse
hænger sammen med den. I Broderforholdet og
senere i Forholdet til Medarbejdere og Samtidige i
det Hele føres den Enkelte ind i den hele store
Kreds af solidariske Interesser, hvor ogsaa han skal
yde sin Skjærv. Sammenhængen træder her ikke
saa klart frem, som indenfor Familien, men jo mere
Kulturen skrider frem, des nøjere forbundne vise
dog alle menneskelige Bestræbelser sig at være, saa
at Fremskridt i én Retning i Længden er umulig uden
tilsvarende Fremskridt i andre Retninger. Faderforholdet
endelig danner Slutningen og peger mod
Fremtiden. Hvad der i sin simpleste Form optræder
indenfor Familien, kan udvides og udvikles til at
omfatte hele Slægten. Der er én stor Kreds af Opgaver,
som væsenlig betinges ved Forholdet til Fortiden;
den omfatter Alt, som tjener til at bevare og
sikre det menneskelige Liv og den Organisation, de
faste Former, uden hvilke Livet her lige saa lidt
% --- p. 81 ---
\setcounter{footnote}{0}%
\opage{81}som noget andet Sted kan bestaa. En anden stor
Kreds af Opgaver peger mod Fremtiden; hertil hører
Alt, hvorved Livet kan modtage højere og videre
Udvikling og Berigelse. Orden og Fremskridt forudsætte
hinanden: Ordenen er Fremskridtets Betingelse,
og Fremskridtet er Ordenens Formaal.\footnote[1]{Politique positive I. 95; 105.} I
begge Kredse af Opgaver finder det personlige Liv
sit Indhold og sin Næring. Det harmoniske Samfund
af fri Personligheder, som ogsaa i denne Sammenhæng
viser sig at være den højeste ethiske Ide,
er ikke et blot Ideal, men er relativt og delvis virkeliggjort
i enhver ethisk udviklet Personlighed. Hvis
Maalet stedse laa i Fremtiden, vilde den ethiske
Handlen faa Karakteren af en tantalisk Jagen efter
Noget, som stedse unddrager sig. Men det Rige,
paa hvis Komme der skal arbejdes, er allerede kommet,
hvor der arbejdes ret, hvad enten saa Arbejdet
i det enkelte Tilfælde nærmest koncentrerer sig om
Udviklingen af Individets eget personlige Liv i de
forskjellige Retninger, eller det har en mere udadgaaende
Karakter. Sammenhængen med Menneskeslægtens
store og fælles Formaal helliger og adler
selv den ringeste Gjerning, naar den udføres med
Troskab.

Det er ikke blot abstrakte Love og Betragtninger,
der træde Tanken imøde ved Anskuelsen af
Menneskelivets Udvikling. Historiens store Skikkelser,
i hvilke de menneskelige Opgaver og Formaal
fremtræde lige som personificerede, staa som
% --- p. 82 ---
\setcounter{footnote}{0}%
\opage{82}Forbilleder, som Typer for, hvad det menneskelige
Liv kan være. Nogle af disse Skikkelser have
endog i den Form, i hvilken de ere os overleverede,
et saadant Fuldendelsens Præg, at de synes hævede
over de Kampe og Vanskeligheder, gjennem hvilke
det er de Flestes Lod at arbejde sig fremad. Mythen
og Legenden have i mange Tilfælde været inspirerede
af den samme Tanke, som \emph{Schiller} udtaler i
et af sine Epigrammer, at ædle Naturer gjælde ved
hvad de \emph{ere}, lavere Naturer ved hvad de \emph{gjøre}.
Det er en Tanke, som har sin store Berettigelse,
da, som nylig bemærket, Vejen og Maalet ethisk
set ikke kunne falde ganske udenfor hinanden. Men
«den stridende Kirke» behøver Forbilleder, som
\emph{virkelig} selv have kæmpet, for hvem Kampen
har været mere end en guddommelig Leg. Forbilledet
kan derfor lige saa lidt være absolut som
Avtoriteten: det vilde være det Samme som at gjøre
det ubrugeligt for relative Væsener. Hvad vi skulle
lære af det, er jo netop, hvad det menneskelige Liv
kan gjøres til ved Anspændelse af Menneskets egne
Kræfter.

Life of great men all remind us,\\
We can make our life sublime. (Longfellow.)

Men saa mægtig en Indflydelse, Tanken om den
fremskridende Menneskehed end kan have, navnlig
naar denne ses repræsenteret i de fremragende Personligheder,
saa fører Betragtningen af det Ethiske
os dog videre. Udviklingen er ikke først begyndt
med Menneskets Optræden. De bedst begrundede
% --- p. 83 ---
\setcounter{footnote}{0}%
\opage{83}Hypotheser\footnote[1]{Smlgn. mit Foredrag „Filosofien og Darwinismen“ („Nær og Fjern“. 1874. Nr. 93---94).} føre os til at se selve Menneskeslægtens
Oprindelse som Resultatet af en lang organisk og
fysisk Udviklingsproces. Gjennem en Række af
Stadier, under de forskjelligste Former, hæver Naturlivet
sig fra de dunkle Taagemasser til den højeste
aandelige Klarhed. Vi kunne overalt finde de
samme Grundlove, den fremskridende Dannelse af
mindre Helheder indenfor det større Hele.\footnote[2]{Kfr. \emph{Herbert Spencer}: First Principles.} Enhver
individuel Form i den organiske som i den uorganiske,
i den menneskelige som i den dyriske Verden betegner
et ejendommeligt Stadium i den store Livsudvikling.
Vi forklare ikke det Højere af det
Lavere, fordi vi se hint fremgaa af dette; tvertimod
forstaa vi først ret Verdenslivet, naar vi se, hvad
det formaar at yde paa sine Højdepunkter, i sine
mest individuelle Former. Mennesket er, netop
ifølge Udviklingshypothesen, ingen Fremmed i Verden,
men har sine Rødder dybt i dens inderste
Væsen; og hvad der for Mennesket staar som nødvendigt,
som sandt og som godt, er lige saa vel en
Realitet som de ydre, materielle Masser. Ved at
arbejde for sine Idealers Virkeliggjørelse, ved at udvikle
og fuldkommengjøre det menneskelige Liv,
spiller han den Rolle, som skal spilles paa det Sted,
hvor han er sat. Han staar som Endemaalet for en
lang Udvikling, og hans naturlige Opgave er at udfylde
denne Plads saaledes, at han ikke blot hævder
% --- p. 84 ---
\setcounter{footnote}{0}%
\opage{84}den, men ogsaa benytter den til at føre Udviklingen
videre. Hans Kamp for Tilværelsen bliver en Kamp
for Menneskets Værdighed. Det er det aandelige
Liv, Tankens, Følelsens og Villiens Liv, hvorved
Mennesket sondrer sig fra de lavere Livstrin, og
som han maa frede om og udvikle, hvis han ikke
vil synke tilbage fra den Plads, han har vundet.
Alle de individuelle Udgangspunkter, vi tidligere
have fremdraget, faa nu deres rette Betydning ved
at føjes ind i den store Sammenhæng, i hvilken
Menneskelivet staar. \emph{Selvopholdelsesdriften},
der ved ensidig at koncentrere sig om det enkelte
Individ bliver Egoisme, faar sin Berettigelse, naar
det ses, at den Enkeltes Kamp for at bestaa er Betingelse
for hele Slægtens Bestaaen og Udvikling.
De Opgaver, der sætte de individuelle Interesser i
Bevægelse, ere tillige de, ved hvis Behandling
Slægten øver og udvikler sine Kræfter. \emph{Sympathien}
bliver det forenende Baand for Slægtens samtidig
og successivt levende Individer; den hæver
sig over det blotte Instinkt, naar den lutres ved
Tanken om den store, fælles Opgave, hvorpaa Alle
arbejde, til alle Tider og under alle Forhold. \emph{Fornuften}
bliver nu mere end en rent formel Evne; den
faar sit levende og fyldige Indhold ved Anskuelsen
af Menneskehedens Levnetsløb, de Love og Betingelser,
under hvilket det finder Sted, de Opgaver
og Pligter, som heraf under givne Forhold følge.

Skjønt Menneskelivet saaledes har sin Grund i
Verdenslivet, er det dog et Forsvindende i Forhold
til dette. De Tider ere forbi, da man kunde slaa
% --- p. 85 ---
\setcounter{footnote}{0}%
\opage{85}sig til Ro ved den Antagelse, at Mennesket var
Verdenslivets Midtpunkt, og Alt til for hans Skyld.
Jo mere Tanken har gjennemtrængt Naturen og
lært dens Love at kjende, og jo mere Mennesket
ved Hjælp af denne Indsigt formaar at benytte
Naturkræfterne i sin Tjeneste, des mere er tillige
den Erkjendelse vaagnet, at han selv, at hele Menneskelivet
kun er et forsvindende Element i Forhold
til det store Hele. Verdenslivets Uendelighed og
de urokkelige Love, efter hvilke det udfolder sig,
overbevise ham om hans Endelighed og Begrænsning.
Den Afhængighedsfølelse, hvori man med
Rette har søgt Religionens Væsen, vil saaledes naturlig
opstaa ogsaa hos den, der ikke lever i Troen
paa overnaturlige Myndigheder. Men al Frygt og
egoistisk Interesse udelukkes ved den klare Erkjendelse,
at det ikke beror paa Menneskenes Ønsker og
Følelser, men paa en bestemt, i faste Love udtrykt
Tingenes Orden, hvad der skal blive til Virkelighed.
Den, som stræber efter Erkjendelse af sig selv og
sine Tilværelsesvilkaar, vil derfor som «frydfuldt Led
i Verdensordenens Kjæde»\footnote[1]{\emph{Dante}: Paradiset, 3die Sang, 54de Vers.} i frejdig Resignation
overlade til Naturens Gang, hvad og hvor meget
der skal virkeliggjøres for ham af hvad han attraar
og haaber. Naar vi holde fast, at det Højeste, vi
kjende, de ideale menneskelige Formaal og Love,
har udviklet sig ved denne Tingenes Orden, bliver
Afhængighedsfølelsen til Ærefrygt; Verdensordenen
bliver ikke blot fysisk, men ethisk. Hvor forsvin%
% --- p. 86 ---
\setcounter{footnote}{0}%
\opage{86}dende end Menneskelivet monne være i Forhold til
Verdenslivet, de højeste Goder, vi kjende, have
deres Plads indenfor dette, og deres Værd forringes
ikke ved, at vi ikke formaa at lade Sol, Maane og
Stjerner staa stille eller forandre deres Gang for at
tjene dem. Vi skylde Livet og Tilværelsen altfor
meget til, at vi skulde stille saadanne Fordringer.
I vor dybe Afhængighed af Naturen føle vi dog
Taknemlighed og Ærefrygt for den store Sammenhæng,
som ogsaa har ladet Aandslivets Frugter
modnes.

Det Ethiske staar saaledes i nøje Sammenhæng
med det Religiøse. Paa begge Omraader spiller
Følelsen af Ærefrygt en Hovedrolle. Forskjellen
beror paa, at indenfor det Ethiske, hvor Mennesket
staar som handlende, knyttes denne Følelse til hans
egen højeste Erkjendelse af Formaal og Opgaver,
men i det Religiøse gaar han længere tilbage og ser
sin Handlen saa vel som hele sit Væsen, alle sine
Ønsker og Formaal, som Led i hele Menneske- og
Naturlivet og underlagt dets Betingelser. Saa længe
den Følelse ikke er udslukt i Menneskeheden, at
saa vel Individet som Slægten ere Led i den uendelige
Helhed, fremgaa og betinges af den, saa længe
vil det Ethiske og det Religiøse ogsaa vedblive at
gjøre sig gjældende under forskjellige Former. Vi
staa her overfor Noget, der har sin faste Grund
i Menneskets Natur og Vilkaar og ikke staar og
falder med vilkaarlige Magters Bud.

% --- p. 87 ---
\setcounter{footnote}{0}%
\chap{--- Den ethiske Lov og Fremskridtet.}
\opage{87}
Avtoritetstroens Tilhængere prise det som et af
de vigtigste Fortrin ved deres Princip, at man af
Sandhedens Kilde umiddelbart faar at vide, hvad
det Gode og Rette er, saa at Afgjørelsen heraf ikke
er overladt Tiders og Individers skiftende Meninger.
Dette vilde ganske vist være et stort Fortrin, naar
de Væsener, der skulde handle, vare ude af Stand
til selv at stræbe efter en Overbevisning i disse
Spørgsmaal, og hvis Hovedvægten var at lægge
paa en vis Sum og Art af Handlinger og ikke paa
det fri Sindelag, hvoraf Handlingerne udspringe.
Men naar man netop i Menneskets Stræben efter
Erkjendelse af sit Væsen og sin Stilling i Naturen
ser hans højeste Kræfter røre sig, vil man give Sir
\emph{William Hamilton} Ret i, at en levende Vildfarelse
er bedre end en død Sandhed, og selve den levende
Stræben efter ethisk Overbevisning bliver et vigtigt
Element i det Ethiske, uden hvilket de harmoniske
% --- p. 88 ---
\setcounter{footnote}{0}%
\opage{88}Personligheders Rige ikke kan opbygges. Det Menneske
er endnu ikke naaet til sand Personlighed,
der staar og lytter efter Befalinger udefra, før han
kan bestemme sig. Det vil desuden altid vise sig
at være en Illusion, naar man mener, at Livet i
Virkeligheden styres af den overnaturlige Avtoritets
Bud. Denne staar for fjærnt til at kunne trænge igjennem
over for de virkelige Forhold, hvor saa mangfoldige
andre --- ædle og uædle --- Kræfter røre sig. «Ikke
én af tusinde Kristne», siger \emph{Stuart Mill} (Om Friheden,
2det Kap.), «lader de Grundsætninger og
Forskrifter, der indeholdes i det nye Testamente,
være en Vejledning og Prøve for hans individuelle
Opførsel. Den Maalestok, hvorefter han retter sig,
er det, der er Skik og Brug i hans Nation, hans
Samfundsklasse, hans religiøse Bekjendelse. Han
har saaledes paa den ene Side en Samling ethiske
Grundsætninger, som han tror en ufejlbarlig Visdom
har givet ham til Regler, hvorefter han skal
rette sin Opførsel; men paa den anden Side har
han en Samling dagligdags Anskuelser og Anvisninger,
der i det Hele taget ere et Kompromis
imellem den kristne Tro og det virkelige Livs Interesser
og Indskydelser, idet de til en vis Grad
stemme overens med nogle af hine Grundsætninger,
ikke stemme saa meget med andre af dem, og endelig
staa i direkte Strid med visse af dem.
Den
første af disse Rettesnore giver han sin Hyldest;
den anden retter han sig i Virkeligheden efter.»
Men det er ikke blot Tradition og Vane, der saaledes
kan hindre den ubetingede Avtoritets Ind%
% --- p. 89 ---
\setcounter{footnote}{0}%
\opage{89}flydelse; ogsaa den fri, selvstændige Overbevisning
om Humanitetens Krav spiller mere og mere en betydelig
Rolle, hvorom hele den nyere Tids sociale
Udvikling tilstrækkelig vidner. Her er da Syn for
Sagen, at der virkelig paa fri og menneskelig Maade
kan erhverves en fast Indsigt i det Ethiske. At
overse denne Kilde, for at søge efter en anden, vilde
være at gaa over Aaen efter Vand.

Kun da kan man med Rette lovprise Avtoritetsprincipet
for dets Sikkerhed og Tydelighed, naar
det i hvert enkelt Tilfælde er muligt at faa Svar
paa Spørgsmaal. Den katholske Kirke yder sine
Tilhængere denne Fordel. Katholiken kan i alle
Tvivlstilfælde henvende sig til sin Skriftefader, der
staar som Kirkens og altsaa i sidste Instans som
Guds Repræsentant, som et Rør, gjennem hvilket
Vældet fra den absolute Sandheds Kilde ledes til
ham, den Enkelte. Protestantismen gaar glip af
denne Sikkerhed og Tydelighed, naar den henviser
hver Enkelt til hans indre Overbevisning, og den
indtager en vaklende og uholdbar Stilling, naar den
vil raade Bod derpaa ved at tilskrive Kirken en relativ
Ufejlbarhed (\emph{Martensens} «Katholicisme og
Protestantisme»). En relativ Ufejlbarhed er ikke
Ufejlbarhed. For Protestanten bliver i Virkeligheden
Samvittigheden den eneste Rettesnor, og han
kan kun følge Kirken, hvor denne stemmer med
Samvittigheden. Protestantismen gjør et Skridt
fremad i Henseende til ethisk Frigjørelse, men bliver
staaende paa Halvvejen. Karakteristisk træder dette
frem hos de ældste protestantiske Moraltheologer,
% --- p. 90 ---
\setcounter{footnote}{0}%
\opage{90}der lærte, at Samvittigheden «i Henseende til Magt
og Myndighed staar midt imellem Gud og Mennesket:
under Gud, men over Mennesket». Her antydes
endnu Muligheden af en Avtoritet, der kunde staa
over Samvittigheden, og det er denne Mulighed,
der adskiller den protestantiske Ethik fra den humane
Ethik.

Staar det fast, at Mennesket kun kan forpligtes
i Samvittigheden, ved fri Anerkjendelse og Følelse
af det Rette, saa er Samvittigheden selv den
højeste Avtoritet, og ikke underordnet nogensomhelst
anden. Samvittighed vil jo sige den højeste Grad
af Klarhed og Inderlighed, hvortil Menneskets
ethiske Tænkning og Følelse er naaet ved al mulig
Anstrengelse. Da det nu er Mennesket selv, der
skal handle, og Mennesket som han er under visse
givne Forhold, kan han kun søge Maalestokken for
sin Handlen i den højeste Forstaaelse, det er ham
muligt at vinde af Handlingens Væsen og Betydning.
Vi maa derfor give \emph{Johann Gottlieb Fichte} Ret
i hans dristige Sætning: «Samvittigheden farer aldrig
vild og kan ikke fare vild». Den gjælder for
Mennesket i Handlingens Øjeblik. Kommer den
bedre Indsigt før Handlingen, vil den om muligt
hindre denne; og kommer den efter Handlingen,
kommer den som en Dommer, der dog maa erkjende,
at Mennesket gjorde ret i at handle efter sin bedste
Overbevisning. Det er Pligt at have Samvittighed;
Mennesket maa stræbe efter en klar og bestemt
Overbevisning om det Rette. Paa ethvert givet
Trin maa han bestemme sig efter den Indsigt, som
% --- p. 91 ---
\setcounter{footnote}{0}%
\opage{91}staar til hans Raadighed, idet han omhyggelig søger
at raade Bod paa sin begrænsede Erkjendelse ved
at foregribe og forestille sig, hvorledes Handlingen
vilde tage sig ud fra et højere Synspunkt. Men da
selve dette højere Synspunkt altid modificeres ved
den Handlendes virkelige Standpunkt og Vilkaar,
som han kun til Dels formaar at se bort fra, er den
Mulighed aldrig udelukket, at Handlingen senere,
under andre Forhold, kan tage sig anderledes ud
for ham. Alligevel vil han, dersom han virkelig
handlede af Overbevisning, indse, at den senere,
klarere og rigtigere Indsigt selv kun blev mulig ved,
at han paa det tidligere Stadium anvendte sin ufuldkomne
Indsigt saa godt, han formaaede. Begge Stadier
ligge altsaa i en og samme Linie, betegne Trin
af Udviklingen henimod et og samme Maal.

Hvad der gjælder for den Enkelte, gjælder for
Menneskeslægten i dens Helhed. De forskjellige
menneskelige Udviklingstrin betegne lige saa mange
forskjellige Trin i Erkjendelsen af den ethiske Lov.
Man kan ikke paastaa, at der hersker Enhed i de
ethiske Grundsætninger hos Menneskene; det gjendriver
Erfaringen og Historien. Enheden er ikke
Udgangspunktet, men Endemaalet. Men hvad der
fremtræder overalt, saa snart en vis menneskelig
Udvikling er naaet, er selve Forpligtelsesforholdet.
Dette kan være til Stede med større eller mindre
Klarhed og Inderlighed, Dybde og Omfang. Hvad
der paa et Sted inddrages under dette Forhold,
betragtes et andet Sted som ligegyldigt. Saa snart
et højere Trin er naaet, fældes Dommen over det
% --- p. 92 ---
\setcounter{footnote}{0}%
\opage{92}lavere, ved hvilket det dog selv er blevet muligt.
For saa vidt er Historien Dommen i Slægtens saa
vel som i den Enkeltes Liv. En absolut Sandhed
gribe vi intetsteds. Det er vor Lod at arbejde paa
Løsningen af en uendelig Gaade, som fremstiller
sig under stedse nye Former, saa længe Livet varer.

Vi føres herved til det Spørgsmaal, hvorvidt
man har Ret til at antage en Fremadskriden i
ethisk Henseende hos Menneskeslægten. Denne
Antagelse, hvortil saa Meget i vore foregaaende
Undersøgelser maa føre os, har man forkastet, idet
man navnlig holdt sig til den indre, subjektive Side
af det Ethiske. Saaledes erklærer den berømte
Naturforsker \emph{Wallace}, der grundig kjender den malayiske
Race, at de Vilde fyldestgjøre deres simple
Morallov i det Mindste lige saa godt, som de civiliserede
Evropæere, ja at Massen af vor Befolkning
ikke er naaet ud over de Vildes Morallov, men i
visse Henseender endog er sunken ned under den.
Man er gaaet et Skridt videre, idet man har afvist
Tanken om et egenlig ethisk Fremskridt som indeholdende
en Uretfærdighed, da de senere Slægter,
der begyndte deres Bane paa et højere ethisk Udviklingstrin,
saa vilde være heldigere stillede end de
tidligere. Da det i det Ethiske væsenlig kommer
an paa, hvad man bringer ud af de givne Betingelser,
gjør det ikke Noget, paa hvilket Trin man
begynder. Under de simpleste og de mest udviklede
Kulturforhold kan Moraliteten altsaa være den
samme. Efter denne Betragtningsmaade, som i
vor Tid er hævdet af \emph{Bouillier}, er det Ethiske lige
% --- p. 93 ---
\setcounter{footnote}{0}%
\opage{93}til alle Tider: Tæller og Nævner kunne vexle,
men Grundforholdet bliver det samme. (1/2 = 2/4
= 3/6 = 4/8 etc.). % the fractions are small solidus fractions in the print

Man gaar her ud fra, at det ethiske Forhold er
til Stede selv paa det laveste menneskelige Udviklingstrin.
Men det ethiske Forhold forudsætter bestemte
Betingelser, som ikke altid kunne have været
til Stede. Det kan ikke bestemt paavises, førend
Samfundet og Avtoritetsforholdet har dannet sig, og
før Menneskets aandelige Evner have naaet en vis
Styrke og et vist Omfang. Oprindelig skjelner
Mennesket kun mellem Godt og Ondt som mellem
det, der bringer ham Glæde, og det, der bringer Smerte.
Saaledes forklarede en Vild, at naar Nogen tog hans
Kone fra ham, var det Ondt, men naar han tog en
Andens Kone, var det Godt. Og efter at det ethiske
Forhold er traadt i Kraft som Anerkjendelse af en
højere Lov og Handlen af Ærefrygt for den, vil det
antage en meget forskjellig Karakter under de forskjellige
historiske Vilkaar. Det Ethiske er ikke
noget Isoleret i Menneskets Væsen, staar ikke ligegyldig
over for hans Udvikling eller Tilbagegang
i alle andre Henseender, men træder i den inderligste
og livligste Vexelvirkning med Alt, hvad der ellers
rører sig hos Mennesket. Det Omfang, den Renhed
og Frihed, hvormed det Ethiske gjør sig gjældende,
vil altid til en vis Grad bero paa det Trin af Kultur
og aandelig Udvikling, Mennesket har opnaaet.
Med en klarere Forstaaelse af Maalet og Vilkaarene
maa ogsaa selve Forpligtelsen antage en højere
Karakter. Selv hvad det ethiske Grundforhold an%
% --- p. 94 ---
\setcounter{footnote}{0}%
\opage{94}gaar, maa vi derfor antage en Fremadskriden. Men
et herfra helt forskjelligt Spørgsmaal er det, om
de senere Slægter derved faa Noget forud for de
tidligere. De Betingelser, under hvilke Udviklingen
begynder for de enkelte Slægter og Individer, raade
de selv ikke for. Vi indtage naturligvis et højere
Standpunkt end de, paa hvis Skuldre vi staa, men at
vi gjøre det, skylde vi netop dem. Og jo mere
Sammenhængen mellem Slægterne træder frem for
Bevidstheden, jo mere navnlig den foregaaende
Slægts opdragende Myndighed over for den opvoxende
ses i sin hele Betydning, des mere vil Pieteten
ogsaa vækkes over for alle dem, der hver paa sin
Maade have kæmpet og arbejdet for at stille os der,
hvor vi nu staa.

Endnu klarere ses Nødvendigheden af at antage
en ethisk Fremadskriden, naar vi se hen til den
objektive Side af Sagen, til de Livsformer og Institutioner,
hvori det Ethiske skaffer sig Udtryk. Selve
den friere og menneskeligere Maade, hvorpaa Livsforholdene
efterhaanden ordnes og lægges til Rette for
Menneskene, maa dog være af højere ethisk Værd
end de raa og primitive Ordninger. Det er opsamlet
Arbejde, der ganske vist som al Kapital baade kan
bruges og misbruges, men som dog repræsenterer
et Fond, Menneskeslægten raader over.

Fra det ethiske Synspunkt betragte vi Menneskeslægtens
Historie som det harmoniske Samfunds
Udviklingshistorie, den fremskridende Virkeliggjørelse
af Humanitetens Rige. Man kan indenfor denne
Udvikling sondre mellem tre Stadier; paa det ene
% --- p. 95 ---
\setcounter{footnote}{0}%
\opage{95}hersker Magten, paa det andet Retten, paa det tredie
Pligten. Disse Stadier kunne paavises saa vel for
Slægtens Vedkommende i det Hele som i hvert
enkelt Individs ethiske Udvikling og indenfor de
enkelte humane og sociale Omraader.

Ingen Opdragelse finder Sted uden Tvang og
uden Anvendelse af Straf. Selv den mest filantropiske
Pædagog maa i Praxis anvende Tvangen, for
at hindre sin Myndling i Handlinger, der vilde være
absolut og øjeblikkelig skadelige for ham selv eller
Andre, og Myndlingen vil ofte foreløbig ikke faa
andet Indtryk heraf end af en hvilken som helst
Anvendelse af Overmagten. I det primitive Samfundsliv,
hvor det Filantropiske kun svagt ytrer sig,
hersker Magten næsten fuldstændig, ledet af de
Mægtiges Selvopholdelsesdrift. Alt Arbejde læsses
over paa de Svage, paa Kvinder og Børn. Den
overvundne Fjende behandles som ethvert andet
Bytte. Kun indenfor Stammen anerkjender man
andre Individers Ret til at bestaa og tilfredsstille
deres Interesser, og selve denne Anerkjendelse er
bleven aftvungen ved Magt. Retten hviler oprindelig
paa Anerkjendelse af det givne Magtforhold.
Indenfor Familien, hvor Retsforholdet ikke gjør
sig gjældende, vedbliver Magten at raade, som det
navnlig ses i den ældre romerske Ret, der fuldstændig
underkaster Hustru og Børn Familiefaderens
Myndighed. Naar Medfølelsen for de Svage, som
netop udvikles og næres indenfor Familien, fører
til, at man ikke længere dræber den forsvarsløse
Fjende, gjøres denne til Slave og bliver som saa%
% --- p. 96 ---
\setcounter{footnote}{0}%
\opage{96}dan endnu lige saa lidt som Hustru og Børn delagtig
i Retsforholdet. Herren staar over for ham
endnu i Naturtilstanden, Krigen ulmer bag ved
Freden, som det ses af saadanne Love som den i
Rom, der paabød, at alle en Mands Slaver skulde
dræbes, naar han var bleven myrdet, og naar de
til den Tid havde befundet sig saa nær, at de kunde
høre hans Raab.\footnote[1]{Smlgn. \emph{Tacitus}: Annales XIV, 42---44 (kfr. XIII, 32.)} Denne Lov udsprang, som \emph{Montesquieu}
træffende bemærker,\footnote[2]{L'esprit des lois XI, 16.} af Krigens Lov,
kun at Fjenderne her vare i Samfundets Indre.
Hvorledes Slaverne behandledes, afhang dels af
Herrens Sindelag, dels af de ydre Forhold. Om
Cato den Ældre fortælles, at han lod sine Slaver
dømme hverandre indbyrdes, altsaa etablerede et
Slags Retsforhold imellem dem af sin gode Villie.
Men naar Mistanken stadig stod bag ved, naar
Herren følte sig som den Ene mellem de Mange,
kunde fuld Tryghed aldrig komme i Stand. Den
Humanitet, hvormed Slaverne alligevel i det Hele
behandledes hos Grækerne og Romerne, havde
sikkert ikke blot sin Grund i disse Folks menneskelige
og liberale Tænkemaade, men ogsaa i Vanskeligheden
ved at skaffe sig nye Slaver, lige som omvendt
Negerslavernes Pinsler for en stor Del stammede
fra den Lethed, hvormed de kunde erstattes,
og som gjorde, at Herrerne fandt deres Regning
ved at lade dem slide sig ihjel og derefter kjøbe nye.\footnote[3]{\emph{Stuart Mill}: Principles of political economy. II, 5,1.}

% --- p. 97 ---
\setcounter{footnote}{0}%
\opage{97}Vigtigheden af, at Magtforholdet afløses af Retsforholdet,
beror paa, at Retten aabner Vejen for
en Begrundelse. Det første Element i Rettens
Væsen er ganske vist Sanktionen fra Samfundsmagtens
Side, idet denne paatager sig at hævde Enhver
i Besiddelsen og Nydelsen af, hvad han er og
har. Men dermed forbinder sig nødvendigvis en
Regulering, en Udjævnen af Tilfældigheder og Uligheder.
Der søges en almindelig Maalestok og bestemte
Principer til Afgjørelse af de enkelte Tilfælde;
men dette maa igjen føre til, at der gives
Grunde for de enkelte Bestemmelser. De gamle
romerske Retslærde søgte derfor at bevise Slaveriets
Retmæssighed.\footnote[1]{\emph{Hugo Grotius}: De jure belli et pacis III, 7, 5; 11, 16.} De gik ud fra, at det er tilladt
at dræbe sin Modstander i Krigen; naar jeg alligevel
skaaner ham efter at have tilintetgjort hans Modstandskraft,
bliver han min Ejendom. De udledte
Ordet servus (Slave) af servare (bevare), og fremstillede
paa denne Maade Slaveriet som et Fremskridt
i Humanitet, hvad det i Virkeligheden ogsaa var i
Forhold til Kannibalismen og den dyriske Rasen
mod den afvæbnede Fjende. Denne Begrundelse
anerkjendtes i nyere Tid endnu af de første Retsfilosofer
(\emph{Grotius}, \emph{Hobbes}, \emph{Locke}). Det var det
18de Aarhundrede forbeholdt at bibringe denne Institution
det afgjørende Stød. \emph{Montesquieu} omstyrtede
den gamle Begrundelse ved den simple Indvending,
at en afvæbnet Fjende, der ikke længere
kan gjøre nogen Skade, har man fornuftigvis ingen
% --- p. 98 ---
\setcounter{footnote}{0}%
\opage{98}anden Ret over end den at sikre sig mod fornyet
Fare.\footnote[1]{„Det 17de Aarhundrede havde ingen Tvivl med Hensyn til Slaveriet. Bossuet accepterede det uden Betænkning som en af Skriften avtoriseret Kjendsgjerning. Locke bekæmpede det rigtignok, men uden Originalitet og Styrke, og han bevarede det endda i visse Tilfælde (Forbrydere og Krigsfanger). Den eneste Diskussion derom før Montesquieu, som fortjener at erindres, er Bodins i det 16de Aarhundrede. Han var den Eneste i dette Aarhundrede, der havde hævet sin Røst mod Slaveriet. Hvad Middelalderen angaar, da var der Enstemmighed med Hensyn til denne Institution, som Kristendommen syntes at have tilintetgjort. Man maa gaa tilbage til Kirkefædrene og Stoikerne for at finde en saa levende Protest som det 18de Aarhundredes. Men Kirkefædrene, der kun støttede sig til Menneskenes religiøse Lighed, og som i den menneskelige Rets Navn tillod, hvad de forkastede i den mystiske og kristelige Rets Navn, havde ikke overskaaret dette fordærvelige Ondes Rod. Deraf kom det, at trods alle Formildelser af Slaveriet, som var omformet til Livegenskab, var Alt beredt til en Fornyelse af dette Onde, der vel var forringet, men ikke tilintetgjort, da Opdagelsen af Amerika og af Mennesker med sort Hudfarve gav Begjærligheden, Overtroen og Uvidenheden Anledning dertil. Doktorernes og Theologernes Stemme hævede sig ikke, med Undtagelse af Las Cacas's, mod dette Attentat paa Menneskeretten. Man maa derfor sige, at det først er det 18de Aarhundrede der har bibragt Slaveriet et dødbringende Slag; det er Montesquieu, der har havt dette Mod og denne Ære; oplyste af ham, af J. J. Rousseau og Andre efter dem, have Evropas Nationer bestemt sig til at befri sig fra denne Plet.“ \emph{Paul Janet}: Histoire de la science politique. 2e. éd. II. p. 504---505.} % PRINTED AS IS: „Las Cacas's“ (for Casas; image)

Ved at gjennemføre den Tankegang, der saaledes
var sat i Bevægelse, og ved den Begejstring for%
% --- p. 99 ---
\setcounter{footnote}{0}%
\opage{99}Humanitet og den Respekt for Individualiteten, der
derved vaktes, kom man i forrige Aarhundrede til
Ideen om almindelige og medfødte Menneskerettigheder,
som formuleredes forskjellig af de forskjellige
Partier og i de forskjellige Forfatninger, og som
\emph{Kant} i sin Retslære reducerede til en eneste: »Frihed
(Uafhængighed af enhver Andens tvingende Vilkaarlighed),
for saa vidt den kan bestaa med enhver
Andens Frihed efter en almindelig Lov, er den
eneste, oprindelige Ret, som tilkommer ethvert
Menneske i Kraft af hans menneskelige Væsen«.
Det er, som man ser, det samme Princip, \emph{Stuart
Mill} i vore Dage saa energisk og indtrængende har
indskærpet i sin Bog »om Friheden«. Kant saa
Betydningen af denne Grundsætning udtrykt deri,
at enhver Indskrænkning af den personlige Frihed
maa forsvares, medens Friheden i sig selv ikke behøver
nogen Begrundelse. Men netop derfor fører
dette Princip os ud over Rettens egenlige Omraade.
Thi for saa vidt der i det Hele kan forbindes nogen
Mening med Ord som »medfødt Ret« eller »Menneskerettighed«,
er det snarere en ethisk end en juridisk
Betydning. At Enhver skal behandles som fri
Personlighed, er en ethisk Fordring, der ikke i
denne Form kan slaas fast som ydre Lov. Men
dertil kommer, at en rent formel Opfattelse af dette
Princip netop vilde hindre dets Gjennemførelse.
Thi det gaar ud fra, at Enhver virkelig er fri Personlighed,
egnet til at have sit Midtpunkt i sig selv,
--- men dette er meget langt fra at være Tilfældet.
Frihed kan ikke være »medfødt«, den maa erhverves.
% --- p. 100 ---
\setcounter{footnote}{0}%
\opage{100}Ved forhastet at anerkjende Mennesker for fri, vil
man netop hindre dem i at blive det. Man indsaa
ogsaa snart, da Begejstringsrusen over »Menneskerettighederne«
var forbi, at det kun lidt kan nytte
at give Mennesket Rang som frit Væsen, naar man
ikke hjælper ham til virkelig at blive det. Her
ligger vore Dages sociale Problem, der i sine mangfoldige
Former berører alle Omraader, lige fra Religion,
Videnskab og Skole til Ejendom, Industri og
Maskinvæsen. Det er saa omfattende, netop fordi
det ikke angaar nogen enkelt Side af Menneskelivet,
men den fulde Udvikling af det personlige Liv
hos Alle. Ingensinde før har Bevidstheden om,
hvad det gjælder, været saa levende.

Vi føres herved til den Overbevisning, at Retsbegrebet
maa underordnes Pligtbegrebet. Retten
bestaar væsenlig i Samfundets Sanktion af den
Enkeltes Magtstilling. Retten kan skille Individerne
ad og kan gjøre, at de betragte Samfundet
som Middel for deres private Formaal. Under
Kampen mod Despotiet fandt man et fast
Holdepunkt i Individets absolute Ejendomsret.
Ejendom er et Atom! --- udraabte den ældre
\emph{Pitt} i Parlamentet. Men vi kjende ingen absolute
Atomer. Intet Atom er løsnet fra de almindelige
Love for hele den Sammenhæng, indenfor hvilken
det befinder sig. Og lige som Retten til Ejendom
forudsætter Samfundsmagtens Sanktion, saaledes maa
Ejendommen selv fra et ethisk Synspunkt ses som
begrundende et Pligtforhold. \emph{Auguste Comte} har
alvorlig bekæmpet den moderne Sondring mellem
% --- p. 101 ---
\setcounter{footnote}{0}%
\opage{101}det Private og det Offenlige; efter hans Anskuelse
bør Enhver betragte sig som offenlig Funktionær.
Ved Ejendommen har den Enkelte, lige som ved
sin ydre Stilling i det Hele, faaet anvist Omraade
og Midler til Arbejde i de kulturhistoriske Opgavers
Tjeneste. De kulturhistoriske Formaal omfatte
Alt, hvad der bidrager til Menneskehedens
Udvikling og Fremgang i indre og ydre Henseende.
Livet og Virksomheden i Familien, i det borgerlige
Samfund og i Staten, i Kunstens og Videnskabens,
i Industriens og Handelens Verden tjener dette
Formaal. Og ved at arbejde paa disse Omraader
lever den Enkelte sig ind i den Følelse, at han er
Led i en stor Organisme, Borger i et Humanitetens
Rige, der arbejder sig frem i Tidernes Løb. Sit
egenlige Jeg finder han da i den Opgave, som tilfalder
ham indenfor denne Udvikling. Personlighedens
Gyldighed og Betydning betinges ved Værdien
af det Indhold og de Interesser, i hvilke den
lever. Det ene Individ kan kun hævde sig og fordre
Anerkjendelse af det andet ved at godtgjøre
Betydningen af hvad det lever i og for. Den sande
Personlighed udvikles ved Arbejdet for Formaal og
Opgaver, som stilles af den virkelige Verden. Virkeligheden
er mere end en Skueplads, hvor Individerne
skulle fremstille sig for at vise deres Virtuositet.
Naar Munkene i den ægyptiske Ørken tilbragte
deres Tid med at vande tørre Stokke, som
plantedes i Sandet, kunde dette være en Lektion i
asketisk Selvopgivelse og tjene som Symbol paa den
menneskelige Virkens Betydningsløshed over for den
% --- p. 102 ---
\setcounter{footnote}{0}%
\opage{102}overnaturlige Verden. Men den humane Ethiks
Livsnerve er, som vi fra flere Sider have set, den
inderlige Forbindelse mellem den Enkelte og Slægten,
og derfor ogsaa mellem den Enkelte og den
virkelige Verden.\footnote[1]{S. \emph{Kierkegaard} har villet hævde en fuldstændig Modsætning mellem Ethik og Historie, og fra sit Standpunkt med Rette. „Gud kunde uden at gjøre Uret og uden at fornægte sit Væsen, hvilket er Kjærlighed, skabe et Menneske, udrustet med Evner som ingen Anden, og sætte ham paa et afsides Sted og sige til ham: Gjennemlev du nu i Anstrængelse som ingen Anden det Menneskelige, arbejd, saa det Halve kunde være nok til at omskabe en Samtid, men du og jeg vi blive ene derom; al din Stræben skal slet ingen Betydning have for noget andet Menneske, og dog skal du, forstaar du, du skal ville det Ethiske, og du skal, forstaar du, du skal være begejstret, fordi dette er det Højeste“. Uvidenskabelig Efterskrift p. 98. --- Fra det asketiske Synspunkt straaler Lydigheden, som her er Grunddyden, desto klarere, jo mindre Resultatet er. Om man omskaber en Samtid eller vander en vissen Stok, kommer i hvert Tilfælde ud paa Et. Smlgn. ovenfor p. 71.} Dette Forhold maa ikke gjøres
til et Skinforhold, hvad det altid maa blive under
dualistiske Forudsætninger. Miraklet strider mod
Ethiken lige saa vel som mod Fysiken og Logiken.

Der er en Harmoni mellem Personligheden og
de ethiske Opgaver. Personligheden udvikles ved
Arbejdet paa Opgavernes Løsning, og omvendt bestaar
Opgaven direkte eller indirekte i det personlige
Livs Udvikling. Et Samfund bestaar i en Flerhed,
en Mangfoldighed, som er forbunden til Enhed
ved ejendommelige Love. Jo friere og selvstændigere
de enkelte Led røre sig, jo flere forskjellige
% --- p. 103 ---
\setcounter{footnote}{0}%
\opage{103}Muligheder, de virkeliggjøre, samtidig med, at Enheden
bevares og stedse faar en inderligere Karakter
og højere Gyldighed, desto fuldkomnere er Samfundet.
Ideen om Menneskeslægten som Personlighedernes
Rige, som levende Enhed af individuelle
Kræfter er derfor den højeste ethiske Ide. Det er
til den, den ethisk Handlende, bevidst eller ubevidst,
under den ene eller den anden Form ser hen, naar
han vejer og vurderer Formaal og Midler. Den er
ikke Frugten af vilkaarlig Spekulation, men er et
Ideal, der har dannet sig for den menneskelige Bevidsthed
under Historiens Gang\footnote[1]{Kfr. min Afhandling: „Forestillingen om Retfærdighed i dens historiske Hovedformer.“ Nationaløkonomisk Tidsskrift. 2det Bind p. 81---116.}, idet Tanken paa
kunstnerisk Vis hentede Elementerne til Billedet
af det Fuldkomne fra den ufuldkomne Virkelighed,
et Ideal, hvis delvise, fremskridende Virkeliggjørelse
kan paavises.

Men Fremskridtet gaar ikke i lige Linie, og ikke
alle Dele af Menneskeslægten ere delagtige i det.
Dette kunde synes at være en vægtig Indvending
mod hele Ideen om Fremskridtet. Dog maa man
ikke over de reaktionære Perioder glemme den
væsenlige Tendens, den Lov, som aabenbarer sig,
naar man sammenligner en Række af Perioder, især
da selve Reaktionen viser sig at være et vigtigt
Led i Udviklingen, der beriger denne med nye Erfaringer
og fastere Former. Al Bevægelse er rhythmisk,
men det udelukker ikke, at den kan være
fremadskridende. »Hvor mange forskjellige Ele%
% --- p. 104 ---
\setcounter{footnote}{0}%
\opage{104}menter tage Del i Bevægelsen, vil den samme Tilstand
aldrig ganske vende tilbage. En politisk Reaktion
bringer aldrig Tingenes gamle Orden tilbage.
Vor Tids Rationalisme er saare forskjellig fra det
forrige Aarhundredes«.\footnote[1]{\emph{Spencer}: First principles. Part II. chap 10: The rhythm of motion.} Og hvad de forskjellige
Racers ulige Deltagelse i Fremskridtet angaar, forholder
det sig her som indenfor Dyreverdenen, hvor
Typens Udvikling i nogle Grupper standser paa et
Punkt, der for andre Grupper kun staar som et
Gjennemgangspunkt. Det menneskelige Foster viser
paa tidlige Stadier Lighed med lavere Typer. Saaledes
gaar det med Folkenes Civilisation; det ene
Folk standser paa et Trin, der kjendes fra andre
Folks Historie som et overvundet Standpunkt. Saaledes
ere mange orientalske Samfund standsede paa
det Udviklingstrin, hvor det borgerlige og det kirkelige
Regimente endnu falde sammen. Man kan
maaske endog, naar man ser bort fra de højeste
Kulturtrin, sige, at Stagneringen er Reglen og
Fremskridtet Undtagelsen.\footnote[2]{\emph{Maine}: Ancient Law p. 22.} Kun en ringe Del
af Menneskeheden har naaet en højere Civilisation.
Men medens den højere udviklede Dyreart ikke formaar
at give de lavere Arter Del i sin Udviklings
Resultater, er en saadan Meddelelse mulig indenfor
den menneskelige Verden. Den evropæiske Kultur
synes at skulle blive universel, at optage de andre
Kulturformer i sig, skjønt Mødet med disse mere
primitive Kulturformer endnu ikke kan siges gjen%
% --- p. 105 ---
\setcounter{footnote}{0}%
\opage{105}nemgaaende at have havt en meget human Karakter.
Evropæerne have over for andre Verdensdeles Beboere
endnu en ikke ringe Grad af den Selvfølelse og
Nationalstolthed, som de Gamle nærede overfor Barbarerne.
Dog vil ogsaa denne Skranke engang falde,
alle Floder ville forene sig i én stor Strøm og et
eneste Kultursamfund omfatte alle Jordens Beboere.

Paa det ethiske Standpunkt har man ikke Tid
til at dvæle ved Tanken om, hvor meget der er
udrettet (»wie wir's so herrlich weit gebracht«), lige
saa lidt som man har Tid til at gruble over, hvad
der kunde have været anderledes. Trykket af, hvad
der staar tilbage, vil endogsaa i Reglen til enhver
Tid være saa stærkt, at Menneskene nødig ville
indrømme, at de ere lykkeligere end deres Forfædre.
Den gode gamle Tid, hvis Laster og Lidelser man
ikke har nogen levende Følelse af, men hvis Dyder
og Goder man tror paa, lovprises derfor i Regelen
paa den nærværende Tids Bekostning. Menneskene
have for det meste følt sig i Jernalderen og henlagt
Guldalderen til Fortiden. Vor Tid synes særlig at
være Pessimismens Tid. Det kan udledes af flere
Aarsager. De store politiske og sociale Reformer i
det sidste Aarhundrede hilsedes med Forventninger
om en Frihedens Guldalder, som ikke ere gaaede i
Opfyldelse, --- bl. A. af den simple Grund, at de,
der frigjordes, ifølge Sagens Natur endnu manglede
den Opdragelse, som kun selve Friheden kan give.
Det store Opsving i Industri og Handel er ikke
blevet ledsaget af den rette Fordeling af alle de
Goder, som derved opnaaedes. Den fremskredne
% --- p. 106 ---
\setcounter{footnote}{0}%
\opage{106}intellektuelle Kultur har fremkaldt en Reflexion
og en Bevidsthed, som Naturfolkene mangle, og som
ofte kan hindre den fri og fulde Nydelse af det
Værdifulde i Livet, medens den skærper Sansen
for det Disharmoniske og for Skyggesiderne. I
ethisk Henseende er der dog netop herved gjort et
stort Fremskridt, i Retning af Sympathi med Andres
Lidelse. En personlig Erfaring af Smerten er nødvendig
for, at Sympathien med Andres Smerte kan
faa den rette Styrke og Udstrækning. Og den indtrængende
Bevidsthed om hvad der staar tilbage,
forudsætter Forestillingen om det høje Maal, der
skulde naaes. Selve Pessimismen vil saaledes, hvor
den er mere end Frugten af den raffinerede Nydelsessyges
skuffede Forhaabninger, føre til en klarere
Opfattelse af Menneskelivets Vilkaar og til Følelsen
af, hvor vigtigt det er, at der arbejdes paa deres
Fuldkommengjørelse. I det ethiske Princip, som
det i det Foregaaende er udviklet, indeholdes Muligheden
af at forlige Arbejdet og Nydelsen, Pligtmoralen
og Nyttemoralen. Den ethiske Grundopgave
er, som vi have set, at udvikle Harmonien
indenfor den individuelle Personlighed saa vel som
mellem de enkelte Personer i Samfundet; men den
højeste Lykke bestaar netop i en saadan Harmoni,
i de personlige Kræfters Udvikling i Samklang med
hverandre.

--- Det ethiske Forhold vil ikke atter forsvinde,
saa længe det menneskelige Livs Grundvilkaar forblive
de samme. Det skylder bestemte Betingelser
sin Opstaaen og sin Virksomhed i Menneskeheden,
% --- p. 107 ---
\setcounter{footnote}{0}%
\opage{107}men disse Betingelser ere, som vi have søgt at vise,
de samme som de, hvorpaa selve Slægtens Bestaaen
og Udvikling beror. Formerne, hvorunder det
Ethiske fremtræder, kunne vexle, uden at det forandrer
sig i sit inderste Væsen. Kun dersom det
faldt ganske sammen med Forholdet til de ydre
Avtoriteter, vilde dets Rolle være udspillet, naar
Emancipationen sejrer; saa vilde vi allerede nu
være Vidne til dets Opløsningsproces. Den fremskridende
Kultur fører til, at Menneskene se med
Ironi paa al Avtoritet, og naar de med denne forkaste
Alt, hvad der i Tidernes Løb opvoxede i Læ
af den, vil ogsaa det Ethiske forsvinde. Men i den
blaserede Stemning, som saa let opstaar ved Overgangen
fra et stort Princip til et andet, maa man
ikke se det endelige Resultat af Udviklingens Gang.
En nøjere Undersøgelse af det Ethiskes psykologiske
og historiske Grundlag fører til Anerkjendelse ikke
blot af dets store og blivende Betydning, men ogsaa
af dets indre Sammenhæng med Menneskehedens
Grundlove.

% --- p. 108 ---
\setcounter{footnote}{0}%
\chap{Villiens Frihed.}
\opage{108}


Anerkjendelsen af den ethiske Lov adskiller sig
fra theoretisk Indsigt og æsthetisk Anskuelse derved,
at den stiller et Krav til os. Men herved rejser
sig da naturlig det Spørgsmaal, hvor vidt Mennesket
er i Stand til at fyldestgjøre dette Krav, eller
Spørgsmaalet om den menneskelige Villies Frihed.
Ogsaa dette er et Problem, som angaar Ethikens
Grundlag, hvorfor vi endnu maa give en kort Antydning
af, hvorledes det er at opfatte fra det Standpunkt,
som i det Foregaaende er angivet. Med
velberaad Hu anstilles denne Undersøgelse først nu,
thi saa omstridte de Spørgsmaal end kunne være,
hvormed vi i de foregaaende Kapitler have sysselsat
os, saa synes dog Problemet om Villiens Frihed at
have en særegen Braad. Dette har til Dels sin
Grund i, at Villiens Psykologi endnu er saa lidet
opklaret; Villien er maaske den mest indviklede
Side af Menneskets Væsen, hvor saare mange bevidste
og ubevidste Faktorer spille en Rolle, medens
% --- p. 109 ---
\setcounter{footnote}{0}%
\opage{109}Erkjendelsens og Følelseslivets Udvikling er simplere
og lettere at følge. Men det maa vistnok
navnlig udledes af, at man i Almindelighed tror at
vide overordenlig god Besked om dette Spørgsmaal,
hvorfor man lettelig føler sig frastødt ved en theoretisk
Undersøgelses afvigende Resultater. Den formentlige
Viden hindrer ikke blot den rette Indsigts
Fremvæxt ved at udelukke Trangen til den, men ogsaa
ved at begunstige Følelser, som gjøre det umuligt at
betragte Sagen paa en uhildet Maade.

Overalt hvor vi søge at forstaa og forklare en
Ting, maa vi gaa tilbage til den større Sammenhæng,
indenfor hvilken den bliver til. Vi maa se den udvikle
sig og betragte dens Forhold til de andre Ting,
der ere Led i samme Kjæde. Denne Fremgangsmaade
have vi fulgt i de foregaaende Undersøgelser
og ville ogsaa anvende den i den nærværende. Ved
saaledes at se Tingen udvikle sig opdager man lettere
de Elementer, som indeholdes i dens Væsen;
de forskjellige Stadier i Udviklingen betegnes netop
ved at nye Elementer slutte sig til de tidligere
givne, og ved at Forbindelsen mellem de givne Elementer
antager en anden Karakter. Men heri ligger
da, at for at forstaa en Ting maa vi opløse (analysere)
den. Som den staar i sin færdige og afsluttede
Enhed er den ubegribelig. Den maa sønderbrydes,
hver af dens Dele maa betragtes for sig, og
vi maa da se, om vi kunne forstaa Delenes Samvirken,
og saaledes sætte sammen igjen, hvad vi have
skilt ad. Ofte gaar det Forskningen, som det gaar
Børnene med deres Legetøj, hvilket de ogsaa rive
% --- p. 110 ---
\setcounter{footnote}{0}%
\opage{110}i Stykker for at studere dets skjulte Væsen, men
hvilket de ikke formaa at sætte sammen igjen. Men
denne Fare opvejes rigelig, naar Opløsningsprocessen
(Analysen) blot lærer os nogle nye Elementer at
kjende; det første Skridt er da gjort og kan drage
flere efter sig; vi kunne jo ikke vente at naa Maalet
paa én Dag. Den umiddelbare Opfattelse føler sig
imidlertid i Reglen frastødt og uhyggelig stemt ved
denne Opløsning af, hvad der for den staar som en afsluttet,
i sig hvilende Helhed, og paa ethvert Omraade
har den forskende Tanke maattet kæmpe en
haard Kamp for at hævde sin Ret. Det synes, som
om al Individualitet trues paa Liv og Død af Forskningen.
For den livløse Verden vil man nu til
Dags ikke længere nedlægge Indsigelse mod Analysens
Anvendelse; man vil her ikke let tale om
Individer. Kun for Fantasien tager et Bjærg eller
en fjærn Ø sig ud som et individuelt Væsen. Men
i Organismen mene endnu mange at have et absolut
Individ, skjønt \emph{Goethes} Ord: «Kein Lebendiges
ist eins, immer ist es vieles», have faaet fuldstændig
Bekræftelse ved Videnskabens Resultater,
og gjælde ikke blot for Organismen som Helhed,
men ogsaa for dens Dele; selv Cellen, i hvilken man
en Tid troede at have fundet det egentlige Individ,
er en lille Verden, og saaledes i det Uendelige.

Skulle vi da aldeles opgive Individualitetens Begreb?
Opløser Naturen sig da ikke for os i et
Kaos, hvor utallige Elementer flyde om imellem
hverandre og kun for et Øjeblik forbindes for at
danne tilsyneladende faste og afsluttede Skikkelser?
% --- p. 111 ---
\setcounter{footnote}{0}%
\opage{111}Denne Slutning udspringer kun af Fantasiens Utaalmodighed
over at skulle opgive den uopløselige Enhed,
hvormed den har udstyret Tingene. Med langt
større Ret maa vi vende Forholdet om og sige, at
Naturens hele Virkomhed er individualiserende. Den % PRINTED AS IS: „Virkomhed“ (for Virksomhed)
analyserende Forsken søger kun de Elementer, som
i Naturen ere forbundne til ejendommelige Helheder.
Paa alle Naturlivets forskjellige Stadier danner der
sig Midtpunkter eller Knudepunkter, hvorfra særlige
Virksomheder udgaa. Enhver af disse ejendommelige
Virksomheders Natur beror paa Vexelforholdet
mellem det givne Midtpunkt og de omgivende Betingelser.
Allerede indenfor den livløse Verden danner
der sig mindre Kredse, som paa ejendommelig
Maade modificere de udenfra kommende Paavirkninger.
Man tænke f. Ex. paa Stoffernes Tilstandsform
og dens Afhængighed af Temperatur og Tryk; i
dette bestemte Forhold til Vilkaarene ligger noget
Individuelt. Men større Selvstændighed og derfor
ogsaa Mulighed for videre og dybere gaaende Vexelvirkning
med Omverdenen finde vi først hos de
levende Væsener, som derfor betegne et højere Trin
af Individualisation. Det levende Væsen formaar
ved det friere og inderligere Forhold mellem sine
Elementer i langt højere Grad at lempe sig efter
Omgivelserne --- ja paa de højere Trin endog at
lempe Omgivelserne efter sig. Den nyere Videnskab,
som lægger en særegen Vægt paa Forholdet mellem
Organismen og dens Tilværelsesvilkaar, har derved
banet Vejen for en fyldigere Opfattelse af Individualiteten,
som man tidligere ofte betragtede som en
% --- p. 112 ---
\setcounter{footnote}{0}%
\opage{112}isoleret Enhed. Det er i den Maade, hvorpaa der
reageres mod visse Betingelser, at Individualiteten
lægger sig for Dagen. Det Individuelle er ikke
Substans, men Funktion. Vexelvirkningens større
eller mindre Rigdom og Mangfoldighed afgiver Maalestokken
for, hvor højt et Væsen staar i Tingenes
Række. Og gjennem denne Vexelvirkning er det
tillige, at Overgangen, efter den sandsynligste Hypothese,
sker fra de lavere til de højere Former. Kampen
for Tilværelsen bestaar deri, at Individet hævder
sig i sin Ejendommelighed under de givne Forhold,
og er altsaa en Kamp for Individualiteten.

Fra dette Synspunkt staar ogsaa det sjælelige
Liv i nøje Sammenhæng med hele Naturudviklingen.
Det betegner det højeste Trin af Individualisation,
vi kjende. Bevidsthed er Betingelsen for, at Vexelforholdet
til Verden antager en omfattende og dybtgaaende
Karakter. De forskjellige Paavirkninger
udefra maa samles, erindres og kombineres; der maa
sondres mellem forskjellige Arter af Paavirkninger,
for at Reaktionen nøjagtig kan svare til det, som
fremkalder den. Men hin Samlen og denne Sondren
kan kun foregaa i en Bevidsthed. De forskjellige
Stadier af Bevidsthedslivets Udvikling --- «fra Indvoldsormen
i et organisk Væv til en Newton eller
Shakespeare, hvis Tanke omfatter Verden» --- ere
ogsaa Stadier af Individualitetens Udvikling.

Det menneskelige Individ har ved sin hele Organisation
Muligheden af at træde i et rigere Vexelforhold
til Omverdenen, end de andre levende Væsener.
Dog er oprindelig ogsaa det kun en Sidebane,
% --- p. 113 ---
\setcounter{footnote}{0}%
\opage{113}ad hvilken Bevægelsen udefra ledes videre, uden at
der i Individet selv sker Andet end en mekanisk
Forplantning af det meddelte Indtryk fra Sanseorganet
til Centralorganet og fra dette til Bevægelsesorganet.
Dette er Tilfældet ved de saakaldte
Reflexbevægelser. Her følger Bevægelse umiddelbart
paa Indtryk, som naar Pupillen trækkes sammen
under Indtrykket af stærkt Lys, eller naar En uden
selv at vide det jager en Flue bort, som har sat sig
paa hans Kind, medens han er fordybet i sit Arbejde.
Saadanne Bevægelser, der ere ligefremme
Udslag eller Reflexer af Paavirkningen udefra, kunne
derfor ogsaa udføres, naar man har berøvet Individet
det bevidste Livs vigtigste Organer, hvad de mange,
almindelig bekjendte Forsøg med Frøer og Salamandre
godtgjøre. Robin saa endogsaa Armen paa
en halshugget Forbryder begynde en afværgende Bevægelse,
da han skrabede ham paa Brystet med en
Skalpel. Endog hvor Bevægelsen udgaar fra Individet
selv, som i de spontane og hensigtsløse Bevægelser,
hvor Nervesystemet søger Afledning for
sin forøgede Energi (se ovenfor p. 11), har den endnu
en rent mekanisk Karakter. Og dog maa vi i Evnen
til spontan og reflex Bevægelse se den Spire,
hvorfra Villien udvikler sig, uden at det er muligt
at angive det bestemte Punkt, hvor det Uvilkaarlige
hører op og det Vilkaarlige begynder. Om Villie
tale vi først, naar der har dannet sig en fast og
sammenhængende Kreds af Forestillinger og Følelser,
hvoraf Individets Handlen bestemmes, saa at
det modtagne Indtryk eller den spontane Impuls
ikke umiddelbart omsættes i udadgaaende Virksom%
% --- p. 114 ---
\setcounter{footnote}{0}%
\opage{114}hed.
Førend der har dannet sig en saadan fast
Kjærne i Bevidstheden, tilskrive vi heller ikke Individet
noget egenligt Jeg. Men det Punkt, hvor
et saadant kan siges at være til Stede, naaes gjennem
en lang Udvikling, som begynder med meget
faa og meget dunkle Forestillinger og Følelser.
Nogle Fysiologer have endogsaa antaget, at der selv
ved de rene Reflexbevægelser finder en Fornemmelse
Sted, skjønt denne ikke kommer til vor Bevidsthed.
Saa meget er i hvert Tilfælde vist, at der ikke kan
drages nogen skarp Grænse mellem de uvilkaarlige
og vilkaarlige Bevægelser. Selv om Fornemmelser
eller Forestillinger optræde imellem Indtrykket og
den deraf følgende Handling, kan denne dog for saa
vidt kaldes rent mekanisk, som det egenlige Jeg,
d. v. s. den hele Kreds af Tanker og Følelser, som
ellers raader i Individets Bevidsthed, slet ikke kommer
til at virke. Dette er Tilfældet ved alle Handlinger,
som foretages af Vane, hvor vi ofte handle
efter en aldeles tydelig Forestilling, uden at denne
træder i Vexelvirkning med den øvrige Forestillingskreds.
Det er i det Hele Tilfældet i «ubevogtede
Øjeblikke», hvor Opmærksomheden ikke er vakt; vi
gjøre da noget Andet end vi egenlig ville. Men
jo større Betydning Mellemrummet mellem Tilskyndelse
og Handling faar, desto mere Grund er der
til at tale om en Villie. I dette Mellemrum kan
Tilskyndelsen bringes i Forbindelse med Bevidsthedens
øvrige Indhold og virke paa den sjælelige
Tilstand i det Hele. Hvilke Tanker og Følelser
% --- p. 115 ---
\setcounter{footnote}{0}%
\opage{115}den herved vækker i det enkelte Tilfælde, beror
dels paa de almindelige psykologiske Love for Forestillingernes
Forbindelse og Reproduktion, dels paa
Individets medfødte Temperament og hans under
Livets Gang udviklede Karakter. En Virkning kan
aldrig udledes af en enkelt Aarsag, men forudsætter
stedse en hel Række af Betingelser. Handlingen
bliver en Følge af den Maade, hvorpaa Individets
Natur modificeres ved Indtrykket eller Tilskyndelsen.
Jo mere indskrænket Forestillingskredsen er, desto
færre Muligheder fremstiller der sig for Bevidstheden,
og desto mere umiddelbart fører Indtrykket til Handling.
Jo vanskeligere det er at sætte Forestillingskredsen
i Bevægelse, saa at Erindringsbilleder kunne
dukke frem eller Tanken vende sig mod det Fremtidige
eller i det Hele mod Gjenstande, der ligge
fjærnt fra den i Øjeblikket givne Situation, desto mindre
er Villien udviklet. Børn og Vilde ere derfor Øjeblikkets
Børn, lade sig umiddelbart henrive af den
første den bedste Tilskyndelse. Rejsende fremhæve
som den mest karakteristiske Egenskab hos de Vilde deres
Mangel paa Evne til stadigt og udholdende Arbejde,
hvis Belønning først Fremtiden skal bringe. Deres
Følelser vækkes hurtig, men vare kun kort. De
have, som man har sagt, Børns Forstand og Mænds
Lidenskaber. Uden Erindring bliver der ingen Sammenhæng
i Bevidsthedslivet og derfor heller intet
egenligt Jeg. I Stedet for at skjelne mellem vilkaarlige
og uvilkaarlige Bevægelser kunde man derfor
skjelne mellem Bevægelser, som forudsætte Erindring
og Besindelse, og saadanne, der foretages
% --- p. 116 ---
\setcounter{footnote}{0}%
\opage{116}umiddelbart og uden Betænkning. \emph{Wundt} har
med Rette fremhævet, at Dyr, som ere berøvede
Bevidsthedens Centralorganer, mangle Erindring,
hvilket ses, naar man lader nogen Tid gaa hen mellem
Forsøgene; de samme uhensigtsmæssige Bevægelser
foretages da paa ny, førend de hensigtsmæssige
findes. Men Grænsen bliver stadig bevægelig,
idet Handlinger, som først udførtes med
Besindelse, efterhaanden foretages aldeles uden Betænkning.


Villien forudsætter altsaa Overvejelse. Forud
for den villede Handling gaar der en Kamp i Bevidstheden
mellem forskjellige Muligheder, forskjellige
Tilskyndelser til Handling, og Villiens Gjenstand
bliver den Tilskyndelse, som sejrer over de
andre. Jo rigere og stærkere udviklet Bevidsthedslivet
er, desto mere vil ogsaa Handlingen kunne kaldes
Individets egen, efter som den afgjort er bestemt
ved hans Natur. Den faste Kjærne, som har dannet
sig i Bevidstheden, vil gjøre sig gjældende ved enhver
væsenlig Handling, der skal udøves, og stemple
den med sit Præg. Der har med andre Ord dannet
sig en Karakter, en Personlighed. Fornuften har,
som Hamlet siger, blandet sig med Blodets Drift,
«saa Lykken ej paa ham som paa en Fløjte, kan
spille hvilken Melodi hun lyster». Beslutningen er
ikke et umiddelbart Udslag, men Frugten af en
Proces, der er foregaaet i Individets inderste Væsen,
og hvortil alle Elementer i dette have ydet deres
Bidrag.

Enhver Beslutning fremgaar af de forhaanden
% --- p. 117 ---
\setcounter{footnote}{0}%
\opage{117}værende psykologiske Forudsætninger med samme
Nødvendighed, som Virkning overalt følger efter
Aarsag. Men det er ikke nogen ydre Magt, der
overvælder Mennesket; hans egen Natur, hans eget
Jeg gjør sig gjældende i det stærkeste Motiv i
en særegen, ved Forholdene bestemt Form. Det er
Mennesket selv, der handler. Den Handlende kan
have en Følelse af, at intet af de Motiver, der have
gjort sig gjældende for hans Bevidsthed, indeholder
den tilstrækkelige Aarsag til Handlingen. Men det
viser kun, at man maa gaa dybere til Bunds for at
finde Handlingens Aarsag. Bag de bevidste Motiver
ligge de ubevidste, nedarvede eller erhvervede Anlæg,
som øve en afgjørende Indflydelse paa vor
Handlen. Heraf kommer det, at vi kun efterhaanden
ved Erfaring lære vor egen naturlige Karakter
at kjende; ved det, vi gjøre, erfare vi, som
\emph{Schopenhauer} siger, hvad vi ere, og ofte overraskes
vi i høj Grad ved de Opdagelser, vi paa
denne Maade gjøre.

I Menneskets Væsen ligger, som vi tidligere
have set, en oprindelig Drift til Virksomhed; hvorledes
den ytrer sig, bestemmes ved den fra Slægtens
og Naturens Haand modtagne Organisation. Individets
egne Erfaringer og Handlinger efterlade sig
ogsaa Spor og Eftervirkninger, der gjøre sig gjældende
som Anlæg og Dispositioner i visse bestemte
Retninger. Vi begynde altsaa aldrig absolut forfra
i vor Handlen. Som Led i en stor Sammenhæng
og i en uendelig Udviklingsrække have vi vort bestemte
Udgangspunkt og vore bestemte Forudsæt%
% --- p. 118 ---
\setcounter{footnote}{0}%
\opage{118}ninger, som ere Et med vort individuelle Væsen,
vor oprindelige Karakter. Derfor er Villiens Frihed
fuldstændig illusorisk, naar man derved forstaar
en Evne til absolut at begynde en Handling, til at
være Aarsag uden at være Virkning. En Aarsag,
som ikke er Virkning, som altsaa er sin egen Aarsag
(causa sui), maa være et absolut Væsen, hævet
over alle Betingelser, løsrevet fra den lovmæssige
Tingenes Orden. Men et saadant Væsen er os
ubegribeligt. Dette have skarpsindige Forsvarere
af Villiens absolute Frihed ogsaa indset; for dem
er den et Postulat, som ikke kan begrundes. Men
hvor meget man end anerkjender den menneskelige
Erkjendelses Begrænsning, kan man dog ikke
saaledes afkaste Broen midt i Fænomenernes Verden.
Hvis Sammenhængen her sprænges paa et enkelt
Punkt, vil det være forbi med al Lovmæssighed og
al Naturerkjendelse. Enten maa man gjøre Mennesket
til en Gud, eller ogsaa maa man anerkjende, at han
opstaar og udvikler sig efter bestemte Love, som
til en vis Grad maa være eller blive os tilgængelige.
Erfaringen viser os nu tydelig nok, at han er et
endeligt Væsen, selv om Naturens individualiserende
Proces i ham har naaet, hvad vi maa betragte som
dens højeste Punkt. Individualiteten beror her
som overalt paa, at der danner sig en mindre, relativt
afsluttet Kreds indenfor den større, og at de
udefra kommende Paavirkninger modificeres ved de
for denne mindre Kreds gjældende Love.

Men lige saa vist som vi først lære vor Karakter
at kjende af vore Handlinger, lige saa vist er det
% --- p. 119 ---
\setcounter{footnote}{0}%
\opage{119}ogsaa, at vore Handlinger omvendt bidrage til at
bestemme og udvikle vor Karakter. Det er en
fysiologisk Kjendsgjerning, at et Organs Struktur
udvikles ved Øvelse. Virksomhedens Resultat er
ikke blot det udadtil fuldbyrdede Værk; det er ogsaa
den forøgede Færdighed, den mere udviklede
Disposition i en vis Henseende, som Individet derved
har opnaaet. Enhver Følelse, enhver Beslutning,
enhver Handling føjer en Sten til Bygningen
af Menneskets Karakter. Og naar der nu er naaet
en saadan aandelig Udvikling, at Individet kan blive
sig dette bevidst, idet han begynder at fyldestgjøre
den sokratiske Fordring: Kjend dig selv! --- saa vil
han ogsaa kunne søge at lede Forestillingernes, Følelsernes
og Beslutningernes Løb i den Retning,
han anser for den bedste. Mennesket er stillet over
for den indre Natur, lige som han er stillet over for
den ydre. Han kan ikke gjøre Mirakler. Han kan
ikke gribe ind i Tingenes Gang for at lede den
efter sine Øjemed, uden for saa vidt han kan lære
de Betingelser at kjende, under hvilke Tingene opstaa
og virke: kan han bringe disse Betingelser tilveje,
da kan han ogsaa lede Naturfænomenerne efter
sin Villie, f. Ex. aflede Lynstraalen, eller han kan
i det Mindste forudse deres Indtræden og træffe
Forholdsregler mod deres Følger. Paa lignende
Maade gaar det ham med Hensyn til hans eget
Væsen. Han kan heller ikke her skabe eller øve
Mirakler. Men har han tilstrækkeligt Kjendskab
til sig selv, det vil sige, til den Maade, hvorpaa
hans Væsen paavirkes af Omstændighederne og
% --- p. 120 ---
\setcounter{footnote}{0}%
\opage{120}Forholdene, saa kan han søge at forandre disse Forhold
for derved ogsaa at forandre sine egne Forestillinger
og Følelser og lede dem i en vis Retning.
Hvor langt man ad denne Vej kan naa, kan kun
Erfaringen vise. Grænsen vil altid være afstukket
ved de oprindelige Anlæg i Menneskets Natur; og
kun Fantasteriet vil vove sig udover den. Ogsaa
her viser Mesterskabet sig i Begrænsningen. Men
meget, der umiddelbart set synes umuligt, kan dog
naaes, om ikke direkte saa indirekte. Vi kunne
ikke umiddelbart indvirke paa Hjærtets Bevægelse,
men ved at holde Vejret og trykke Brystkassen
sammen eller ved kemisk Paavirkning kunne vi
modificere dets Bevægelse, naar vi ville. --- »Kan man
ikke springe over det atlantiske Ocean som over en
Grav, saa kan man dog sejle over det, og skjønt
man forsikrede Kolumbus, at det ikke gik an, og
at Enhver, der forsøgte det, nødvendigvis maatte
omkomme, gjorde han det dog og omkom ikke.
Paa lignende Maade ere Karakterforvandlinger mulige,
men rigtig nok kun indenfor et vist givet Naturels
Grænser, og ogsaa her kræver Udførelsen
Tid, Mod og Udholdenhed«.\footnote[1]{\emph{Fortlage}: Psykologiske Foredrag. Kbhvn. 1876. p. 187.}

Alt beror paa, at Opmærksomheden og Interessen
er vakt. Paa Erkjendelsens Omraade ledes
Mennesket først umiddelbart af Indtrykkene og de
Ideassociationer, som fremkaldes af dem, men disse
have ofte aldeles Intet at gjøre med Tingenes virke%
% --- p. 121 ---
\setcounter{footnote}{0}%
\opage{121}lige Sammenhæng. Naar han af Erfaringen ledes
til at indse dette, vækkes Opmærksomheden, og han
stræber at danne og forbinde sine Forestillinger
saaledes, at de svare til Virkeligheden. Bevidstheden
er intet Steds aldeles passiv. Aktivitet er
Betingelsen for Receptivitet. Opmærksomheden,
som er af saa stor Betydning for Erkjendelsen, er
allerede en Villie.\footnote[1]{„Apperceptionen [Sindets opmærksomme Opfattelse af Noget] og Impulsen til frivillig Bevægelse ere kun forskjellige Former for Villiens Vækkelse. --- --- De vilkaarlige Bevægelser udgaa højst sandsynlig fra de forreste Dele af Hjærnen, medens de centrale Sansningsorganer fortrinsvis synes at findes i Hjærneoverfladens længere tilbage liggende Dele. Paa den anden Side kan man neppe tvivle om, at de højere aandelige Funktioner navnlig ere bundne til Forhjærnens Udvikling. Denne Sammenhæng bliver først forstaalig, naar vi overveje, at Villiesinnervationens Udgangspunkter tillige beherske Sansningscentrerne, og saaledes ikke blot bestemme Bevægelserne, men Opfattelsen af Sanseindtrykkene og de reproducerede Forestillingers Løb.“ \emph{Wundt}: Physiologische Psychologie. p. 765 f.} Smertelig Erfaring er ofte
Betingelsen for at den kan vækkes. Paa Handlingens
som paa Erkjendelsens Omraade bliver Mennesket
klogt af Skade. Den umiddelbare Naturdrifts
Herredømme maa ofte lede til afgjort Strid med
hvad Mennesket fastholder som sin rette Opgave,
sit sande Formaal, førend en indtrængende Bevidsthed
om Handlingens Betydning kan opstaa. Naar
Tanken er vaagnet og Bevidstheden om en højere
Lov har udviklet sig, er der vaagnet en ny Drift i
Mennesket, der stræber efterhaanden at beherske
% --- p. 122 ---
\setcounter{footnote}{0}%
\opage{122}hans Handlen. Men det er kun Skridt for Skridt,
at Modsætningen mellem de to Tendenser bliver
tydelig; selve Fristelsen vil da ifølge Ideassociationens
Lov vække Forestillingen om den Lov, mod
hvilken den strider, og derved gjøre en Kamp mulig.
Den ved Øjeblikkets Indtryk vakte Lidenskab
vil længe vedblive at fordunkle Bevidstheden om
den væsenlige Lov.

„Noch ist, bei tiefer Neigung für das Wahre,\\
Ihm Irrthum eine Leidenschaft“. (Goethe).

En Betingelse for, at Forestillingen om en væsenlig
Lov kan danne sig, er, som vi tidligere have
set, at Individet har levet sig ind i Tanken om
den større Helhed, til hvilken det hører. Dette
sker gjennem Opdragelsen, og derfor er Opdragelse
ogsaa Forudsætning for Selvopdragelse. Mennesket
maa ledes, for at kunne lede sig selv. Friheden i
Betydning af Evne til Selvopdragelse kan kun blive
til, hvor denne Betingelse er fyldestgjort. Atter
her se vi, hvor nøje den Enkelte er knyttet til det
Hele; hvad han formaar, er afhængigt af hans
Forhold til dette. Det er kun det fuldt udviklede
Individ, som kan være frit i den Betydning, i hvilken
vi her tage dette Ord.

Til Forsvar for Antagelsen af den absolute Frihed,
af Evnen til at være Aarsag uden at være Virkning,
beraaber man sig i Reglen dels paa den almindelige
Bevidsthed, den umiddelbare Følelse, dels
paa Angeren og Tilregneligheden.

Den umiddelbare Bevidsthed og Følelse kan her
% --- p. 123 ---
\setcounter{footnote}{0}%
\opage{123}Intet afgjøre. Den lærer os kun, hvad vi føle og
forestille os, men ikke, hvorledes Følelsen og Forestillingen
er opstaaet eller i hvad Forhold de staa
til den virkelige Sammenhæng. Den begrunder
ingen Theori, hverken i den ene eller i den anden
Retning; det er først den reflekterende Tænkning,
som stiller Problemet og søger en Forklaring. Den
umiddelbare Bevidsthed har ikke Uret i at paastaa,
at vi selv ere Aarsag til vore Handlinger; hvad
den har Uret i, er, at vi skulde være absolut Aarsag,
Aarsag uden at være Virkning. Den antager
dette, fordi den ikke føler sig tvungen eller sat i
Bevægelse udefra, naar den beslutter, paa samme
Maade som den antager, at Jorden staar stille, fordi
Bevægelsen ikke mærkes. Vi mærke ikke selv vor
Villies Bevægelse, lige saa lidt som vi mærke Jordens;
begge Steder slutte vi os til Bevægelsen.
Intet Steds er der Noget, som bevæger, uden selv
at bevæges.

Vi have tidligere vist, hvorledes det højere Livstrin
fælder Dommen over det lavere. Lige som
Verdenshistorien er en Verdensdom, saaledes er det
enkelte Menneskes Historie Dommen over ham selv.
Vi maale vor tidligere Handlen efter den fuldkomnere
Erkjendelse og Følelse af det Rette, som vi
siden have opnaaet. I den rolige og harmoniske
Stemning fælde vi Dommen over hvad vi gjorde,
medens Lidenskaben beherskede os. Saaledes anlægge
vi en ideal Maalestok paa os selv, en Maalestok
fra et andet Standpunkt end det, hvorfra der
handledes. Den Smerte, der opstaar ved Følelsen
% --- p. 124 ---
\setcounter{footnote}{0}%
\opage{124}af Modsætningen mellem vor virkelige Handlen og
den ideale Handlen, kalde vi Anger. Angeren
bliver altsaa forstaalig og ses i sin store ethiske
Betydning, selv om vi ikke antage nogen absolut
Frihed. Den bestaar efter vor Opfattelse i den
gjennemtrængende og pinlige Følelse af \emph{endnu ikke}
at være naaet videre. Hvor stærk og inderlig denne
Følelse er, beror paa, med hvilken Styrke og Klarhed
Forestillingen om Idealet, om den forpligtende
Lov fremtræder. Angeren peger fremad; den er
den ethiske Karakters Fødselsveer, et Vidnesbyrd
om vort Anlæg til Udvikling. Dersom den skulde
forudsætte, at vi ogsaa kunde have handlet anderledes
under de givne Forhold, vilde det i hvert
Tilfælde være en aldeles ørkesløs Betragtning at
dvæle ved, og \emph{Fichte} vilde have Ret i sin energiske
Ytring, at vi ikke have Tid til at angre.
Angerens Værd beror dels paa, at Forestillingen om
det Rette bliver klarere og stærkere ved den dybe
Følelse af det Urette, dels paa det nøjere Kjendskab
til sig selv, Mennesket vinder gjennem den.
Det Skete kan ikke gjøres ugjort, men de nye
Muligheder, som fremstille sig ved Selverkjendelsen,
virke ansporende og fremaddrivende. Fra den
humane Ethiks Standpunkt kan det derfor gjælde:
felix culpa! (lykkelige Brøde!)

Hvad endelig Tilregneligheden angaar, da beror
den paa, at Mennesket er saa udviklet, at der har
dannet sig et Jeg, en fast Kjærne i hans Bevidsthed,
som kan bestemme hans Handlen. Derimod
forudsætter den ingen absolut Frihed. Selv om jeg
% --- p. 125 ---
\setcounter{footnote}{0}%
\opage{125}nok saa meget kan give en udtømmende Forklaring
af min Handling, kan indse alle Motivers Sammenhæng
og Nødvendighed, maa jeg dog forkaste den
og tilregne mig den, naar jeg indser dens Uoverensstemmelse
med hvad jeg anerkjender som den
menneskelige Handlens højeste Lov, og naar jeg
var ved min fulde Fornuft, da jeg foretog den.
Motiverne vare mig jo ikke fremmede; de vare netop
Udtryk for mit Væsen og fremgik med Nødvendighed
af det under de givne Forhold. Hvorledes
kunde ellers Handlingen være min egen? Dersom
jeg i et enkelt Øjeblik kunde handle saaledes, at
Handlingen aldeles Intet havde at gjøre med mit
virkelige Væsen, ikke med Nødvendighed fremgik
af dette, vilde min Personlighed spalte sig i mange
Dele. Vel skjelner man med Rette mellem Menneskets
egenlige Jeg eller sande Villie og den Maade,
hvorpaa han viser sig i sin virkelige Handlen, og
Apostelen Paulus siger jo endog: »Jeg gjør ikke
det Gode, som jeg vil, men det Onde, som jeg
ikke vil; men naar jeg gjør, hvad jeg ikke vil, saa
er det ikke længere mig, der gjør det, men Synden,
som boer i mig.« (Brevet til Romerne VII, 19---20).
Men man maa ikke forvexle en billedlig Modsætning
med en virkelig Adskillelse. Hvad vi kalde
det egenlige eller sande Jeg, kunde ogsaa kaldes
det ideale Jeg, Karakteren som den vilde være,
naar den befriedes fra sine Ufuldkommenheder. Men
det har ingen Existens uden i selve det virkelige
Jeg. Vore Fejl ere »vore Dyders Fejl«, og det er
% --- p. 126 ---
\setcounter{footnote}{0}%
\opage{126}kun for at undskylde os selv, at vi antage vort
»bedre Jeg« for vort egenlige Jeg.

Tilregneligheden peger lige som Angeren fremad,
ikke tilbage. At jeg erkjender mig som tilregnelig
med Hensyn til en vis Handling, vil sige, at jeg
føler Forpligtelsen til at arbejde videre for at
komme ud over den Ufuldkommenhed, Handlingen
har lagt for Dagen. Jeg føler mig hjemfalden til
fornyet Opdragelse. Frihed i Betydning af Evne
til Selvopdragelse er, som vi have bemærket, først
mulig, naar der er gaaet en Opdragelse forud. Det
er en Evne, som forudsætter visse bestemte Betingelser.
Naar Bevidsthedens Klarhed og Livlighed
enten paa Grund af den unge Alder, af Sindssygdom
eller af forbigaaende abnorme Tilstande ikke
er til Stede i en saadan Grad, at Mennesket kan
»se fremad og tilbage«, overveje Handlingens Væsen
og Følger, saa betragter heller ikke den juridiske
Lov ham som tilregnelig, og det er et Tegn paa
den voxende psykologiske Indsigt, at de nyere
Straffelove lægge større Vægt paa saadanne Omstændigheders
Indflydelse end de ældre (sammenlign
f. Ex. Straffeloven af 1866 med Kristian V.s danske
Lov). Et normalt og fuldt udviklet Bevidsthedsliv
er Betingelse for Tilregneligheden. Denne Betingelse
er ikke til Stede hos Børn, hvis Opdragelse
ikke er fuldendt, ikke heller hos Sindssyge, eller
rettere, den findes hos dem kun paa et lavere Trin.
Thi ogsaa Barnet kaldes jo til Regnskab og straffes
for, hvad der ligger indenfor dets ethiske Omraade;
og hvad de Sindssyge angaar, har \emph{Maudsley} træffende
% --- p. 127 ---
\setcounter{footnote}{0}%
\opage{127}bemærket, at man maa antage en vis Ansvarlighed
hos dem, idet man paa Sindssygeanstalterne virker
paa de Afsindige gjennem de sædvanlige Motiver,
skjønt man ikke straffer dem som fuldt ansvarlige,
naar disse Motiver ikke faa Magt over dem.\footnote[1]{\emph{Maudsley}: Le crime et la folie. p. 190. --- Smlgn. \emph{Griesinger}: Pathologie und Therapie der psychischen Krankheiten. 2te Aufl. p. 479: „Paa Sindssygeanstalterne hæmmes den Syges ydre Uro og den lydelige Ytring af hans sygelige Drifter allerede ved de andre Syges Exempel, ved den herskende Fredens og Ordenens Aand; han drages af sig selv ind i hele Husets rolige Bevægelse; mulig Modstand fra hans Side mødes ikke saa meget ved direkte Tvang som af hans egen Følelse af Underkastelse under det Heles imponerende Magt. Han finder her Skaansomhed og Opmærksomhed, men mærker ogsaa, at Gjenstridighed ikke vilde hjælpe; han lærer snart at rette sig efter Lægens Anordninger og ser, hvorledes den Maade, han behandles paa, det Maal af Frihed og Nydelse, som kan blive ham til Del, afhænger af hans Fatning og Adfærd. Saaledes finder han her væsenlig Hjælp til Selvbeherskelse.“} Tilregneligheden
er altsaa relativ, varierer med Betingelser
og Forhold. Fuldt udviklede (»voxne og
sjælssunde«) Personer tvinges af Samfundet til at
adlyde dets Love. Ogsaa dette er en Opdragelse
til et bestemt Formaal og forudsætter kun Tilregnelighed
i den Betydning, at der ikke hos Personerne
maa være Noget, som udelukker Muligheden af at
naa dette Maal. Den, som under ellers normale
Vilkaar i Gjerning viser sig at mangle Agtelse for
Samfundets Lov, maa bringes til at anerkjende den
ved at føle dens Magt. Vort psykologisk-ethiske
Synspunkt fører os saaledes til Samstemmen med
% --- p. 128 ---
\setcounter{footnote}{0}%
\opage{128}den Maade, hvorpaa Professor Goos i sin »Indledning
til den danske Strafferet« (p. 60---64) har begrundet
Statens Straffemyndighed.\footnote[1]{Nærværende Afhandling var i det Væsenlige færdig, da Prof. Goos's Bog udkom, men det var mig til stor Opmuntring og Bestyrkelse at finde en beslægtet Tankegang i dette skarpsindige Værk.} Fra et juridisk
Synspunkt kan man ganske vist ikke betragte Opdragelse
i Betydning af ethisk Forbedring som
Straffens Formaal, da Statsmagten ikke har Ret til
at gribe ind i den Enkeltes inderste personlige Liv.
Heller ikke Lægebehandlingen af de Sindssyge gaar
jo ud paa at forbedre dem, men kun paa at restituere
det gamle Jeg i dets Sundhed; og \emph{Griesinger}
finder heri Forskjellen mellem Opdragelse og Sindssygebehandling.
Hvad Samfundet maa forlange, er
praktisk Anerkjendelse af dets Myndighed, og ved
Straffen fremtvinges en saadan, hvor den har vist
sig at mangle. Ved Straffen bøjes, som Prof. Goos
siger, Forbryderens Karakter til Lydighed mod
Loven. »Det Indgreb i Personens Retsgoder, den
Lidelse, som Straffen indeholder, sigter til at bibringe
Personen ny og kraftigere Motiver til sin
Karakters Optugtelse til Lydighed mod Loven, efter
at det har vist sig, at den ham selv overladte Optugtelse
ikke har hidført, hvad Samfundsmagten i
denne Henseende kræver«. Prof. Goos drager af
dette Resultat den Slutning, at der som Forudsætning
for Straffens Berettigelse ikke kræves Tilregnelighed
i Betydning af Evne til absolut Frihed,
»Villiesfrihed i den Forstand, at Kausalitetslovens
% --- p. 129 ---
\setcounter{footnote}{0}%
\opage{129}Anvendelse paa de menneskelige Villiesakter eller
Opfattelsen af Forholdet mellem Motiver og Beslutning
som Forholdet mellem Aarsag og Virkning er
udelukket«. Den nødvendige Forudsætning er kun,
at Optugtelse ved menneskelig Bestræbelse ikke er
umulig, og det er den ikke, saa længe det endnu
er muligt at vække Menneskets Selvfølelse, hans
Opmærksomhed overfor sig selv, hans Evne til
Energi og Resignation. Kun den kan være hjemfalden
til Opdragelse, hos hvem denne i det Hele
er mulig. Forbrydelsen giver Samfundet Ret til at
opdrage Forbryderen. Dennes Forhold til Avtoriteten
synker paa Grund af hans egen Gjerning tilbage
til Udgangspunktet; da han ikke føler Ærefrygt,
maa han føle Frygt.

Vi have tidligere indskærpet, at det enkelte Individ
aldrig kan udsondres af den større Sammenhæng,
hvori det opstaar og udvikler sig. Alligevel
var det en saadan Udsondring, som forudsattes ved
den Maade, hvorpaa man i ældre Tider betragtede
Straffen. Hvad enten man ved denne vilde true
og afskrække, gjengjælde Øje for Øje og Tand for
Tand eller forsone en vred Guddom, stod Forbryderen
som hele Samfundets Fjende, som den, over
for hvem man ikke længere havde nogen Pligt. For
en stor Del udsprang denne Betragtningsmaade vel
af den lidenskabelige Hævnfølelse, der ser Krænkelsen
som absolut og derfor ingen Begrænsning
tilsteder; men ogsaa Dogmet om Villiens absolute
Frihed har her øvet sin Indflydelse. Det opløser
Slægten i Individer, i absolut uafhængige Udgangs%
% --- p. 130 ---
\setcounter{footnote}{0}%
\opage{130}punkter. Saa snart man derimod indser den dybe
Sammenhæng mellem den Enkelte efter hans Naturs
og hans Karakters hele Omfang og hele Slægten
og Samfundet, maa det ogsaa stille sig som umuligt,
at Noget som helst aldeles skulde kunne løsne
denne Sammenhæng. Som selve Forbrydelsen ofte
viser tilbage til Misforhold og Brøst i Samfundslivet,
saaledes maa det paa den anden Side være
Samfundets Opgave ikke at tilintetgjøre, men fremhjælpe
og udvikle de Spirer til sandt menneskeligt
Liv, som kunne skjules selv i det tilsyneladende
mest forhærdede Sind; det arbejder derved paa sin
egen Sag, sin egen Udvikling.

Villiens Frihed, saaledes som vi anerkjende den,
strider ikke mod Nødvendigheden. De naturlige
Loves Gyldighed omfatter ogsaa den ethiske Handlen.
Vi forstaa ved Villiens Frihed den Tilstand,
hvor Menneskets Fornuft og Følelse og i det Hele
hans aandelige Evner ere udviklede saaledes, at et
enkelt Indtryk, en enkelt Følelse eller Lidenskab
ikke fuldstændig behersker ham eller hindrer Besindelse
og Overvejelse, Anlæggelse af den Maalestok,
der for ham staar som den højeste. Frihed
er her altsaa ikke Modsætning til Nødvendighed,
men til Blindhed og Tilfældighed. Den muliggjør
Selvbeherskelse og Selvopdragelse. Men hvad enten
vi betragte Friheden som Evne eller som Tilstand,
er den stedse relativ. Vor Frihed bestaar i en
Stræben efter at blive fri. Kun faa ere de Punkter,
hvor den Enkelte virkelig udfolder fuld Aktivitet.
Jo videre hans Udvikling i denne Henseende er
% --- p. 131 ---
\setcounter{footnote}{0}%
\opage{131}skreden fremad, desto større vil ogsaa det Bidrag
være, som hans selvbevidste Handlen yder til Karakterens
Dannelse. Og naar Følelsen skærpes for
Betydningen af dette Bidrag, naar Opmærksomheden
med Styrke drages i denne Retning, danner Bevidstheden
sig et idealt Forbillede, hvor Bidraget
fra Personlighedens bevidste Handlen bliver saa
stort, at de andre Elementer i Forhold dertil forsvinde,
--- Idealet af den fuldkomne Frihed, af hvad
den menneskelige Personlighed vilde være, naar alle
hindrende Vilkaar og Indflydelser tænktes borte.
Dette Ideal er kun en særegen Form for den almindelige
ethiske Lov, der, som tidligere vist, bestaar
i Forestillingen om det fuldkomne Menneskeliv
som det, der udvikler sig gjennem Tiderne. I
denne Forestilling er ogsaa Billedet af den frit ɔ:
af Menneskets bevidste Personlighed udsprungne
Handlen et nødvendigt Element; det udtrykker den
Virksomhedsform, som stemmer med Menneskets
Plads i Naturvæsenernes Række.

Men dette ethiske \emph{Ideal} forvexler den almindelige
Opfattelse med den \emph{virkelige} Villie, som altid
er betinget. Den fuldkomne Frihed tænke vi os
ved Abstraktion og Idealisation. Tanken om den
virker, som ethvert Ideal, ansporende og tilskyndende;
men det beror paa en Illusion, naar man
omdanner den til en Evne, der skulde være til
Stede hos Alle og under alle Forhold. En saadan
Illusion er ikke noget enestaaende, men har paa ethvert
Omraade hindret Opfattelsen af den naturlige
Udvikling. Man vægrer sig ved at erkjende det
% --- p. 132 ---
\setcounter{footnote}{0}%
\opage{132}Fuldkomne som Resultatet af en Udvikling og betragter
det som fra først af værende. Saaledes
tænker man sig Sjælen som et Væsen for sig, der
indplantes i Legemet, og det fuldt udviklede Legeme
som allerede til Stede i Spiren en miniature.
Hvad selve Udviklingen først fører til, skal paa en
eller anden Maade have været til fra Begyndelsen
af. Derved er man kommen til at tilskrive alle
Mennesker i hvert Øjeblik som Evne, hvad der i
sig selv kun er Maalet for en lang Udvikling.

De fleste Forsøg paa at hævde Villiens absolute
Frihed have just ikke en saadan videnskabelig Form,
at de indbyde til nærmere Drøftelse. Kun én Form
for den absolute Frihedslære vil det lønne sig at
fremdrage her, nemlig \emph{Kants}, da han er den mest
fremragende af de Forskere, man vil kunne beraabe
sig paa i denne Sag. Her er nu for det Første
den Mærkelighed, at Kant selv med stort Eftertryk
fremhæver den menneskelige Villens og Handlens
Naturnødvendighed, for saa vidt Menneskets Væsen
er os tilgængeligt gjennem Erfaringen. Kant er i
Psykologien afgjort Determinist. »Alle Menneskets
Handlinger i Erfaringens Verden,« siger han i »Kritik
der reinen Vernunft,« »bestemmes ved hans
empiriske Karakter og de andre medvirkende Aarsager
efter Naturens Orden, og naar vi til sidste
Grund kunde udforske alle hans vilkaarlige Handlinger,
vilde der ikke gives en eneste menneskelig
Handling, som vi ikke skulde kunne forudsige og
erkjende som nødvendig Følge af de forudgaaende
Betingelser. For denne empiriske Karakters Ved%
% --- p. 133 ---
\setcounter{footnote}{0}%
\opage{133}kommende gives der altsaa ingen Frihed, og kun
efter denne Karakter kunne vi dog betragte Mennesket,
naar vi blot ville \emph{iagttage}, og, som det sker
i Anthropologien, fysiologisk udforske de bevægende
Aarsager til hans Handlen.« Det er den samme
Grundopfattelse, vi i det Foregaaende have søgt at
gjennemføre. Kant var saa gjennemtrængt af Determinismens
Aand, at han i sin Methodelære hævder,
at selv de vildeste Hypotheser, naar de blot holde
sig indenfor de naturlige Aarsagers Omraade, ere
at foretrække for Appellen til det Overnaturlige,
der kun tjener den »dovne Fornuft« som Paaskud
for at opgive Arbejdet med at finde Aarsagerne ad
Erfaringens Vej. Men ifølge Kants hele Opfattelse
af det Ethiske hører dette aldeles ikke til Fænomenernes,
til Erfaringens Verden. Den ethiske
Lov, som kundgjør sig i vor praktiske Fornuft, kan,
som vi tidligere have set, ifølge Kants Lære slet
ikke modtage nogen Begrundelse. Den lyder til
os fra en højere Verden end Fænomenernes, absolut
uafhængig som den er af al Erfaring. For saa vidt
Menneskets Villie bestemmes af den ethiske Lov,
er han altsaa hævet over Erfaringsverdenen. Foruden
sin empiriske Karakter, som kan blive Gjenstand
for Erfaring og Iagttagelse, har Mennesket
derfor ogsaa en »intelligibel« Karakter, ifølge hvilken
hans Handlen ledes af Fornuften. »Fornuften selv«,
siger Kant, »er intet Fænomen og ikke underkastet
nogen af den sanselige Erfarings Betingelser.« Ved
den faar Mennesket Evne til at handle med absolut
% --- p. 134 ---
\setcounter{footnote}{0}%
\opage{134}Frihed ɔ: uafhængig af alle i Erfaringen givne Motiver
og Tilskyndelser.

Man vil strax se, at Kant aldeles ikke antager
nogen Beslutning eller Villen uden nødvendig Aarsag.
Hvad han forstaar ved Frihed er en Bestemmelse
ved ideale Grunde (Fornuft, Pligt) i Modsætning
til Bestemmelsen ved empiriske Grunde (Drifterne,
Lyst og Smerte). Ogsaa for den »intelligible«
Karakters Vedkommende er der et nødvendigt Forhold
mellem Motiver og Handlinger, hvorfor Kant
i sin »Kritik der practischen Vernunft« (§ 6) stiller
sig følgende Opgave: »Forudsat at en Villie er fri
(sic), at finde den Lov, som ene egner sig til med
Nødvendighed (sic) at bestemme den.« Han undersøger
fremdeles omhyggelig, »hvorledes den ethiske
Lov kan blive Villiens Drivefjeder,« og finder, at
det sker gjennem Følelsen af Agtelse og Ærefrygt.
Vi have altsaa her kun en anden eller højere Art
af Nødvendighed, ikke Frihed i den Betydning, i
hvilken vi ovenfor have bekæmpet den. Ogsaa vi
have opfattet Friheden som et Ideal, der ikke er
umiddelbar givet. Forskjellen mellem vor Lære og
Kants bestaar deri, at han i den fornuftbestemte Frihed
ser Noget, der ligger ud over Erfaringens Verden,
medens den Frihed, vi tale om, netop staar
som Maalet for Menneskets Stræben og Udvikling i
Erfaringens Verden. Vi kunne her trøste os med,
at vi vise større Troskab mod Kants filosofiske
Princip end han selv gjør. Hans Storhed beror
paa, at han har paavist vor Erkjendelses Subjektivitet
og dens Begrænsning til Erfaringen. Alene
% --- p. 135 ---
\setcounter{footnote}{0}%
\opage{135}i den praktiske Fornuftlov mente han at have en
Undtagelse. Men med hvilken Ret kalder han da
denne Lov en »Kjendsgjerning«? En Kjendsgjerning
maa jo dog være tilgængelig for Erfaringen.
Naar han siger, at »Fornuften selv ikke er noget
Fænomen«, hvorledes lærer han den da at kjende,
da Erfaringen jo ifølge hans eget Princip er den
eneste Erkjendelseskilde? Og naar han siger, at
»Fornuften ikke er underkastet den sanselige Erfarings
Betingelser«, hvorledes kan Fornuften da blive
til og udvikle sig?\footnote[1]{Kant kommer her i Modsigelse med sig selv, idet han andre Steder (Kritik der pract. Vern. § 7 Anm. 2) siger, at den praktiske Fornuft er „en naturlig erhvervet Evne.“} Vi kjende ingen anden Fornuft
end den, der er Resultatet af Menneskets
aandelige Udvikling i Erfaringens Verden. Den er
altsaa selv et Fænomen og maa opfattes ifølge de
for alle Erfaringer gjældende Love. Grunden til,
at Kant hensatte den ethiske Fornuftlov og Friheden
i en anden Verden end Erfaringens, var deres
ideale Karakter i Forhold til de andre Motiver,
Evner og Kræfter i Mennesket. Det er en Rest
af Dogmatisme hos Kant, at han her anser det for
nødvendigt at tænke sig Idealet existerende. I
Virkeligheden vil det, at Mennesket idealt set er
frit, ikke sige Andet end, at det er hans Opgave og
Formaal at erhverve Friheden. Hvad kunde det
hjælpe, at Idealet existerede i en hinsidig Verden,
naar det ikke kunde tilstræbes og virkeliggjøres i
den Verden, vi leve i, Erfaringens Verden?

% --- p. 136 ---
\setcounter{footnote}{0}%
\opage{136}Dersom man vilde indvende herimod, at et saadant
Ideal umulig kan hævdes, eftersom Erfaringen
stadig godtgjør vor Mangel paa Frihed, da maa det
bemærkes, at ethvert som helst Ideal ifølge sit
Væsen peger ud over Virkeligheden. Uden denne
Modsætning vilde det aldrig komme til Handling.
Driften til at handle opstaar ved Forestillingen om
en Tilstand, som ikke er til Stede, men hvis Billede
Erindringen bevarer; naar Bevidsthedslivet har naaet
en vis Udvikling, træder i Stedet for den virkelige
Erindring Forestillingen om en fuldkomnere Tilstand,
som saa rigtig nok ofte forvexles med virkelig
Erindring; heraf Mythen om ideal Uskyldighed og
Lykke i et forsvundet Paradis. Enhver Drift viser
ud over det Givne, og Idealet er det fuldendte og
harmoniske Udtryk for, hvad der allerede delvis og
uklart dæmrer i den primitive Driftytring; derfor
har det sin Rod i Virkeligheden og fører dog ud
over den til enhver Tid givne Virkelighed. Kun
ved at bygge paa Determinismens Grund vil Ethiken
paa én Gang kunne hævde sit ideale Synspunkt og
bevare dets nøje Sammenhæng med det Virkelige.
Derfor have ogsaa store ethiske Tænkere søgt at
gjennemføre den deterministiske Tankegang. Særlig
kan her henvises til den engelske Skole (\emph{Hume},
\emph{Mill}, \emph{Spencer}), \emph{Herbart}, \emph{Schleiermacher},
\emph{Spinoza}, \emph{Stoikerne} og \emph{Kant}, naar vi tage ham
efter Aanden, ikke efter Bogstaven.

Man maa gaa tilbage til den humane Ethiks
Stifter, \emph{Sokrates}, og den fra ham udgaaende Bevægelse,
for at finde Forklaringen af de Vanskelig%
% --- p. 137 ---
\setcounter{footnote}{0}%
\opage{137}heder, hvori Spørgsmaalet om Villiens Frihed er
blevet indviklet. Sokrates førte den menneskelige
Handlen tilbage til Selverkjendelsen og hævdede
Viden som Dydens Væsen. Lidenskaben blev for
ham Uvidenhed eller urigtig Viden. Heraf udsprang
den Definition af Villien, som siden \emph{Platon}
og \emph{Aristoteles} har været den herskende, idet den
gjennem den kristelige Theologi ogsaa gjennemtrængte
den populære Opfattelse: Villien er den
af Fornuften ledede Drift. Hos Grækerne fremhævedes
vel stedse, ifølge deres hele naturalistiske
og æsthetiske Verdensopfattelse, den oprindelige
Begavelse, Temperamentet og Naturdriften som
væsenlig indgribende og betingende ved Menneskets
Handlen. Platon udleder al Slethed af sygelig Legemsbeskaffenhed
og slet Opdragelse og har et for
hans Tid overraskende skarpt Blik for Sindssygens
Væsen og Indflydelse paa Villien. Aristoteles fastholder
en lykkelig Naturbegavelse som nødvendig
Forudsætning for den ethiske Udvikling. Men den
kristelige Betragtning af Mennesket tog i sin Strenghed
ikke Hensyn til saadanne Betingelser. Den
betragtede Ansvarligheden som Menneskets inderste
Væsen. Selv de dunkle Drifter og Rørelser, der
endnu før Bevidsthedslivets Udvikling til Klarhed
og Styrke beherske Mennesket, betragtes som udsprungne
af en Villie, der kan drages til Ansvar.
»Alle Sindsbevægelser«, siger \emph{Augustinus}, »ere
ikke Andet end Villier. Thi hvad er Drift og Lyst
Andet end Villie, idet vi bifalde, hvad vi ville?
Og hvad er Frygt og Sorg Andet end Villie, idet
% --- p. 138 ---
\setcounter{footnote}{0}%
\opage{138}vi afsky, hvad vi ikke ville?«\footnote[1]{De civitate dei. XIV, 6.} »Ved det Sanselige,
Kjødelige,« siger S. \emph{Kierkegaard},\footnote[2]{Kjærlighedens Gjerninger (3die Udg.) I, 60.} »forstaar
Kristendommen det Selviske; der lader sig heller
ingen Strid tænke mellem Aand og Kjød, med
mindre der er en oprørsk Aand paa Kjødets Side,
med hvilken saa Aanden strider; der lader sig saaledes
ingen Strid tænke mellem Aand og en Sten
eller mellem Aand og et Træ. Altsaa det Selvkjærlige
er det Sanselige.« Der skjuler sig for
den konsekvente kristelige Betragtning en Villie i
enhver sjælelig Rørelse hos Mennesket, lige som
der bag enhver Naturbegivenhed ligger en guddommelig
Villiesakt. Her komme saaledes to absolute
Friheder til at staa over for hinanden, af
hvilke den ene er den almægtige, Alt skabende,
over for hvem den anden synes at maatte blive Intet.
Denne Vanskelighed, som Apostlen Paulus slaar
til Jorden ved at foreholde Mennesket hans absolute
Afmagt (hvorved han altsaa netop indrømmer
Vanskeligheden), søger Augustinus paa sin skarpsindige
Vis at løse ved at stille Villien udenfor
Naturen: »Gud er Ophav til al Natur og alle
Kræfter, men ikke til alle Villier. Det Onde har
ingen Natur; den onde Villie udspringer slet ikke
af nogen som helst Aarsag: det er jo først ved selve
Villien, at det Onde bliver til!« Naar Villien ikke
er Natur, ligger den ogsaa ud over al Legemlighed.
Sin Hovedindvending mod Platons Lære om Le%
% --- p. 139 ---
\setcounter{footnote}{0}%
\opage{139}gemsbeskaffenhedens Indflydelse paa Villien henter
Augustinus fra, at Syndens Ophav intet Legeme
har. Den sidste Kilde til det Onde er en naturløs
Villie, som vælger det Onde med fuld Viden og
Bevidsthed.\footnote[1]{De civ. dei. V, 9; XI, 9; XII, 6; XIV, 3.}

Til saadanne Konsekvenser føres man ved at
fastholde en Opfattelse af Villien, hvorved en ideal
Tilstand umiddelbar forudsættes hos Mennesket
som virkende Evne paa ethvert Trin af hans Udvikling
og i enhver Situation, i hvilken han befinder
sig. Vi ende her i det psykologisk Selvmodsigende:
thi den fulde Viden og Bevidsthed om
det Ondes Væsen, som er en Forudsætning for
det djævelske Standpunkt, vilde netop umuliggjøre
en Villen af det Onde. Det er umuligt, at Mennesket
i en og samme aandelige Tilstand lige saa godt
skulde kunne vælge den ene som den anden af to
modsatte Ting. Hovmodet, der skal være Syndens
Begyndelse, og som bestaar i, at Jeget vælger at
være sit eget Princip i Stedet for at have sit Princip
i en højere Magt, har lige saa vel som enhver
anden Karakteregenskab sin naturlige Oprindelse,
og de Handlinger, hvortil den fører, udspringe ifølge
de almindelige psykologiske Love. Den theologiske
Opfattelse bliver staaende ved det Faktum, at Jeget
vil gjøre sig selv til det Første og unddrage sig
Forholdet til den højere Magt. Den psykologiske
Opfattelse gaar længere tilbage og undersøger, hvorfor
og hvorledes Jeget kommer til at ville dette.
% --- p. 140 ---
\setcounter{footnote}{0}%
\opage{140}Værdien af en Karakters digteriske Fremstilling
beror paa, hvorvidt den giver os den fulde Forklaring
af alle Beslutninger og Handlinger. Selv en
Karakter som \emph{Shakspeares} Richard den Tredie, ved
hvilken Digteren maaske har overskredet den psykologiske
Sandhed og Rimeligheds Grænse, finder
sin forsonende Forklaring ved den allerførste Replik.
Foragt for Blødagtigheden og Usselheden omkring
ham i Forening med Følelsen af hans egen legemlige
og sjælelige Hæslighed, der gjør ham det umuligt
at deltage i den herskende letfærdige Livsglæde,
samtidig med at han i sig føler en Energi, der
kunde besejre alle Skranker, og hvorved han staar
højt over dem, der betragte ham som et Udskud,
--- alt dette avler den trodsige Beslutning at spille
Djævelens Rolle:

„og derfor, da jeg ikke kan som Bejler
behage denne tungerappe Tid,
saa har jeg sat mig for at spille Skurk
og hade disse Dages tomme Glæder.“
Digteren har tillige klart set, at man ved den psykologiske
Begrundelse ikke kan blive staaende ved
det enkelte Individ, men maa gaa ud over det til
hele Tidsalderen og de foregaaende Generationer.
Richard staar som Produkt af den røde og den
hvide Roses Kampe, af al den Grusomhed og Lumskhed,
som de avlede. Lidenskaberne have deres
Naturhistorie, og den ethiske Betragtning kan ikke
overspringe alle de Mellemled, hvorved Karakteren
og Handlingen først blive forstaalige, --- Mellemled,
som derfor egenlig ere Dele af selve Karak%
% --- p. 141 ---
\setcounter{footnote}{0}%
\opage{141}terens og Handlingens Væsen. En fransk Forfatter
har sagt: tout comprendre, ce serait tout pardonner,
og med Rette, for saa vidt som Handlingen ved at
ses i sin Nødvendighed ikke længere staar som det
rent Vilkaarlige. Men skjønt den ethiske Betragtning
ikke kan stride mod den psykologiske, saa
falder den dog ikke sammen med denne. I Psykologien
se vi tilbage, søge de forudgaaende Aarsager,
udlede de senere sjælelige Tilstande af de tidligere.
I Ethiken se vi fremad, maale Stemninger, Beslutninger og Handlinger efter et vist Forbillede og
spørge kun om deres harmoniske eller disharmoniske
Forhold til dette. Viser det sig, at de ikke stemme
med Normen, da bliver det Opgaven at finde Aarsagen
til denne Uoverensstemmelse, og naar den
er funden, at undersøge, hvor vidt det er muligt at
rydde denne Hindring bort.

\emph{Hobbes} har træffende bemærket, at hvis den
aristotelisk-skolastiske Definition af Villien som
fornuftig Drift (appetitus rationalis) var rigtig, saa
kunde der ikke gives nogen frivillig Handling, som
stred imod Fornuften. Hvad denne Definition
overser, er, at saa vel Driften som Fornuften maa
gjennemløbe en lang Udvikling, førend de kunne
forene sig og danne fornuftig Villie. Man kan derfor
ikke betragte denne som en Evne, der fra først
af tilhører Enhver, saa at det beror paa ham selv,
om han vil anvende den. I en vis Forstand handler
Mennesket aldrig mod Fornuften, nemlig mod
den Fornuft, han i Handlingens Øjeblik raader over;
derimod kan han handle mod den Fornuft, han
% --- p. 142 ---
\setcounter{footnote}{0}%
\opage{142}burde have, og som det er Maalet for hans Stræben
at tilegne sig. Deraf at man ikke kan, følger ikke,
at man ikke skal. Det Onde bestaar i, at Mennesket
\emph{endnu ikke} er godt, ikke i, at han er ophørt
med at være det. Den Grad af Karakter og Fornuftudvikling,
som Mennesket har opnaaet, kan
fordunkles og forringes, naar de rette Motiver og
Impulser ikke vedligeholdes, og naar Opmærksomheden
sløves. Men netop dette godtgjør jo, at han
endnu ikke har opnaaet en saadan Styrke, Sammenhæng
og Konsekvens i sin Karakter, at han kan
hævde det Erhvervede. Vi kjende, som allerede
bemærket, kun vor Karakter af Erfaringen, og den
er aldrig afsluttet. Heri ligger paa den anden Side
ogsaa, at vi aldrig med Vished kunne vide, om vi
have brugt alle vore Kræfter. Som vi i mange
Henseender ere tilbøjelige til at tiltro os selv for
meget, saaledes tro vi her maaske for ringe om,
hvad vi formaa. Vi maa her sige med Aladdin:
»Om jeg er Herre, gjælder et Forsøg!« Kun ad
experimental Vej kunne vi faa Sikkerhed her, som
paa saa mange andre Omraader. Hvis man vilde
lade sig afskrække fra saadanne Forsøg ved den
Tanke, at de ikke nytte, da den nødvendige Tingenes
Orden jo i hvert Tilfælde gaar sin Gang, saa
begaar man en Fejlslutning, som Logiken endog
har givet et særeget Navn (»den dovne Fornuft,«
ignava ratio), og som vilde svare til, at man af den
ydre Naturs Lovmæssighed udledte, at al vor Indgriben
og Indvirken paa Tingene var frugtesløs.
Netop fordi vor Villen og Handlen er Led i de
% --- p. 143 ---
\setcounter{footnote}{0}%
\opage{143}naturlige Begivenheders Række kunne vi modificere
dem ved at ændre de Betingelser, under hvilke
de indtræde.

Menneskenaturen er i sig selv hverken god eller
ond. Begreber som Godt og Ondt ere Forholdsbegreber,
som forudsætte Anvendelsen af et vist
Synspunkt, en vis Maalestok. Gaa vi ud fra, at
den menneskelige Natur formaar at hæve sig til en
høj Grad af Fornuft og Frihed, af Sandhed og Kjærlighed,
ville vi med Rette kalde den god. Tage vi
derimod Hensyn til den Modstand, der maa overvindes,
og den Kamp, der maa udholdes, for at
dette Maal kan naaes, kan man kalde den ond.
Den kristelige Lære om Arvesynden og den kantiske
Lære om det radikale Onde ere fra dette
Synspunkt fuldt berettigede; kun en glat, overfladisk
Optimisme kan tvivle derom Den Modstand, der % PRINTED AS IS: „derom Den“ (the point wanting; both copies)
ydes, saa vel som de Kræfter, der arbejde fremad,
udspringe af Menneskehedens Natur. J. G. \emph{Fichte}
søgte at vise, at det radikale Onde bestaar i Træghed,
Inerti. Lige som der i den ydre Natur hersker
en Inertiens Lov, idet enhver Ting bestaar i
sin forhaandenværende Tilstand, saa længe den ikke
ved Paavirkninger udefra føres ud over den, saaledes
vil det menneskelige Individ uden stærke Impulser
udefra ikke udvikle sig til Frihed og Sandhed.
Derfor er Livet i Samfundet, det inderlige Forhold
til Slægten, hele den historiske Side af Menneskelivet
saa afgjørende for det ethiske Liv. Vore tidligere
Undersøgelser have fra flere Sider lært os
Opdragelsens Nødvendighed at kjende. Uden Avtori%
% --- p. 144 ---
\setcounter{footnote}{0}%
\opage{144}teternes Indflydelse vilde Menneskeslægtens aandelige
Væxt ikke gaa for sig. Menneskehedens Historie
er væsenlig Avtoriteternes Historie. Under
denne store Opdragelsesproces gjør det Givne, det
hidtil Opnaaede, der er kommet i Ligevægt og
stivner i faste Former, Modstand mod de højere
Magter, der skulle arbejde sig frem, lige som ethvert
Legeme gjør Modstand mod et andet, der vil indtage
dets Plads. Paa det aandelige og særlig paa
det ethiske Omraade skal nu oven i Kjøbet den
højere Magt vinde Indgang gjennem Menneskets
Frihed, og vi have set, at denne bestaar i energisk
Opmærksomhed og Anspændelse; hertil maa Mennesket
altsaa først vækkes: de Kræfter, han skal
bruge, kan han ikke give sig selv, skjønt han skal
finde dem i sin egen Natur. Af Trægheden udvikler sig ikke blot Fejghed og Falskhed, som Fichte
har vist, men ogsaa Trods, det mest intensive Udtryk
for Inertien, der ikke lader sig føre ud over
sig selv for at deltage i en større Helheds Liv.
Trodsen forudsætter Inertiens mindre energiske
Former, og kun ved atter her at gjøre Enden til
Begyndelsen kan man betragte den som Syndens
Kilde. For at overvinde Modstanden maa man opsøge
de Aarsager, som oprindelig bragte Udviklingen
ind paa denne Vej; man maa her, som ved al Opdragelse,
finde Motiver og Midler, hvorved Jegets
Isolation kan ophæves, idet dets Interesse knyttes
til et væsenligt, almengyldigt Formaal.

Den humane Ethik, for hvem Guldalderen eller
Paradiset ikke er et tabt Gode, men et Maal, der,
% --- p. 145 ---
\setcounter{footnote}{0}%
\opage{145}saa vidt det er muligt og berettiget under menneskelige
Vilkaar, skal naaes ved Slægtens Arbejde, kan
ikke betragte Synden eller det Onde som noget
Absolut. Den kjender et »endnu ikke«, ej et »ikke
mere«. Den véd, at Maalet ifølge sin Natur maa
ligge ud over det Givne, og at der derfor altid maa
være en Afstand mellem Virkelighed og Ideal.
Denne Afstand tager den som given og forsøger
ikke at udfylde den ved Fantasiens Hjælp. Uden
en saadan Afstand var det jo heller ikke værd at
leve. (Af lignende Grund bad Baggesen om altid
at beholde syv uopfyldte Ønsker tilbage). For saa
vidt Afstanden skal udfyldes paa anden Maade end
ved stedse fornyet Stræben, er denne Udfyldning
Poesiens, ikke Tænkningens Sag. Men Afstanden
maa ikke forvexles med uforsonlig Modsætning;
tvertimod, selve Ufuldkommenhederne bane Vejen
for det Fuldkomnere. Og jo mere det bliver indlysende,
at Mennesket virkelig kan yde et Bidrag
til Dannelsen af sin egen og derved ogsaa af hele
Slægtens Karakter, desto mere vil ogsaa Interessen
koncentrere sig om dette Punkt, og desto flere
Omveje og Lidelser vil der spares. En fremragende
engelsk Sindssygelæge, \emph{Henry Maudsley}, har nylig
indskærpet Vigtigheden af Opmærksomhed og Interesse
for Karakterens Udvikling som Middel til at
bevare for Afsindighed. Han ser i Karakterens
Svækkelse, i den almindelige Tilbøjelighed til at
lade sig drive af Strømmen i Stedet for selvvirksom
at gribe ind for en stor Del Aarsagen til
Sindssygens Udbredelse. Men han fremhæver, at et
% --- p. 146 ---
\setcounter{footnote}{0}%
\opage{146}Resultat  her naaes langsomt  og maaske først bliver
synligt, naar Arbejdet fortsættes gjennem en Række
af Generationer\footnote[1]{Se det smukke Slutningskapitel af hans Bog „Le crime et la folie“. (Chap. IX: des moyens de se préserver de la folie).}. Afsindigheden staar i nøje Forbindelse
med de andre Onder, hvorunder Menneskeheden
lider; den er i sig selv kun en af de Former,
hvorunder Skrankerne for den højere Udvikling optræde.
Den rationelle Lære om Villiens Frihed,
som vi her have forsøgt at hævde, viser Muligheden
af at flytte disse Skranker, selv om det i det enkelte
Tilfælde kun var et Hanefjed. Vi maa i
Ethiken som paa andre Omraader ofte regne med
uendelig smaa Størrelser. Den indskærper tillige,
at det kun er muligt ved indtrængende psykologisk
Kjendskab til Karakteren og dens Betingelser,
lige som vi ogsaa kun betvinge den ydre Natur ved
Lydighed mod dens Love. Derimod vilde Antagelsen
af en absolut Frihed let føre til fantastiske
Forsøg paa at omskabe sig selv, til Fortvivlelse
over at man ikke kan det, eller til Tvivl om al
lovmæssig Sammenhæng paa det aandelige Omraade
i det Hele.

\begin{center}\rule{6em}{0.4pt}\end{center}% an ornamental rule (a small lozenge on a hairline) ends the book, p. 146

\end{document}
